\documentclass[11pt,a4paper,oneside]{report}
\usepackage[T1]{fontenc}
\usepackage[utf8]{inputenc}
\usepackage[italian]{babel}
\usepackage{lmodern}
\usepackage{amsmath,amssymb,amsthm,mathtools}
\usepackage[a4paper,margin=27mm,headheight=14pt]{geometry}
\usepackage{microtype}
\usepackage{enumitem}
\usepackage{booktabs}
\usepackage{needspace}
\usepackage{tikz}
\usetikzlibrary{arrows.meta}
\usepackage{fancyhdr}
\usepackage{hyperref}
\usepackage{bookmark}

\definecolor{coloreschema}{RGB}{30,77,112}
\hypersetup{
  unicode=true,
  hidelinks,
  pdftitle={Appunti di analisi 1},
  pdfauthor={Jude De Sa'Pinto},
  pdfsubject={Analisi 1 -- Prof. Federico Cacciafesta},
  bookmarksnumbered=true
}
\setcounter{tocdepth}{2}
\setcounter{secnumdepth}{2}
\setlist[enumerate]{itemsep=.35em,topsep=.4em}
\setlist[itemize]{itemsep=.25em,topsep=.4em}
\raggedbottom
\emergencystretch=1.5em
% Conserva titoli, provenienza e prime righe sulla stessa pagina.
\AddToHook{cmd/section/before}{\Needspace{9\baselineskip}}
\AddToHook{cmd/subsection/before}{\Needspace{7\baselineskip}}

\pagestyle{fancy}
\fancyhf{}
\fancyhead[L]{\small\nouppercase{\leftmark}}
\fancyfoot[C]{\thepage}
\renewcommand{\headrulewidth}{.3pt}
\renewcommand{\chaptermark}[1]{\markboth{#1}{}}

\theoremstyle{plain}
\newtheorem{teorema}{Teorema}[chapter]
\newtheorem{assioma}[teorema]{Assioma}
\newtheorem{proposizione}[teorema]{Proposizione}
\theoremstyle{definition}
\newtheorem{definizione}[teorema]{Definizione}
\newtheorem*{esempio}{Esempio}
\theoremstyle{remark}
\newtheorem*{osservazione}{Osservazione}

% Evita di lasciare l'inizio di un enunciato isolato a fondo pagina.
\AddToHook{env/proposizione/before}{\Needspace{6\baselineskip}}
\AddToHook{env/teorema/before}{\Needspace{6\baselineskip}}
\AddToHook{env/assioma/before}{\Needspace{6\baselineskip}}

\newcommand{\N}{\mathbb{N}}
\newcommand{\Z}{\mathbb{Z}}
\newcommand{\Q}{\mathbb{Q}}
\newcommand{\R}{\mathbb{R}}
\newcommand{\fonti}[1]{\par\nobreak\noindent{\small\itshape #1.}\par\medskip}
\newenvironment{schema}
  {\par\medskip\noindent\begin{minipage}{\linewidth}\centering\small}
  {\end{minipage}\par\medskip}
\tikzset{
  asse/.style={-{Stealth[length=2mm]},draw=black!55,line width=.5pt},
  intervallo/.style={draw=coloreschema,line width=1.5pt},
  estremo/.style={draw=coloreschema,fill=white,line width=.9pt}
}


\title{Appunti di analisi 1\\[.8em]\large Prof. Federico Cacciafesta}
\author{Jude De Sa'Pinto}
\date{}

\begin{document}
\pagenumbering{roman}
\hypersetup{pageanchor=false}
\maketitle
\hypersetup{pageanchor=true}
\tableofcontents
\clearpage

\chapter*{Convenzioni e notazioni}
\addcontentsline{toc}{chapter}{Convenzioni e notazioni}
Scriviamo $\N=\{0,1,2,\ldots\}$ e $\N^*=\{1,2,3,\ldots\}$; con $\Z$, $\Q$ e $\R$ indichiamo rispettivamente gli interi, i razionali e i reali. Il simbolo $\subseteq$ ammette l'uguaglianza; per l'inclusione stretta usiamo $\subsetneq$. Il complemento $A^c$ è sempre relativo all'insieme ambiente specificato; nel capitolo sulla topologia l'ambiente è $\R$.

Usiamo $(a,b)$ per l'intervallo aperto, $[a,b]$ per quello chiuso e le notazioni miste $[a,b)$ e $(a,b]$ quando è incluso un solo estremo. La notazione alternativa $]a,b[$ ha lo stesso significato di $(a,b)$.

\clearpage
\pagenumbering{arabic}
% !TeX root = ../main.tex
% Fonte letta integralmente: Lezione 1.pdf, pp. 2--6.
% Le spiegazioni, la tabella di verità e gli esempi sviluppano i passaggi
% sintetici della fonte senza attribuire al docente formulazioni aggiunte.
\chapter{Logica e metodi di dimostrazione}
\label{cap:logica}
\fonti{Lezione 1}

\section{Proposizioni e connettivi logici}
% Fonte: Lezione 1, pp.~2--3

Una \emph{proposizione} è un'affermazione alla quale si può attribuire un
valore di verità: vero oppure falso. Per esempio, «$2+3=5$» è una
proposizione vera, mentre «$2+3=6$» è una proposizione falsa. Una frase
contenente una variabile libera, come «$x>0$», richiede invece di sapere
quale valore si assegna a $x$ prima di poter stabilire se sia vera o falsa.

Per costruire proposizioni complesse a partire da proposizioni $P$ e $Q$
si usano i \emph{connettivi logici}:
\begin{description}[style=nextline]
  \item[Negazione: $\neg P$.]
    È vera esattamente quando $P$ è falsa; si legge «non $P$».
  \item[Congiunzione: $P\land Q$.]
    È vera quando sono vere sia $P$ sia $Q$; si legge «$P$ e $Q$».
  \item[Disgiunzione: $P\lor Q$.]
    È vera quando è vera almeno una delle due proposizioni. Si tratta
    della disgiunzione inclusiva: è vera anche se sono vere entrambe.
  \item[Implicazione: $P\Rightarrow Q$.]
    È falsa soltanto quando $P$ è vera e $Q$ è falsa. $P$ è l'ipotesi
    e $Q$ la tesi; si legge «se $P$, allora $Q$».
  \item[Doppia implicazione: $P\Leftrightarrow Q$.]
    È vera quando $P$ e $Q$ hanno lo stesso valore di verità: entrambe
    vere oppure entrambe false. Si legge «$P$ se e solo se $Q$».
\end{description}

La tabella seguente riassume tutti i casi; $V$ significa vero e $F$ falso.
\begin{center}
\renewcommand{\arraystretch}{1.2}
\begin{tabular}{ccccccc}
\toprule
$P$ & $Q$ & $\neg P$ & $P\land Q$ & $P\lor Q$ & $P\Rightarrow Q$ & $P\Leftrightarrow Q$\\
\midrule
$V$ & $V$ & $F$ & $V$ & $V$ & $V$ & $V$\\
$V$ & $F$ & $F$ & $F$ & $V$ & $F$ & $F$\\
$F$ & $V$ & $V$ & $F$ & $V$ & $V$ & $V$\\
$F$ & $F$ & $V$ & $F$ & $F$ & $V$ & $V$\\
\bottomrule
\end{tabular}
\end{center}

\section{Implicazioni e metodi di dimostrazione}
% Fonte: Lezione 1, p.~3

L'implicazione $P\Rightarrow Q$ esprime che la verità di $P$ basta a
garantire la verità di $Q$. Per questo $P$ è detta \emph{condizione
sufficiente} per $Q$ e $Q$ \emph{condizione necessaria} per $P$.
L'implicazione non afferma che $P$ sia vera: esclude soltanto il caso
in cui l'ipotesi sia vera e la tesi falsa. In particolare,
\[
  P\Rightarrow Q\quad\Longleftrightarrow\quad \neg P\lor Q.
\]
Per dimostrare una doppia implicazione si devono verificare entrambi
i versi, $P\Rightarrow Q$ e $Q\Rightarrow P$.

\begin{proposizione}[Contrapposizione]
Un'implicazione è equivalente alla sua contronominale:
\[
  (P\Rightarrow Q)\quad\Longleftrightarrow\quad
  (\neg Q\Rightarrow\neg P).
\]
\end{proposizione}
\begin{proof}
La prima implicazione è falsa esattamente nel caso $P$ vera e $Q$ falsa.
La seconda è falsa esattamente quando $\neg Q$ è vera e $\neg P$ è falsa,
cioè nello stesso caso. Le due implicazioni hanno quindi sempre lo
stesso valore di verità.
\end{proof}

Questa equivalenza fornisce un metodo di dimostrazione: se è più semplice
partire dalla negazione della tesi, si dimostra $\neg Q\Rightarrow\neg P$
e si conclude $P\Rightarrow Q$.

Il \emph{ragionamento per assurdo} consiste nell'assumere la negazione di ciò che si
vuole dimostrare e nel ricavarne una contraddizione. Per provare
$P\Rightarrow Q$ si può dunque assumere che $P$ sia vera e $Q$ falsa e
mostrare che queste due assunzioni non possono valere contemporaneamente.
I due metodi sono strettamente collegati, ma conviene distinguerne
il modo di procedere.

\section{Predicati e quantificatori}
% Fonte: Lezione 1, pp.~4--5

Un \emph{predicato} $P(x)$ è un'affermazione che dipende da una variabile
$x$. Il suo significato va accompagnato dall'insieme nel quale $x$ può
variare: per esempio, $P(x)$ può essere la proprietà «$x^2=1$» con
$x\in\R$. I quantificatori permettono di trasformare un predicato
in una proposizione.

\begin{definizione}[Quantificatore universale]
La proposizione
\[
  \forall x\in U,\quad P(x)
\]
è vera se $P(a)$ è vera per ogni valore $a$ appartenente all'insieme
ambiente $U$. Il simbolo $\forall$ si legge «per ogni» o «qualunque sia».
\end{definizione}

\begin{definizione}[Quantificatore esistenziale]
La proposizione
\[
  \exists x\in U:\ P(x)
\]
è vera se esiste almeno un elemento $a\in U$ per il quale $P(a)$ è vera.
Il simbolo $\exists$ si legge «esiste». Non si esclude che gli elementi
con questa proprietà siano più di uno.
\end{definizione}

\begin{definizione}[Esistenza e unicità]
La proposizione
\[
  \exists!\,x\in U:\ P(x)
\]
è vera quando esiste esattamente un elemento di $U$ che soddisfa $P$.
Per dimostrarla occorre provare sia l'esistenza sia l'unicità: dopo
aver trovato un elemento $a$ con la proprietà richiesta, si mostra che
ogni altro elemento che la soddisfa coincide con $a$.
\end{definizione}

Per esempio, $\exists x\in\R:x^2=1$ è vera, ma
$\exists!\,x\in\R:x^2=1$ è falsa, perché sia $1$ sia $-1$ risolvono
l'equazione. Se imponiamo anche $x>0$, la soluzione diventa unica.

\section{Negare una proposizione con quantificatori}
% Fonte: Lezione 1, p.~6

Negare un'affermazione universale significa trovare almeno un caso
in cui la proprietà non vale; negare un'affermazione esistenziale
significa affermare che la proprietà non vale per nessun elemento.
Formalmente,
\begin{align*}
  \neg\bigl(\forall x\in U,\ P(x)\bigr)
    &\quad\Longleftrightarrow\quad \exists x\in U:\ \neg P(x),\\
  \neg\bigl(\exists x\in U:\ P(x)\bigr)
    &\quad\Longleftrightarrow\quad \forall x\in U,\ \neg P(x).
\end{align*}
Quando la negazione attraversa un quantificatore, dunque, $\forall$
si trasforma in $\exists$ e $\exists$ si trasforma in $\forall$;
anche la proprietà all'interno deve essere negata.

Per esempio, la negazione di «tutti gli elementi di $U$ sono positivi»
è «esiste un elemento di $U$ che è minore o uguale a zero».
Non è necessario che siano tutti negativi: basta un controesempio.
Se sono presenti più quantificatori, si applica la regola a ciascuno,
mantenendo il loro ordine. Per esempio,
\[
  \neg\bigl(\forall x\in\R\ \exists y\in\R:\ y>x\bigr)
  \quad\Longleftrightarrow\quad
  \exists x\in\R\ \forall y\in\R:\ y\le x.
\]

% !TeX root = ../main.tex
% Fonte letta integralmente: Lezione 1.pdf, pp. 7--13.
% Gli esempi finiti e il diagramma di Venn sono aggiunte esplicative.
\chapter{Insiemi e relazioni}
\label{cap:insiemi}
\fonti{Lezione 1}

\section{Insiemi, elementi e insieme ambiente}
% Fonte: Lezione 1, pp.~7--8

Il concetto di \emph{insieme} viene assunto come primitivo: intuitivamente indica una raccolta di oggetti, detti \emph{elementi} o \emph{punti}. Un insieme $A$ è assegnato quando, per ogni oggetto considerato, la frase «$x$ appartiene ad $A$» ha un valore di verità. Scriviamo
\[
  x\in A\quad\text{se $x$ appartiene ad $A$},\qquad
  x\notin A\quad\text{se non vi appartiene}.
\]
\textbf{Uguaglianza e rappresentazione.}
\begin{itemize}
  \item Due insiemi sono uguali se e solo se hanno gli stessi elementi.
  \item L'ordine dell'elenco e le eventuali ripetizioni non modificano l'insieme.
\end{itemize}

\textbf{Descrizione per elenco o per proprietà.} Si può assegnare un insieme elencandone gli elementi oppure indicando una proprietà che li caratterizza. Per esempio,
\[
  A=\{1,3,5,7,9\}
   =\{x\in\N:0\le x\le10\ \land\ x\text{ è dispari}\}.
\]
In generale, una descrizione per proprietà ha la forma $A=\{x\in U:P(x)\}$: $U$ è l'\emph{insieme ambiente} e si selezionano gli elementi di $U$ che soddisfano $P$.

\textbf{Insieme vuoto.} Si assume l'esistenza di un insieme privo di elementi, detto \emph{insieme vuoto} e indicato con $\varnothing$. Esso è diverso dal singoletto $\{\varnothing\}$: quest'ultimo contiene un elemento, che è a sua volta l'insieme vuoto.

\section{Sottoinsiemi e insieme delle parti}
% Fonte: Lezione 1, p.~9

\begin{definizione}[Inclusione]
Un insieme $B$ è un \emph{sottoinsieme} di $A$ se ogni elemento di $B$ è anche un elemento di $A$. Si scrive $B\subseteq A$, ossia
\[
  B\subseteq A\quad\Longleftrightarrow\quad
  \forall x,\ (x\in B\Rightarrow x\in A).
\]
\end{definizione}

\begin{itemize}
  \item \textbf{Sottoinsiemi impropri.} $\varnothing$ e $A$ sono sempre sottoinsiemi di $A$. L'inclusione $\varnothing\subseteq A$ vale perché il vuoto non ha elementi che possano violarla.
  \item \textbf{Inclusione stretta.} Se $B\subseteq A$ e $B\ne A$, scriviamo $B\subsetneq A$.
\end{itemize}

\textbf{Insieme delle parti.} L'\emph{insieme delle parti} di $A$, indicato con $\mathcal P(A)$, è l'insieme di tutti i suoi sottoinsiemi:
\[
  \mathcal P(A)=\{B:B\subseteq A\}.
\]
\begin{esempio}
Se $A=\{1,2\}$, allora $\mathcal P(A)=\{\varnothing,\{1\},\{2\},\{1,2\}\}$. È utile distinguere l'appartenenza dall'inclusione: $1\in A$, mentre $\{1\}\subseteq A$ e $\{1\}\in\mathcal P(A)$.
\end{esempio}

\section{Operazioni tra insiemi}
% Fonte: Lezione 1, pp.~10--11

Fissiamo due insiemi $A$ e $B$ contenuti nello stesso ambiente $U$. Le operazioni fondamentali traducono le corrispondenti operazioni logiche.

\begin{definizione}[Unione, intersezione e differenza]
Si definiscono
\begin{align*}
  A\cup B&=\{x\in U:x\in A\lor x\in B\},\\
  A\cap B&=\{x\in U:x\in A\land x\in B\},\\
  A\setminus B&=\{x\in U:x\in A\land x\notin B\}.
\end{align*}
L'unione corrisponde all'appartenenza ad almeno uno dei due insiemi; l'intersezione all'appartenenza a entrambi; la differenza all'appartenenza ad $A$ ma non a $B$.
\end{definizione}

\begin{schema}
\begin{tikzpicture}[x=1cm,y=1cm]
  \begin{scope}
    \fill[coloreschema!20] (-.45,0) circle[radius=.85];
    \fill[coloreschema!20] (.45,0) circle[radius=.85];
    \draw (-.45,0) circle[radius=.85];
    \draw (.45,0) circle[radius=.85];
    \node at (-.8,.35) {$A$}; \node at (.8,.35) {$B$};
    \node at (0,-1.2) {$A\cup B$};
  \end{scope}
  \begin{scope}[xshift=4cm]
    \begin{scope}
      \clip (-.45,0) circle[radius=.85];
      \fill[coloreschema!20] (.45,0) circle[radius=.85];
    \end{scope}
    \draw (-.45,0) circle[radius=.85];
    \draw (.45,0) circle[radius=.85];
    \node at (-.8,.35) {$A$}; \node at (.8,.35) {$B$};
    \node at (0,-1.2) {$A\cap B$};
  \end{scope}
  \begin{scope}[xshift=8cm]
    \fill[coloreschema!20] (-.45,0) circle[radius=.85];
    \fill[white] (.45,0) circle[radius=.85];
    \draw (-.45,0) circle[radius=.85];
    \draw (.45,0) circle[radius=.85];
    \node at (-.8,.35) {$A$}; \node at (.8,.35) {$B$};
    \node at (0,-1.2) {$A\setminus B$};
  \end{scope}
\end{tikzpicture}
\par\smallskip La regione colorata rappresenta l'operazione indicata sotto ciascun disegno.
\end{schema}

\begin{esempio}
Con $A=\{1,2,3\}$ e $B=\{3,4\}$ si ha $A\cup B=\{1,2,3,4\}$, $A\cap B=\{3\}$ e $A\setminus B=\{1,2\}$. Scambiando i due insiemi si ottiene invece $B\setminus A=\{4\}$: la differenza dipende dall'ordine.
\end{esempio}

\textbf{Complementare e insieme ambiente.} Se $B\subseteq A$, l'insieme $A\setminus B$ è detto \emph{complementare di $B$ in $A$}. Quando l'ambiente è $U$, si scrive semplicemente $B^c=U\setminus B$. L'ambiente è essenziale: il complemento dello stesso insieme cambia se cambia $U$.

\section{Le leggi di De Morgan}
% Fonte: Lezione 1, p.~11

\begin{proposizione}[Leggi di De Morgan]
Per $A,B\subseteq U$, con i complementi calcolati nello stesso ambiente,
\[
  (A\cup B)^c=A^c\cap B^c,\qquad
  (A\cap B)^c=A^c\cup B^c.
\]
\end{proposizione}
% Evita che la sola frase conclusiva della seconda identità passi alla pagina seguente.
\enlargethispage{\baselineskip}
\begin{proof}
\textbf{Prima identità.} Per ogni $x\in U$,
\[
  \begin{aligned}
    x\in(A\cup B)^c
    &\iff x\notin A\cup B\\
    &\iff x\notin A\ \land\ x\notin B\\
    &\iff x\in A^c\cap B^c.
  \end{aligned}
\]
Le due condizioni di appartenenza sono equivalenti per ogni elemento dell'ambiente, quindi i due insiemi sono uguali.

\textbf{Seconda identità.} Analogamente, per ogni $x\in U$,
\[
  \begin{aligned}
    x\in(A\cap B)^c
    &\iff x\notin A\cap B\\
    &\iff x\notin A\ \lor\ x\notin B\\
    &\iff x\in A^c\cup B^c.
  \end{aligned}
\]
Si ottiene anche la seconda uguaglianza.
\end{proof}

Le formule esprimono a livello di insiemi le equivalenze logiche $\neg(P\lor Q)\Leftrightarrow(\neg P\land\neg Q)$ e $\neg(P\land Q)\Leftrightarrow(\neg P\lor\neg Q)$. Saranno utili per passare dalle proprietà degli aperti a quelle dei chiusi.

\section{Prodotto cartesiano e relazioni}
% Fonte: Lezione 1, p.~12

\begin{definizione}[Prodotto cartesiano]
Il prodotto cartesiano di $A$ e $B$ è l'insieme delle coppie ordinate con primo elemento in $A$ e secondo elemento in $B$:
\[
  A\times B=\{(a,b):a\in A,\ b\in B\}.
\]
\end{definizione}

\textbf{Coppie ordinate.} Una coppia è \emph{ordinata}: $(a,b)=(a',b')$ se e solo se $a=a'$ e $b=b'$. In particolare, scambiare le componenti modifica in generale la coppia. Per esempio,
\[
  \{1,2\}\times\{a,b\}
   =\{(1,a),(1,b),(2,a),(2,b)\}.
\]

\begin{definizione}[Relazione]
Dati due insiemi non vuoti $A$ e $B$, una relazione tra $A$ e $B$ è un sottoinsieme $R\subseteq A\times B$. Si dice che $a\in A$ e $b\in B$ sono in relazione, e si scrive $a\,R\,b$, quando $(a,b)\in R$. Se $A=B$, si parla di una relazione \emph{su} $A$.
\end{definizione}

\begin{esempio}
Specificare una relazione significa scegliere quali coppie considerare. Su $A=\{1,2,3\}$ la relazione «essere minore o uguale» è il sottoinsieme
\[
  R=\{(1,1),(1,2),(1,3),(2,2),(2,3),(3,3)\}\subseteq A\times A.
\]
\end{esempio}

\section{Relazioni d'ordine e ordine totale}
% Fonte: Lezione 1, p.~13

\begin{definizione}[Relazione d'ordine]
Una relazione $R$ su un insieme non vuoto $A$ si dice \emph{relazione d'ordine} se soddisfa le seguenti proprietà:
\begin{enumerate}[label=(\alph*)]
  \item \textbf{Riflessività:} per ogni $a\in A$, $a\,R\,a$.
  \item \textbf{Antisimmetria:} per ogni $a,b\in A$, $(a\,R\,b\land b\,R\,a)\Rightarrow a=b$.
  \item \textbf{Transitività:} per ogni $a,b,c\in A$, $(a\,R\,b\land b\,R\,c)\Rightarrow a\,R\,c$.
\end{enumerate}
\end{definizione}

\textbf{Antisimmetria e simmetria.} L'antisimmetria non richiede di vietare ogni relazione in entrambi i versi: richiede che ciò avvenga soltanto per elementi uguali. È una condizione diversa dalla simmetria, che richiederebbe invece la validità della relazione in entrambi i versi per ogni coppia in relazione.

\begin{definizione}[Ordine totale]
Un ordine si dice \emph{totale} se ogni coppia di elementi è confrontabile:
\[
  \forall a,b\in A,\qquad a\,R\,b\ \lor\ b\,R\,a.
\]
\end{definizione}

\begin{esempio}
\leavevmode
\begin{itemize}
  \item L'ordine usuale $\le$ su $\N$, $\Z$, $\Q$ o $\R$ è totale.
  \item L'inclusione è un ordine su $\mathcal P(A)$, ma in generale non è totale. Per $A=\{1,2\}$, i sottoinsiemi $\{1\}$ e $\{2\}$ non sono confrontabili: nessuno contiene l'altro. Riflessività, antisimmetria e transitività restano valide.
\end{itemize}
\end{esempio}

% !TeX root = ../main.tex
% Fonti lette sulle immagini: Lezione 2 AL.pdf, pp. 2--11;
% Lezione 3.pdf, pp. 2--4. Nessuna pagina di contenuto di questi intervalli omessa.
% La prova del binomio di Newton e le interpretazioni combinatorie sviluppano
% gli enunciati sintetici della fonte. I registri sono in lavoro/SESSIONI.md.
\chapter{Numeri naturali, induzione e combinatoria}
\label{cap:naturali-combinatoria}
\fonti{Lezioni 2 e 3}

\section{I naturali e il principio del minimo intero}
% Fonte: Lezione 2, p.~2

Gli insiemi numerici che incontreremo si dispongono nella catena
\[
  \N\subseteq\Z\subseteq\Q\subseteq\R\subseteq\mathbb C,
\]
dove $\mathbb C$ indica l'insieme dei numeri complessi, richiamato qui
soltanto per collocare i diversi sistemi numerici. Partiamo dai naturali,
che comprendono lo zero:
\[
  \N=\{0,1,2,3,\ldots\}.
\]
L'ordine usuale $\le$ rende $\N$ totalmente ordinato: dati due naturali,
è sempre possibile confrontarli. Per lavorare sui naturali si assume
inoltre la seguente proprietà, più forte del solo ordine totale.

\begin{assioma}[Principio del minimo intero]
Ogni sottoinsieme non vuoto di $\N$ ammette un minimo. In formule,
\[
  S\subseteq\N,\quad S\ne\varnothing
  \quad\Longrightarrow\quad
  \exists m\in S:\ \forall n\in S,\ m\le n.
\]
\end{assioma}

Il minimo deve appartenere a $S$ ed essere minore o uguale a tutti gli
altri suoi elementi. Il principio permette di parlare del \emph{primo}
naturale che soddisfa una certa proprietà; in una dimostrazione per assurdo
permetterà, per esempio, di scegliere il primo eventuale controesempio.

\section{Il principio di induzione}
% Fonte: Lezione 2, pp.~3 e~8

Il principio di induzione formalizza un'idea semplice: se una proprietà
vale all'inizio e passa da ogni naturale al suo successivo, allora vale
per tutti i naturali a partire da quello iniziale. Possiamo esprimerlo
in due forme equivalenti.

\begin{proposizione}[Induzione: forma insiemistica]
Sia $S\subseteq\N$. Se $0\in S$ e, per ogni $n\in S$, anche
$n+1\in S$, allora $S=\N$.
\end{proposizione}

\begin{proposizione}[Induzione: forma con un predicato]
Sia $n_0\in\N$ e sia $P(n)$ un predicato definito per i naturali
$n\ge n_0$. Supponiamo che valgano:
\begin{enumerate}
  \item la \emph{base induttiva}: $P(n_0)$ è vera;
  \item il \emph{passo induttivo}: per ogni $n\ge n_0$,
    $P(n)\Rightarrow P(n+1)$.
\end{enumerate}
Allora $P(n)$ è vera per ogni $n\ge n_0$.
\end{proposizione}

\begin{proof}
Ragioniamo per assurdo.
Supponiamo che esista almeno un naturale $n\ge n_0$ per cui $P(n)$ è falsa.
L'insieme dei controesempi
\[
  F=\{n\in\N:n\ge n_0\text{ e }P(n)\text{ è falsa}\}
\]
è allora non vuoto. Per il principio del minimo intero esiste
$\bar n=\min F$. Non può essere $\bar n=n_0$, perché $P(n_0)$ è vera
per la base induttiva. Dunque $\bar n>n_0$, e il naturale precedente
$\bar n-1$ appartiene ancora all'intervallo di indici considerato.

Poiché $\bar n$ è il primo controesempio, $\bar n-1\notin F$:
$P(\bar n-1)$ è vera. Applicando il passo induttivo con $n=\bar n-1$
otteniamo che anche $P(\bar n)$ è vera, in contraddizione con
$\bar n\in F$. Pertanto $F$ è vuoto e la proprietà vale per tutti
i naturali $n\ge n_0$.
\end{proof}

La forma insiemistica si ricava scegliendo $n_0=0$ e ponendo
$P(n)$ uguale alla proprietà «$n\in S$». Viceversa, l'insieme degli indici
per cui un predicato è vero rende esplicito il collegamento fra le due
formulazioni. È essenziale verificare entrambi i passaggi: il solo passo
induttivo non garantisce che esista un punto da cui iniziare.

\section{La somma dei primi naturali}
% Fonte: Lezione 2, pp.~4--6

Consideriamo la somma
\[
  1+2+\cdots+99+100.
\]
Accoppiando il primo termine con l'ultimo, il secondo con il penultimo
e così via, si ottengono $50$ coppie di somma $101$:
\[
  (1+100)+(2+99)+\cdots+(50+51)
  =50\cdot101=\frac{100\cdot101}{2}=5050.
\]
Questo suggerisce una formula valida per un numero qualsiasi di termini.
La notazione $\sum_{k=1}^{n}k$ abbrevia la somma $1+2+\cdots+n$:
$k$ è l'indice che percorre i termini, mentre $n$ indica dove la somma
si arresta.

\begin{proposizione}
Per ogni $n\in\N^*$ vale
\[
  \sum_{k=1}^{n}k=\frac{n(n+1)}2.
\]
\end{proposizione}
\begin{proof}
Procediamo per induzione a partire da $n_0=1$. La base è vera, perché
$\sum_{k=1}^{1}k=1=1(1+1)/2$.

Per il passo induttivo fissiamo $n\ge1$ e assumiamo l'\emph{ipotesi
induttiva} $\sum_{k=1}^{n}k=n(n+1)/2$. La tesi al passo successivo è
\[
  \sum_{k=1}^{n+1}k=\frac{(n+1)(n+2)}2.
\]
Separiamo l'ultimo addendo e usiamo l'ipotesi soltanto per la somma
dei primi $n$ termini:
\begin{align*}
  \sum_{k=1}^{n+1}k
  &=\left(\sum_{k=1}^{n}k\right)+(n+1)\\
  &=\frac{n(n+1)}2+(n+1)
    =\frac{n^2+3n+2}{2}
    =\frac{(n+1)(n+2)}2.
\end{align*}
Abbiamo dimostrato $P(n)\Rightarrow P(n+1)$; per il principio di
induzione la formula vale per ogni $n\ge1$.
\end{proof}

\section{La disuguaglianza di Bernoulli}
% Fonte: Lezione 2, pp.~6--7

\begin{proposizione}[Disuguaglianza di Bernoulli]
Per ogni $n\in\N$ e ogni reale $x>-1$ si ha
\[
  (1+x)^n\ge1+nx.
\]
\end{proposizione}
\begin{proof}
Fissiamo un qualunque $x>-1$ e procediamo per induzione su $n$.
Per $n=0$, entrambi i membri valgono $1$.

Supponiamo ora che $(1+x)^n\ge1+nx$. Poiché $1+x>0$, possiamo
moltiplicare questa disuguaglianza per $1+x$ senza cambiarne il verso:
\begin{align*}
  (1+x)^{n+1}
  &=(1+x)^n(1+x)\\
  &\ge(1+nx)(1+x)\\
  &=1+(n+1)x+nx^2\\
  &\ge1+(n+1)x.
\end{align*}
L'ultimo passaggio usa $nx^2\ge0$. La disuguaglianza è così dimostrata
per $n+1$, e quindi per ogni $n\in\N$. Poiché $x>-1$ era arbitrario,
l'enunciato vale per tutti i valori di $x$ richiesti.
\end{proof}

La condizione su $x$ entra precisamente nel passaggio in cui si moltiplica
l'ipotesi induttiva per $1+x$: anche quando $x$ è negativo, il fattore
$1+x$ resta positivo.

\section{Permutazioni, disposizioni e combinazioni}
% Fonte: Lezione 2, pp.~9--10

Nel calcolo combinatorio si contano modi diversi di ordinare o scegliere
elementi di un insieme finito. Le formule seguenti riguardano
oggetti \emph{distinti} e scelte \emph{senza ripetizione}: un oggetto
già scelto non viene scelto di nuovo.

\subsection{Permutazioni e fattoriale}

Una \emph{permutazione} di $n$ oggetti è un modo di disporli tutti in fila.
Per il primo posto ci sono $n$ possibilità, per il secondo $n-1$ e così
via fino all'ultimo. Il numero di permutazioni è dunque
\[
  P_n=n!=n(n-1)(n-2)\cdots2\cdot1.
\]
Il simbolo $n!$ si legge \emph{fattoriale di $n$}. Si pone $0!=1$:
anche l'insieme vuoto ha un solo ordinamento, quello senza elementi.
Per esempio, tre oggetti distinti hanno $3!=6$ possibili ordinamenti.

\subsection{Disposizioni: scegliere e ordinare}

Una \emph{disposizione} di $k$ oggetti scelti fra $n$ consiste nel
selezionare $k$ elementi distinti e disporli in fila, tenendo conto
dell'ordine. Per $0\le k\le n$ il numero di disposizioni è
\[
  D_{n,k}=n(n-1)\cdots(n-k+1)=\frac{n!}{(n-k)!}.
\]
Per $k=0$ il prodotto è vuoto e vale $1$; per $k=n$ ritroviamo le
permutazioni. Per esempio, per assegnare un primo e un secondo posto
fra cinque persone ci sono $D_{5,2}=5\cdot4=20$ possibilità.

\subsection{Combinazioni: scegliere senza ordinare}

Una \emph{combinazione} di $k$ oggetti fra $n$ è una scelta nella quale
l'ordine non conta: equivale a un sottoinsieme di $k$ elementi.
Indichiamo il numero di queste scelte con $C_{n,k}$. Ogni scelta di
$k$ elementi produce $k!$ disposizioni, perché i suoi elementi possono
essere ordinati in $k!$ modi. Quindi
\[
  D_{n,k}=C_{n,k}\,k!,\qquad
  C_{n,k}=\frac{n!}{k!(n-k)!}.
\]
Per esempio, scegliendo due persone fra cinque senza assegnare ruoli
si ottiene $C_{5,2}=20/2=10$.

\section{Coefficienti binomiali e loro proprietà}
% Fonte: Lezione 2, pp.~10--11; lezione 3, p.~2

Per $n\in\N$ e $k\in\Z$ definiamo il \emph{coefficiente binomiale}
\[
  \binom nk=
  \begin{cases}
    \dfrac{n!}{k!(n-k)!},&0\le k\le n,\\[.8em]
    0,&k<0\text{ oppure }k>n.
  \end{cases}
\]
Quando $0\le k\le n$, esso coincide con $C_{n,k}$. Il valore zero
fuori da questo intervallo consente di scrivere le identità senza
dover separare continuamente i casi agli estremi.

\begin{proposizione}
Per ogni $n\in\N$ valgono le seguenti proprietà:
\begin{enumerate}
  \item $\displaystyle\binom n0=\binom nn=1$;
  \item $\displaystyle\binom nk=0$ se $k<0$ oppure $k>n$;
  \item $\displaystyle\binom nk=\binom n{n-k}$ per ogni $k\in\Z$;
  \item per ogni $k\in\Z$ vale la \emph{relazione di Pascal}:
    \[
      \binom{n+1}k=\binom nk+\binom n{k-1}.
    \]
\end{enumerate}
\end{proposizione}

\begin{proof}
Le prime due proprietà seguono subito dalla definizione e da $0!=1$.
Per $0\le k\le n$, la simmetria si verifica scambiando i fattoriali
al denominatore:
\[
  \binom n{n-k}
   =\frac{n!}{(n-k)!\,[n-(n-k)]!}
   =\frac{n!}{(n-k)!\,k!}=\binom nk.
\]
Se $k$ è fuori dall'intervallo $[0,n]$, lo è anche $n-k$ e i due
coefficienti sono entrambi zero.

Per la relazione di Pascal consideriamo prima $1\le k\le n$ e
portiamo le due frazioni allo stesso denominatore:
\begin{align*}
  \binom nk+\binom n{k-1}
  &=\frac{n!}{k!(n-k)!}
    +\frac{n!}{(k-1)!(n+1-k)!}\\
  &=\frac{n!\bigl[(n+1-k)+k\bigr]}{k!(n+1-k)!}\\
  &=\frac{(n+1)!}{k!(n+1-k)!}
   =\binom{n+1}k.
\end{align*}
Dividendo l'identità per il fattore positivo $n!/[(k-1)!(n-k)!$, la si riduce a
\[
  \frac{n+1}{k(n+1-k)}
   =\frac1k+\frac1{n+1-k}
   =\frac{n+1-k+k}{k(n+1-k)}.
\]
Per $k=0$ e $k=n+1$ la formula diventa rispettivamente $1=1+0$
e $1=0+1$; per gli altri indici esterni entrambi i membri sono zero.
\end{proof}

La simmetria ha anche un significato combinatorio: scegliere $k$
elementi equivale a scegliere gli $n-k$ da escludere. Analogamente,
per scegliere $k$ elementi fra $n+1$ si può fissare un elemento e
distinguere le scelte che non lo contengono, in numero $\binom nk$,
da quelle che lo contengono, in numero $\binom n{k-1}$.

\section{Il triangolo di Tartaglia}
% Fonte: Lezione 3, p.~3

Scrivendo in ogni riga i coefficienti $\binom n0,\binom n1,\ldots,
\binom nn$ si ottiene il \emph{triangolo di Tartaglia}. Le estremità
valgono $1$ e ciascun elemento interno è la somma dei due che lo
precedono nella riga superiore: è proprio la relazione di Pascal.

\begin{schema}
\begin{tikzpicture}[x=1.05cm,y=.85cm]
  \foreach \riga/\valori in {0/{1},1/{1,1},2/{1,2,1},3/{1,3,3,1},4/{1,4,6,4,1}} {
    \foreach \valore [count=\colonna from 0] in \valori {
      \node (t-\riga-\colonna) at ({\colonna-\riga/2},-\riga) {$\valore$};
    }
    \node[anchor=east,black!65] at (-3,-\riga) {$n=\riga$};
  }
  \draw[coloreschema,-{Stealth[length=1.8mm]},shorten <=3pt,shorten >=3pt]
    (t-3-1) -- (t-4-2);
  \draw[coloreschema,-{Stealth[length=1.8mm]},shorten <=3pt,shorten >=3pt]
    (t-3-2) -- (t-4-2);
  \node[coloreschema,anchor=west] at (2.9,-3.7) {$3+3=6$};
\end{tikzpicture}
\par\smallskip
Ogni riga contiene i coefficienti binomiali di un dato ordine $n$.
Le frecce mostrano come due coefficienti della riga $n=3$ generano
quello centrale della riga $n=4$.
\end{schema}

\section{La formula del binomio di Newton}
% Fonte: Lezione 3, p.~4

\begin{proposizione}[Binomio di Newton]
Per $a,b\in\R$ e $n\in\N$ vale
\[
  (a+b)^n=\sum_{k=0}^{n}\binom nk a^{n-k}b^k.
\]
I fattori con esponente zero si interpretano come $1$, secondo
la convenzione usata nelle identità polinomiali.
\end{proposizione}

Il termine di indice $k$ contiene $k$ fattori $b$ e $n-k$ fattori $a$.
Per ottenerlo dal prodotto di $n$ parentesi $(a+b)$ si scelgono le
$k$ parentesi dalle quali prendere $b$: questo spiega la presenza
del coefficiente $\binom nk$.

\begin{proof}
Procediamo per induzione. Per
$n=0$ la formula vale, perché entrambi i membri sono $1$ con la
convenzione appena precisata. Supponiamo ora che valga per $n$.
Moltiplicando per $a+b$ e distribuendo i due fattori otteniamo
\begin{align*}
  (a+b)^{n+1}
  &=\sum_{k=0}^{n}\binom nk a^{n+1-k}b^k
    +\sum_{k=0}^{n}\binom nk a^{n-k}b^{k+1}\\
  &=\sum_{k=0}^{n}\binom nk a^{n+1-k}b^k
    +\sum_{k=1}^{n+1}\binom n{k-1}a^{n+1-k}b^k.
\end{align*}
Nella seconda somma abbiamo cambiato indice, in modo che le potenze
di $a$ e $b$ abbiano la stessa forma. Per gli indici $1\le k\le n$
il coefficiente complessivo è allora
$\binom nk+\binom n{k-1}=\binom{n+1}k$.
Ai due estremi restano $a^{n+1}$ e $b^{n+1}$, entrambi con coefficiente
$1$. Si ricompone quindi esattamente
$\sum_{k=0}^{n+1}\binom{n+1}k a^{n+1-k}b^k$, che è la formula al passo
successivo. Il principio di induzione conclude la dimostrazione.
\end{proof}

Per esempio, dalla riga $n=4$ del triangolo si legge subito
\[
  (a+b)^4=a^4+4a^3b+6a^2b^2+4ab^3+b^4.
\]

% !TeX root = ../main.tex
% Fonti lette sulle immagini: Lezione 3.pdf, pp. 5--7;
% Lezione 4.pdf, pp. 2--14; Lezione 5.pdf, pp. 2--11.
% Precisazione rispetto alla formulazione sintetica del manoscritto:
% il gruppo moltiplicativo dei razionali esclude lo zero.
\chapter{Sistemi numerici e completezza}
\label{cap:sistemi-numerici}
\fonti{Lezioni 3, 4 e 5}

\section{Gli interi e la struttura di gruppo}
% Fonte: Lezione 3, p.~5

Gli interi estendono i naturali aggiungendo i numeri negativi:
\[
  \Z=\{0,\pm1,\pm2,\pm3,\ldots\}.
\]
Il vantaggio fondamentale è che ogni elemento ammette un opposto.
Nei naturali, per esempio, non esiste una soluzione di $1+y=0$;
negli interi la soluzione è $y=-1$.

La somma è un'operazione interna a $\Z$, cioè sommando due interi
si ottiene ancora un intero, e soddisfa:
\begin{enumerate}
  \item \emph{associatività}: $(x+y)+z=x+(y+z)$ per ogni $x,y,z\in\Z$;
  \item \emph{esistenza dell'elemento neutro}: $x+0=0+x=x$ per ogni $x\in\Z$;
  \item \emph{esistenza dell'opposto}: per ogni $x\in\Z$ esiste $-x\in\Z$
    tale che $x+(-x)=(-x)+x=0$.
\end{enumerate}
Queste proprietà esprimono che $(\Z,+)$ è un \emph{gruppo}.
Poiché vale anche la commutatività $x+y=y+x$, si tratta di un
gruppo \emph{commutativo}, detto anche \emph{abeliano}.

\section{I razionali come campo ordinato}
% Fonte: Lezione 3, p.~6

I numeri razionali sono quelli rappresentabili mediante una frazione
di interi con denominatore non nullo:
\[
  \Q=\left\{\frac pq:p,q\in\Z,\ q\ne0\right\}.
\]
Frazioni diverse possono rappresentare lo stesso numero:
$p/q=r/s$ se e solo se $ps=rq$, con $q,s\ne0$. Questo distingue
il numero razionale dalla particolare scrittura usata per esprimerlo.

La somma e il prodotto di razionali sono ancora razionali.
La somma rende $(\Q,+)$ un gruppo abeliano, come accade per gli interi.
Il prodotto rende invece $(\Q\setminus\{0\},\cdot)$ un gruppo abeliano:
l'elemento neutro è $1$ e l'inverso moltiplicativo di un razionale
non nullo $p/q$ è $q/p$. Lo zero è escluso dal gruppo moltiplicativo
perché non ha inverso. Inoltre il prodotto è distributivo rispetto
alla somma:
\[
  x(y+z)=xy+xz.
\]
Queste proprietà fanno di $(\Q,+,\cdot)$ un \emph{campo}.

L'ordine usuale $\le$ è totale anche su $\Q$ ed è compatibile con
le operazioni: sommare lo stesso numero conserva le disuguaglianze,
e il prodotto di due numeri non negativi è non negativo. Perciò
si parla di \emph{campo ordinato}.

\section{La proprietà di Archimede nei razionali}
% Fonte: Lezione 3, p.~6

Per ogni coppia di razionali $x,y>0$ esiste un naturale $n$ tale che
\[
  nx>y.
\]
Questa è la \emph{proprietà di Archimede}: per quanto piccolo sia
$x>0$ rispetto a $y$, un multiplo intero di $x$ riesce a superare $y$.

Nei razionali la proprietà si può verificare direttamente. Scriviamo
$y/x=p/q$, con $p,q$ interi positivi. Poiché $q\ge1$, si ha
$p/q\le p<p+1$. Basta allora scegliere il naturale $n=p+1$:
$n>y/x$ e, moltiplicando per $x>0$, si ottiene $nx>y$.

\section{La densità fra razionali}
% Fonte: Lezione 3, p.~7

Fra due razionali distinti esistono infiniti altri razionali:
\[
  \forall x,y\in\Q,\quad x<y
  \quad\Longrightarrow\quad
  \text{esistono infiniti }z\in\Q\text{ con }x<z<y.
\]
Per trovare un primo esempio basta la media $(x+y)/2$, che è ancora
razionale e si trova strettamente fra $x$ e $y$. Per verificare anche
che i valori possibili sono infiniti, consideriamo
\[
  z_m=x+\frac{y-x}{m+1},\qquad m\in\N^*.
\]
Sono tutti razionali, tutti distinti e tutti interni a $(x,y)$:
infatti $0<(y-x)/(m+1)<y-x$. A differenza dei naturali o degli interi,
i razionali non hanno quindi due elementi consecutivi.

\section{La numerabilità dei razionali}
% Fonte: Lezione 3, p.~7

Un insieme infinito si dice \emph{numerabile} quando i suoi elementi
si possono mettere in corrispondenza biunivoca con i naturali:
equivalentemente, si possono elencare senza omissioni e senza
ripetizioni in una lista indicizzata da $\N$.
Nonostante la densità appena vista, $\Q$ è numerabile.

Per costruire una lista cominciamo dai razionali positivi. Disponiamo
le frazioni $p/q$, con $p,q\in\N^*$, in una griglia: il numeratore
cambia lungo le colonne e il denominatore lungo le righe.
Percorriamo le diagonali di somma $p+q$ costante, partendo da $p+q=2$,
poi $3$, poi $4$, e così via. Su ciascuna diagonale prendiamo i
numeratori in ordine crescente.

\begin{schema}
\begin{tikzpicture}[x=1.45cm,y=.9cm]
  \foreach \numeratore in {1,...,4} {
    \node[black!65] at ({\numeratore-1},.7) {$p=\numeratore$};
    \foreach \denominatore in {1,...,4} {
      \node[fill=white,inner sep=4pt] (f-\numeratore-\denominatore)
        at ({\numeratore-1},{1-\denominatore})
        {$\dfrac{\numeratore}{\denominatore}$};
    }
  }
  \foreach \denominatore in {1,...,4} {
    \node[anchor=east,black!65] at (-.55,{1-\denominatore}) {$q=\denominatore$};
  }
  \foreach \somma in {3,4,5} {
    \pgfmathtruncatemacro{\massimo}{\somma-2}
    \foreach \numeratore in {1,...,\massimo} {
      \pgfmathtruncatemacro{\denominatore}{\somma-\numeratore}
      \pgfmathtruncatemacro{\successivo}{\numeratore+1}
      \pgfmathtruncatemacro{\precedente}{\denominatore-1}
      \draw[coloreschema,-{Stealth[length=1.7mm]},shorten <=2pt,shorten >=2pt]
        (f-\numeratore-\denominatore) -- (f-\successivo-\precedente);
    }
  }
  \draw[coloreschema,line width=.9pt] (.86,-1.21) -- (1.14,-.79);
\end{tikzpicture}
\par\smallskip
Le frecce percorrono le prime diagonali. La frazione $2/2$ si scarta,
perché rappresenta lo stesso numero di $1/1$, già incontrato.
\end{schema}

La lettura comincia quindi con
\[
  \frac11,\quad\frac12,\quad\frac21,\quad\frac13,\quad
  \frac22,\quad\frac31,\quad\frac14,\quad\ldots
\]
e si eliminano ogni volta i numeri già comparsi. Dopo questa
eliminazione, l'inizio della lista è
\[
  1,\quad\frac12,\quad2,\quad\frac13,\quad3,\quad
  \frac14,\quad\frac23,\quad\frac32,\quad4,\quad\ldots
\]
Ogni razionale positivo compare, perché una sua rappresentazione
$p/q$ si trova sulla diagonale $p+q$, raggiunta dopo un numero finito
di passi. Il procedimento non si esaurisce, dato che i razionali
positivi comprendono tutti i naturali positivi. Chiamando la lista
ottenuta $r_1,r_2,r_3,\ldots$, otteniamo infine una lista di tutti
i razionali, senza ripetizioni, nell'ordine
\[
  0,\ r_1,\ -r_1,\ r_2,\ -r_2,\ r_3,\ -r_3,\ldots
\]
e quindi una corrispondenza biunivoca con $\N$. Densità e numerabilità
descrivono dunque due aspetti differenti: la prima riguarda ciò che
si trova fra due numeri, la seconda la possibilità di elencare
tutti gli elementi dell'insieme.

% !TeX root = ../main.tex
% Lezione 4.pdf, pp. 2--5. Le precisazioni sulle sezioni e sui quantificatori
% completano la formulazione sintetica del manoscritto.
\section{I limiti dei razionali: l'irrazionalità della radice di due}
% Fonte: Lezione 4, p.~2

I razionali permettono di eseguire le quattro operazioni, con la sola
esclusione della divisione per zero, ma non risolvono ogni equazione
algebrica. L'esempio fondamentale è $x^2=2$: la sua soluzione positiva,
indicata con $\sqrt2$, non è razionale.

\begin{proposizione}
Non esiste $d\in\Q$ tale che $d^2=2$; in particolare $\sqrt2\notin\Q$.
\end{proposizione}
\begin{proof}
Supponiamo per assurdo che $d=p/q$, con $p,q\in\Z$, $q\ne0$ e
la frazione ridotta ai minimi termini, soddisfi $d^2=2$. Si avrebbe
\[
  p^2=2q^2.
\]
Quindi $p^2$ è pari e anche $p$ deve essere pari: infatti il quadrato
di un dispari $2h+1$ è $4h(h+1)+1$, che è dispari. Scriviamo allora
$p=2k$. Sostituendo nella relazione precedente otteniamo
\[
  4k^2=2q^2,\qquad q^2=2k^2.
\]
Anche $q^2$, e dunque $q$, è pari. Ma allora $p$ e $q$ sono entrambi
divisibili per $2$, in contraddizione con la scelta di una frazione
ridotta ai minimi termini.
\end{proof}

\section{I reali e l'assioma di Dedekind}
% Fonte: Lezione 4, pp.~3--4

Per descrivere la retta numerica si introduce il campo dei numeri reali
$\R$, che contiene $\Q$. La somma e il prodotto sono operazioni interne
e l'ordine $\le$ è totale. Le proprietà algebriche e di compatibilità
con l'ordine sono le seguenti:
\begin{enumerate}
  \item la somma è associativa e commutativa e ha elemento neutro $0$;
  \item ogni $x\in\R$ ha un opposto $-x$, con $x+(-x)=0$;
  \item il prodotto è associativo e commutativo e ha elemento neutro $1$,
    con $1\ne0$;
  \item ogni $x\ne0$ ha un reciproco $x^{-1}$, con $xx^{-1}=1$;
  \item il prodotto è distributivo rispetto alla somma:
    $x(y+z)=xy+xz$;
  \item se $a\le b$, allora $a+c\le b+c$ per ogni $c\in\R$;
  \item se $a\le b$ e $c\ge0$, allora $ac\le bc$.
\end{enumerate}
Queste proprietà fanno di $\R$ un campo ordinato. La proprietà che lo
distingue da $\Q$ è la \emph{completezza}, espressa qui dall'assioma
di Dedekind.

\begin{assioma}[Dedekind]
\label{ass:dedekind}
Siano $A,B\subseteq\R$ non vuoti e tali che $a\le b$ per ogni
$a\in A$ e ogni $b\in B$. Allora esiste $c\in\R$ per cui
\[
  a\le c\le b\qquad\text{per ogni }a\in A\text{ e ogni }b\in B.
\]
\end{assioma}

Il numero $c$ separa i due insiemi: si trova a destra di tutti gli
elementi di $A$ e a sinistra di tutti quelli di $B$. L'ordine dei
quantificatori richiede un valore di $c$ che funzioni
contemporaneamente per tutte le coppie:
\[
  \exists c\in\R\ \forall a\in A\ \forall b\in B:\ a\le c\le b.
\]
L'assioma afferma l'esistenza, senza richiedere che il separatore
sia determinato in modo univoco. Per esempio,
se $A=\{0\}$ e $B=\{1\}$, ogni $c\in[0,1]$ è un separatore.

\subsection{Perché l'assioma fallisce nei razionali}

Consideriamo gli insiemi
\[
  A=\{x\in\Q:x\ge0,\ x^2\le2\},\qquad
  B=\{x\in\Q:x\ge0,\ x^2\ge2\}.
\]
Sono entrambi non vuoti: $1\in A$ e $2\in B$. Inoltre $a\le b$
per ogni $a\in A$ e $b\in B$, perché sono non negativi e
$a^2\le2\le b^2$. Tuttavia nessun razionale può separare $A$ e $B$.

Per verificarlo, un eventuale separatore razionale $c$ dovrebbe
soddisfare $1\le c\le2$. Se $c^2<2$, scegliamo un razionale
\[
  0<h<\min\left\{1,\frac{2-c^2}{2c+1}\right\}.
\]
Allora $(c+h)^2<c^2+(2c+1)h<2$, quindi $c+h\in A$ ma $c+h>c$:
il separatore non sarebbe a destra di tutti gli elementi di $A$.
Se invece $c^2>2$, scegliamo un razionale
\[
  0<h<\min\left\{1,\frac{c^2-2}{2c}\right\}.
\]
Si ha $c-h>0$ e $(c-h)^2>c^2-2ch>2$, perciò $c-h\in B$ ma
$c-h<c$, ancora una contraddizione. Rimane soltanto la possibilità
$c^2=2$, esclusa dall'irrazionalità appena dimostrata. In $\R$ il
separatore è il punto $\sqrt2$ che manca a $\Q$.

\subsection{L'idea delle sezioni di Dedekind}
% Fonte: Lezione 4, p.~5

È possibile costruire i reali a partire dai razionali mediante le
\emph{sezioni di Dedekind}. Si considera una coppia $(A,B)$ di sottoinsiemi non vuoti
di $\Q$ che soddisfi
\[
  A\cap B=\varnothing,\qquad A\cup B=\Q,\qquad
  a\le b\quad\forall a\in A,\ \forall b\in B.
\]
La coppia descrive un \emph{taglio} della retta razionale in una parte
inferiore e una superiore. Poiché le due parti sono disgiunte,
in realtà $a<b$ per ogni $a\in A$, $b\in B$.

Per rendere la rappresentazione univoca occorre fissare una convenzione
sull'estremo: scegliamo che $A$ non abbia massimo. Infatti, per un
razionale $r$, sia la coppia
$\bigl(\{q:q<r\},\{q:q\ge r\}\bigr)$ sia la coppia
$\bigl(\{q:q\le r\},\{q:q>r\}\bigr)$ soddisfano le tre condizioni
precedenti, pur descrivendo lo stesso punto. La convenzione sceglie
la prima. Con questa precisazione, tali sezioni possono essere usate,
con opportune operazioni e ordine, per costruire un modello di $\R$.
Un taglio può rappresentare anche un punto irrazionale: in quel caso
nessuna delle due parti contiene un estremo razionale nel punto
di separazione.

% !TeX root = ../main.tex
% Lezione 4.pdf, pp. 6--14. Le prove dei criteri operativi e della
% proprietà duale dell'infimo rendono espliciti i passaggi della fonte.
\section{Intervalli e retta reale estesa}
% Fonte: Lezione 4, pp.~6--8

\begin{definizione}[Intervallo]
Un sottoinsieme non vuoto $I\subseteq\R$ è un \emph{intervallo} se,
ogni volta che $x,y\in I$ e $x<z<y$, anche $z$ appartiene a $I$.
\end{definizione}

In un intervallo si trovano quindi tutti i punti compresi fra due
suoi elementi. Per estremi finiti $a<b$ le principali possibilità sono
\begin{align*}
  [a,b]&=\{x\in\R:a\le x\le b\},&
  (a,b)&=\{x\in\R:a<x<b\},\\
  [a,b)&=\{x\in\R:a\le x<b\},&
  (a,b]&=\{x\in\R:a<x\le b\}.
\end{align*}
I primi due sono detti rispettivamente intervallo \emph{chiuso} e
\emph{aperto}; negli altri due è incluso un solo estremo. Anche
$[a,a]=\{a\}$ è un intervallo.\footnote{La definizione richiede
che un intervallo sia non vuoto. Altre convenzioni ammettono anche
l'insieme vuoto come intervallo.}

Per $I=[a,b]$, i numeri $a,b$ sono gli \emph{estremi}, $b-a$ è
l'\emph{ampiezza}, $(b-a)/2$ il \emph{raggio} e $(a+b)/2$ il
\emph{centro}. Le stesse espressioni individuano ampiezza e centro
anche se uno o entrambi gli estremi non sono inclusi.

Per descrivere intervalli illimitati introduciamo i simboli
$-\infty$ e $+\infty$, ordinati in modo che
\[
  -\infty<x<+\infty\qquad\forall x\in\R.
\]
La \emph{retta reale estesa} è
$\overline\R=\R\cup\{-\infty,+\infty\}$. Le semirette della retta
reale si scrivono, per esempio,
\begin{align*}
  (-\infty,c]&=\{x\in\R:x\le c\},&
  (-\infty,c)&=\{x\in\R:x<c\},\\
  [a,+\infty)&=\{x\in\R:x\ge a\},&
  (a,+\infty)&=\{x\in\R:x>a\}.
\end{align*}
In queste scritture $\pm\infty$ non appartengono agli intervalli,
che sono sottoinsiemi di $\R$; anche tutta $\R=(-\infty,+\infty)$
è un intervallo.

\section{Limitatezza, maggioranti e minoranti}
% Fonte: Lezione 4, pp.~8--10

\begin{definizione}
Un insieme non vuoto $A\subseteq\R$ è \emph{superiormente limitato}
se esiste $M\in\R$ tale che $a\le M$ per ogni $a\in A$.
Un tale numero è un \emph{maggiorante} di $A$.
È \emph{inferiormente limitato} se esiste $m\in\R$ tale che
$m\le a$ per ogni $a\in A$; un tale numero è un \emph{minorante}.
\end{definizione}

L'insieme è \emph{limitato} se lo è sia superiormente sia inferiormente:
\[
  \exists m,M\in\R:\quad m\le a\le M\quad\forall a\in A.
\]
Equivalentemente, esiste $N\ge0$ per cui $|a|\le N$ per ogni
$a\in A$. Infatti, da due estremi $m,M$ basta scegliere
$N=\max\{|m|,|M|\}$; viceversa, $|a|\le N$ equivale a
$-N\le a\le N$. Un insieme che non è limitato si dice
\emph{illimitato}; si può anche specificare la direzione in cui
manca un limite, parlando di insieme superiormente o inferiormente
illimitato.

Maggioranti e minoranti non devono appartenere all'insieme e non
sono determinati in modo univoco. Per esempio, sono possibili scelte
anche molto lontane dagli estremi:
\begin{itemize}
  \item per $A=[0,1)$ si possono prendere $m=-38$ e $M=251$;
  \item per $B=\{1/n:n\in\N^*\}$ si possono prendere
    $m=-1719{,}3$ e $M=16$, perché $0<1/n\le1$;
  \item per $C=\N$, $0$ è un minorante ma non esistono maggioranti,
    come seguirà dal principio di Archimede
    (teorema~\ref{teo:archimede-reali}).
\end{itemize}

Per $A=[0,3)$, in particolare, sono maggioranti $3$, $5$, $11$,
$197$ e, più in generale, tutti e soli i reali $M\ge3$.
Sono minoranti $0$, $-13$, $-\pi$, $-171$ e, in generale, tutti e soli
i reali $m\le0$. La condizione vale per ogni elemento di $A$,
anche se $3\notin A$.

\section{Massimo e minimo}
% Fonte: Lezione 4, p.~11

\begin{definizione}
Un \emph{massimo} di un insieme non vuoto $A\subseteq\R$ è un
maggiorante che appartiene ad $A$; un \emph{minimo} è un minorante
che appartiene ad $A$. Si scrive
\begin{align*}
  M=\max A&\quad\Longleftrightarrow\quad
    M\in A\text{ e }a\le M\ \forall a\in A,\\
  m=\min A&\quad\Longleftrightarrow\quad
    m\in A\text{ e }m\le a\ \forall a\in A.
\end{align*}
\end{definizione}

Se esistono, massimo e minimo sono unici. Per esempio, se $M_1$ e
$M_2$ fossero entrambi massimi, l'appartenenza di ciascuno ad $A$
imporrebbe $M_1\le M_2$ e $M_2\le M_1$, quindi $M_1=M_2$.
Per il minimo si ragiona allo stesso modo.

L'intervallo $[0,3)$ ha minimo $0$ e non ha massimo: per ogni
$a\in[0,3)$ il punto $(a+3)/2$ è ancora nell'intervallo ed è
strettamente maggiore di $a$. Quindi l'essere limitato non assicura
che un insieme possieda un massimo o un minimo.

\section{Estremo superiore ed estremo inferiore}
% Fonte: Lezione 4, p.~12

\begin{definizione}[Supremo e infimo]
Il numero $S\in\R$ è l'\emph{estremo superiore}, o \emph{supremo},
di $A$ se è il più piccolo dei suoi maggioranti. Il numero $s\in\R$
è l'\emph{estremo inferiore}, o \emph{infimo}, se è il più grande
dei suoi minoranti. In formule,
\begin{align*}
  S=\sup A&=\min\{M\in\R:a\le M\ \forall a\in A\},\\
  s=\inf A&=\max\{m\in\R:m\le a\ \forall a\in A\},
\end{align*}
quando tali estremi esistono.
\end{definizione}

Anche questi valori, quando esistono, sono unici. A differenza del
massimo e del minimo, possono non appartenere ad $A$. Se un massimo
esiste, è anche il più piccolo maggiorante: essendo un elemento di
$A$, è minore o uguale a qualunque altro maggiorante. Pertanto
\[
  \max A=\sup A\quad\text{se il massimo esiste},\qquad
  \min A=\inf A\quad\text{se il minimo esiste}.
\]
Viceversa, se $\sup A\in A$ esso è il massimo, e se $\inf A\in A$
esso è il minimo.

\subsection{Il teorema di esistenza degli estremi}
% Fonte: Lezione 4, pp.~13--14

\begin{teorema}[Esistenza del supremo e dell'infimo]
\label{teo:esistenza-estremi}
Ogni sottoinsieme non vuoto di $\R$ superiormente limitato ha
un estremo superiore reale. Ogni sottoinsieme non vuoto di $\R$
inferiormente limitato ha un estremo inferiore reale.
\end{teorema}
\begin{proof}
Sia $A\subseteq\R$ non vuoto e superiormente limitato. L'insieme
dei suoi maggioranti
\[
  \mathcal M_A=\{M\in\R:a\le M\ \forall a\in A\}
\]
è non vuoto. Per definizione, ogni elemento di $A$ è minore o
uguale a ogni elemento di $\mathcal M_A$. Applichiamo l'assioma
di Dedekind ai due insiemi $A$ e $\mathcal M_A$: esiste $S\in\R$
tale che
\[
  a\le S\le M\qquad\forall a\in A,\ \forall M\in\mathcal M_A.
\]
La prima disuguaglianza dice che $S$ è un maggiorante di $A$;
la seconda che è minore o uguale a ogni maggiorante. Dunque
$S=\min\mathcal M_A=\sup A$.

Per l'infimo si può applicare lo stesso ragionamento all'insieme
$-A=\{-a:a\in A\}$. Se $A$ è inferiormente limitato, $-A$ è
superiormente limitato e ha supremo $T$. Cambiando segno alle
disuguaglianze si ottiene che $-T$ è il più grande minorante di
$A$, ossia $\inf A=-\sup(-A)$.
\end{proof}

Il teorema è una conseguenza della completezza. Sostituendo $\R$
con $\Q$ l'enunciato diventa falso: l'insieme
\[
  A=\{x\in\Q:x\ge0,\ x^2\le2\}
\]
è non vuoto e superiormente limitato, ma non ha supremo in $\Q$.
Un eventuale minimo maggiorante razionale sarebbe infatti un
separatore razionale fra gli insiemi $A$ e $B$ già considerati
nell'esempio sull'assioma di Dedekind. Considerando lo stesso
insieme come sottoinsieme di $\R$, il suo supremo è $\sqrt2$;
il suo infimo è $0$, che è anche minimo.

\subsection{Estremi infiniti}

Per gli insiemi non vuoti si estende la notazione ponendo
$\sup A=+\infty$ quando $A$ è superiormente illimitato e
$\inf A=-\infty$ quando è inferiormente illimitato.
Si tratta di valori nella retta estesa, non di massimi o minimi
reali. Per esempio, $\sup\N=+\infty$ e $\inf\N=\min\N=0$.

\section{La caratterizzazione operativa di supremo e infimo}
% Fonte: Lezione 4, p.~14; lezione 5, p.~2

\begin{proposizione}
\label{prop:criterio-estremi}
Sia $A\subseteq\R$ non vuoto. Un numero $S\in\R$ è il suo supremo
se e solo se valgono entrambe le condizioni
\[
  \begin{cases}
    a\le S&\forall a\in A,\\
    \forall\varepsilon>0\ \exists a\in A:\ a>S-\varepsilon.
  \end{cases}
\]
Analogamente, $s\in\R$ è l'infimo se e solo se
\[
  \begin{cases}
    s\le a&\forall a\in A,\\
    \forall\varepsilon>0\ \exists a\in A:\ a<s+\varepsilon.
  \end{cases}
\]
\end{proposizione}
\begin{proof}
Se $S=\sup A$, esso è un maggiorante. Se per qualche
$\varepsilon>0$ non esistesse $a\in A$ con $a>S-\varepsilon$,
tutti gli elementi sarebbero minori o uguali a $S-\varepsilon$:
quest'ultimo sarebbe un maggiorante più piccolo di $S$, assurdo.

Viceversa, supponiamo che le due condizioni valgano. $S$ è un
maggiorante e ogni $M<S$ non può esserlo: scegliendo
$\varepsilon=S-M>0$ si trova $a\in A$ con $a>M$. Quindi $S$ è
il più piccolo dei maggioranti.

Per l'infimo, se $s=\inf A$ e mancassero elementi sotto
$s+\varepsilon$, quest'ultimo sarebbe un minorante più grande.
Viceversa, se le due condizioni per $s$ valgono, ogni $m>s$ viene
escluso dai minoranti scegliendo $\varepsilon=m-s$ e trovando
$a\in A$ con $a<m$. Dunque $s$ è il più grande minorante.
\end{proof}

Il criterio traduce il ruolo di estremo in due verifiche separate:
essere un limite per tutti gli elementi e poter essere avvicinato
arbitrariamente dall'insieme. L'elemento $a$ può dipendere da
$\varepsilon$. Se l'estremo appartiene ad $A$, lo si può scegliere
per ogni $\varepsilon$; il criterio è particolarmente utile quando
l'estremo non è raggiunto.

\begin{schema}
\begin{tikzpicture}[x=1cm,y=1cm]
  \fill[coloreschema!12] (4.3,.52) rectangle (6.8,.78);
  \fill[coloreschema!12] (1.2,-.98) rectangle (3.7,-.72);
  \foreach \quota in {.65,-.85} {
    \draw[asse] (0,\quota) -- (8,\quota);
  }
  \node[anchor=east] at (-.2,.65) {$S=\sup A$};
  \node[anchor=east] at (-.2,-.85) {$s=\inf A$};
  \foreach \pos/\testo in {4.3/{S-\varepsilon},6.8/S} {
    \draw (\pos,.55) -- (\pos,.75);
    \node[below=6pt] at (\pos,.65) {$\testo$};
  }
  \foreach \pos/\testo in {1.2/s,3.7/{s+\varepsilon}} {
    \draw (\pos,-.95) -- (\pos,-.75);
    \node[below=6pt] at (\pos,-.85) {$\testo$};
  }
  \fill[coloreschema] (6.1,.65) circle[radius=2pt];
  \node[above=6pt,coloreschema] at (6.1,.65) {$a\in A$};
  \fill[coloreschema] (2.1,-.85) circle[radius=2pt];
  \node[above=6pt,coloreschema] at (2.1,-.85) {$a\in A$};
\end{tikzpicture}
\par\smallskip
Per ogni $\varepsilon>0$, la fascia $(S-\varepsilon,S]$ deve
contenere almeno un punto di $A$; analogamente per $[s,s+\varepsilon)$.
Il colore evidenzia la fascia di ricerca, senza affermare che tutti
i suoi punti appartengano ad $A$.
\end{schema}

% !TeX root = ../main.tex
% Lezione 5.pdf, pp. 2--11. La proprietà di Archimede precede gli esempi
% che la usano. La parte sulle funzioni, pp. 12--16, è nel capitolo 05.
% Corretto il passaggio del manoscritto su 1/n < -1+epsilon:
% per 0<epsilon<1 il secondo membro è negativo e non vi sono soluzioni.
\section{Il principio di Archimede nei reali}
% Fonte: Lezione 5, pp.~8--9

\begin{teorema}[Principio di Archimede]
\label{teo:archimede-reali}
Fissati $a,b\in\R$, con $a>0$, esiste $n\in\N$ tale che $na>b$.
\end{teorema}
\begin{proof}
\textbf{Riduzione al supremo dei multipli.} Fissiamo $a>0$ e consideriamo l'insieme
\[
  A=\{na:n\in\N\}.
\]
È non vuoto, perché $0\in A$. Dimostriamo che è superiormente illimitato. Se per assurdo fosse superiormente limitato, il teorema di esistenza degli estremi garantirebbe un numero reale $S=\sup A$.

\textbf{Il successivo supera il supremo.} La caratterizzazione operativa, con la scelta $\varepsilon=a>0$, darebbe un elemento $\bar n a\in A$ per cui
\[
  \bar n a>S-a.
\]
Sommando $a$ ai due membri otteniamo $(\bar n+1)a>S$. Ma $\bar n+1\in\N$, quindi anche $(\bar n+1)a\in A$, in contraddizione con il fatto che $S$ sia un maggiorante. L'insieme dei multipli è dunque superiormente illimitato: per ogni $b\in\R$ contiene un elemento $na>b$, come richiesto.
\end{proof}

\textbf{Conseguenze.}
\begin{itemize}
  \item Con $a=1$ si ottiene che i naturali superano qualunque reale: $\N$ è superiormente illimitato.
  \item I reciproci dei naturali positivi possono essere arbitrariamente piccoli:
  \[
    \forall\varepsilon>0\ \exists n\in\N^*:\quad \frac1n<\varepsilon.
  \]
  Basta infatti scegliere $n>1/\varepsilon$.
\end{itemize}
La completezza ha quindi una conseguenza concreta: non esiste un reale positivo minore o uguale a tutti i reciproci dei naturali positivi.

\addtocontents{toc}{\protect\Needspace{3\protect\baselineskip}}
\section{Esempi ed esercizi su supremo e infimo}
% Fonte: Lezione 5, pp.~2--7

\subsection{Un intervallo e un insieme con un punto isolato}
% Fonte: Lezione 5, p.~2

\begin{esempio}[Intervallo con estremo escluso]
Per $A=[0,1)$ si ha
\[
  \inf A=\min A=0,\qquad \sup A=1.
\]
\textbf{Minimo.} Il minimo è $0$, perché $0\in A$ e tutti gli elementi sono non negativi.

\textbf{Supremo non raggiunto.} Il numero $1$ è un maggiorante e, dato $\varepsilon>0$, il punto $a=1-\min\{\varepsilon/2,1/2\}$ appartiene ad $A$ e soddisfa $a>1-\varepsilon$. Il criterio operativo dimostra dunque che $1=\sup A$. Tuttavia $1\notin A$, quindi il massimo non esiste.
\end{esempio}

\begin{esempio}[Un massimo isolato]
Consideriamo
\[
  B=[-1,2]\cup\{5\}.
\]
In questo caso entrambi gli estremi sono raggiunti:
\[
  \inf B=\min B=-1,\qquad \sup B=\max B=5.
\]
Per verificare il criterio del supremo basta scegliere $a=5$ per ogni $\varepsilon>0$, poiché $5>5-\varepsilon$. Il punto che realizza il massimo può dunque essere isolato; la presenza dell'intero intervallo fra $2$ e $5$ non è richiesta.
\end{esempio}

\begin{schema}
\begin{tikzpicture}[x=1.45cm,y=1cm]
  \foreach \quota/\nome in {.7/A,-.8/B} {
    \draw[asse] (-1.5,\quota) -- (5.5,\quota);
    \node[anchor=east] at (-1.8,\quota) {$\nome$};
  }
  \draw[intervallo] (0,.7) -- (1,.7);
  \fill[coloreschema] (0,.7) circle[radius=2.5pt];
  \draw[estremo] (1,.7) circle[radius=2.5pt];
  \node[below=6pt] at (0,.7) {$0$};
  \node[below=6pt] at (1,.7) {$1$};
  \draw[intervallo] (-1,-.8) -- (2,-.8);
  \foreach \pos in {-1,2,5} {
    \fill[coloreschema] (\pos,-.8) circle[radius=2.5pt];
    \node[below=6pt] at (\pos,-.8) {$\pos$};
  }
\end{tikzpicture}
\par\smallskip In $A$ l'estremo superiore $1$ è escluso; in $B$ il punto $5$ è presente ed è il massimo, anche se è separato dall'intervallo.
\end{schema}

\subsection{I reciproci dei naturali positivi}
% Fonte: Lezione 5, pp.~3--4

\begin{esempio}[Reciproci dei naturali positivi]
Sia
\[
  A=\left\{\frac1n:n\in\N^*\right\}.
\]
La rappresentazione dei suoi punti suggerisce $\sup A=1$ e $\inf A=0$. Verifichiamo entrambe le affermazioni.

\textbf{Massimo e supremo.} Il numero $1$ appartiene ad $A$ per $n=1$ e maggiora tutti gli elementi, poiché $1/n\le1$ per ogni $n\ge1$. Dunque
\[
  1=\max A=\sup A.
\]
\textbf{Infimo non raggiunto.} Il numero $0$ non appartiene ad $A$: nessun reciproco di un naturale positivo è zero. Per dimostrare che è l'infimo usiamo le due condizioni operative. È un minorante, perché $1/n>0$; inoltre, per ogni $\varepsilon>0$, Archimede fornisce un naturale $n>1/\varepsilon$, e quindi
\[
  \frac1n<\varepsilon=0+\varepsilon.
\]
Ne segue $\inf A=0$. Poiché $0\notin A$, il minimo non esiste.
\end{esempio}

\subsection{Come scartare un candidato sbagliato per l'infimo}
% Fonte: Lezione 5, pp.~5--6

Per lo stesso insieme $A=\{1/n:n\in\N^*\}$ proviamo a proporre un candidato $k$ diverso da $0$.

\textbf{Candidato positivo: non è un minorante.} Se $k>0$, $k$ non è un minorante. Infatti esiste $n>1/k$, quindi $1/n<k$.

\textbf{Candidato negativo: manca l’avvicinabilità.} Se $k<0$, esso è un minorante, perché tutti gli elementi di $A$ sono positivi, ma non è il più grande: $0$ è un minorante maggiore di $k$. Lo stesso difetto si vede con il criterio operativo. Per il candidato $k=-1$, prendendo $0<\varepsilon<1$ avremmo bisogno di un elemento per cui
\[
  \frac1n<-1+\varepsilon<0,
\]
che è impossibile. In generale, per un qualunque $k<0$, basta scegliere $\varepsilon=-k/2$: allora $k+\varepsilon=k/2<0$ e nessun elemento di $A$ può trovarsi sotto questa soglia.

\subsection{Due esercizi: cubi traslati e unione di insiemi}
% Fonte: Lezione 5, p.~7

\begin{esercizio}[Cubi traslati]
Determinare supremo e infimo dell'insieme
\[
  A_1=\{n^3-1:n\in\N\}.
\]
Stabilire se esistono massimo e minimo.
\end{esercizio}
\begin{svolgimento}
Poiché $n^3\ge0$ per ogni $n\in\N$ e il valore $-1$ si ottiene per $n=0$, si ha $\inf A_1=\min A_1=-1$.

\textbf{Illimitatezza superiore.} L'insieme è superiormente illimitato: dato un qualunque $M\in\R$, scegliamo un naturale $n>\max\{1,M+1\}$. Allora $n^3-1\ge n-1>M$. Dunque, nella retta estesa,
\[
  \sup A_1=+\infty.
\]
L'illimitatezza superiore esclude l'esistenza di un massimo.
\end{svolgimento}

\begin{esercizio}[Unione di due insiemi]
Determinare supremo, infimo ed eventuali massimo e minimo dell'insieme $A_2=\{x\in\R:x^2<1\}\cup\{x\in\R:(x-3)^2=0\}$.
\end{esercizio}
\begin{svolgimento}
\textbf{Descrizione dell'insieme.} Si ha
\[
  A_2=\{x\in\R:x^2<1\}
      \cup\{x\in\R:(x-3)^2=0\}
     =(-1,1)\cup\{3\}.
\]
Il numero $3$ appartiene ad $A_2$ e maggiora tutti gli elementi, quindi $\sup A_2=\max A_2=3$.

\textbf{Infimo non raggiunto.} Il numero $-1$ è un minorante. Dato $\varepsilon>0$, scegliamo
\[
  a=-1+\frac{\varepsilon}{1+\varepsilon}.
\]
Il punto $a$ appartiene a $(-1,0)\subseteq A_2$ e soddisfa $a<-1+\varepsilon$, perché $0<\varepsilon/(1+\varepsilon)<\varepsilon$. Perciò $\inf A_2=-1$; il minimo non esiste, perché $-1\notin A_2$.
\end{svolgimento}

% Mantiene insieme la dimostrazione e il disegno della densità.
\Needspace{36\baselineskip}
\section{La densità dei razionali nella retta reale}
% Fonte: Lezione 5, pp.~10--11

Fra due razionali distinti si trovano altri razionali. Il risultato seguente estende la proprietà a qualunque intervallo aperto con estremi reali.

\begin{teorema}[Densità di $\Q$ in $\R$]
\label{teo:densita-razionali}
Se $a,b\in\R$ e $a<b$, esiste $r\in\Q$ tale che $a<r<b$.
\end{teorema}
\begin{proof}
\textbf{Dilatazione dell’intervallo.} Poiché $b-a>0$, il principio di Archimede permette di scegliere $q\in\N^*$ tale che
\[
  q(b-a)>1.
\]
L'intervallo $(qa,qb)$ ha dunque ampiezza maggiore di $1$ e contiene un intero $p$.

\textbf{Scelta dell’intero.} Per ogni reale $t$ esiste un intero $p$ con $t<p\le t+1$. Infatti, scegliamo un naturale $N>-t$ e consideriamo il minimo naturale $m>t+N$, che esiste per Archimede e per il principio del minimo intero. Si ha $m\ge1$ e $m-1\le t+N$, altrimenti $m$ non sarebbe il minimo. Ponendo $p=m-N$ si ottiene proprio $t<p\le t+1$.

\textbf{Ritorno alla scala iniziale.} Applicando il fatto appena verificato a $t=qa$, troviamo
\[
  qa<p\le qa+1<qb.
\]
Dividendo per $q>0$ segue $a<p/q<b$. Il numero $r=p/q$ è razionale, perché $p\in\Z$ e $q\in\N^*$, e appartiene all'intervallo richiesto.
\end{proof}

\begin{schema}
\begin{tikzpicture}[x=1.6cm,y=1cm]
  \draw[asse] (-.3,0) -- (4.4,0);
  \draw[intervallo] (1.35,0) -- (3.25,0);
  \draw[estremo] (1.35,0) circle[radius=2.5pt];
  \draw[estremo] (3.25,0) circle[radius=2.5pt];
  \node[above=8pt,coloreschema] at (1.35,0) {$qa$};
  \node[above=8pt,coloreschema] at (3.25,0) {$qb$};
  \foreach \pos/\testo in {0/{p-2},1/{p-1},2/p,3/{p+1},4/{p+2}} {
    \draw[black!55] (\pos,-.08) -- (\pos,.08);
    \node[below=7pt] at (\pos,0) {$\testo$};
  }
  \fill[coloreschema] (2,0) circle[radius=2.5pt];
\end{tikzpicture}
\par\smallskip Dilatando l'intervallo con un fattore intero $q$ si ottiene un'ampiezza maggiore di $1$, che assicura la presenza di un intero $p$. Tornando alla scala iniziale si trova il razionale $p/q$.
\end{schema}

\textbf{Infinità dei punti razionali.} I razionali in $(a,b)$ sono infiniti: trovato un primo razionale $r_1$ nell'intervallo, si può cercarne un secondo in $(a,r_1)$, un terzo fra $a$ e il secondo, e continuare. Ogni intervallo aperto non vuoto della retta incontra quindi $\Q$, anche quando i suoi estremi sono irrazionali.



% !TeX root = ../main.tex
\chapter{Funzioni: definizioni e proprietà generali}
\label{cap:funzioni}
\fonti{Lezioni 5 e 6}

\section{Dominio, codominio e immagine}
% Fonte: Lezione 5, p.~12

\begin{definizione}
Una \emph{funzione} $f\colon X\to Y$ è una legge che associa a ogni
elemento $x\in X$ uno e un solo elemento $y\in Y$. Scriviamo
\[
f\colon X\longrightarrow Y,\qquad x\longmapsto y=f(x).
\]
L'insieme $X$ è il \emph{dominio}, $Y$ il \emph{codominio}; $f(x)$ è
l'\emph{immagine di $x$ tramite $f$}. L'\emph{insieme immagine} è
\[
\operatorname{Im}(f)=f(X)=\{f(x):x\in X\}\subseteq Y.
\]
\end{definizione}

La condizione «uno e un solo» riguarda il valore associato a ciascun
elemento del dominio. Elementi distinti possono avere la stessa immagine;
inoltre un elemento del codominio può non essere raggiunto. Perciò
codominio e immagine non coincidono necessariamente.

\section{Iniettività, suriettività e biettività}
% Fonte: Lezione 5, p.~13

\begin{definizione}
Una funzione $f\colon X\to Y$ si dice:
\begin{itemize}
  \item \emph{iniettiva} se elementi distinti hanno immagini distinte:
  \[
  \forall x_1,x_2\in X,\quad
  x_1\ne x_2\ \Longrightarrow\ f(x_1)\ne f(x_2);
  \]
  equivalentemente, $f(x_1)=f(x_2)$ implica $x_1=x_2$;
  \item \emph{suriettiva} se ogni elemento del codominio è raggiunto:
  \[
  \forall y\in Y\ \exists x\in X:\ f(x)=y,
  \qquad\text{ossia }\operatorname{Im}(f)=Y;
  \]
  \item \emph{biettiva} se è sia iniettiva sia suriettiva.
\end{itemize}
\end{definizione}

Per una biezione, dunque, l'equazione $f(x)=y$ ha esattamente una
soluzione $x\in X$ per ogni $y\in Y$: la suriettività dà l'esistenza,
l'iniettività dà l'unicità.

\begin{esempio}
La funzione $f\colon\R\to\R$, $f(x)=x^2$, non è iniettiva perché
$f(1)=f(-1)=1$, e non è suriettiva perché non assume valori negativi.
Se il codominio diventa $[0,+\infty)$, la funzione è suriettiva.
Restringendo anche il dominio a $[0,+\infty)$ diventa biettiva.
Dominio e codominio sono quindi parte essenziale della descrizione di
una funzione.
\end{esempio}

\section{Composizione e suo dominio}
% Fonte: Lezione 5, pp.~14--15

Siano $f\colon X\to Y$ e $g\colon V\to Z$. Per calcolare $g(f(x))$
occorre che il risultato della prima funzione appartenga al dominio
della seconda. Poniamo pertanto
\[
X'=\{x\in X:f(x)\in V\}.
\]
La \emph{funzione composta} è
\[
g\circ f\colon X'\longrightarrow Z,\qquad
(g\circ f)(x)=g(f(x)).
\]
Se $f(X)\subseteq V$, il dominio della composta è tutto $X$; se c'è
soltanto un'intersezione non vuota fra $f(X)$ e $V$, la composta è
definita sul sottoinsieme $X'$. In assenza di intersezione, $X'$ è vuoto.
L'ordine di lettura è da destra verso sinistra: si applica prima $f$ e
poi $g$.

\begin{schema}
\begin{tikzpicture}[>=Stealth]
  \node (X) at (0,0) {$X'$};
  \node (V) at (4,0) {$V$};
  \node (Z) at (8,0) {$Z$};
  \draw[->,coloreschema] (X) -- node[above] {$f|_{X'}$} (V);
  \draw[->,coloreschema] (V) -- node[above] {$g$} (Z);
  \draw[->] (X) to[bend right=22] node[below] {$g\circ f$} (Z);
\end{tikzpicture}
\end{schema}

\begin{esempio}
Prendiamo $f(x)=1+x^2$, con dominio $\R$, e $g(y)=\sqrt y$, con
dominio $[0,+\infty)$. Si ottiene
\[
(f\circ g)(x)=1+(\sqrt x)^2=1+x,
\qquad \operatorname{dom}(f\circ g)=[0,+\infty),
\]
mentre
\[
(g\circ f)(x)=\sqrt{1+x^2},
\qquad \operatorname{dom}(g\circ f)=\R.
\]
La composizione, in generale, non è commutativa: qui differiscono sia
le formule sia i domini. Anche quando una formula si semplifica, il
dominio della composta resta quello imposto dalle operazioni originarie.
\end{esempio}

\section{Grafico, controimmagini e funzione inversa}
% Fonte: Lezione 5, p.~16; lezione 6, p.~2

\begin{definizione}
Il \emph{grafico} di $f\colon X\to Y$ è l'insieme
\[
\operatorname{graf}(f)=\{(x,f(x)):x\in X\}\subseteq X\times Y.
\]
La \emph{controimmagine}, o \emph{fibra}, di $y\in Y$ è
\[
f^{-1}(\{y\})=\{x\in X:f(x)=y\}.
\]
Più in generale, per $B\subseteq Y$ si pone
$f^{-1}(B)=\{x\in X:f(x)\in B\}$.
\end{definizione}

La controimmagine è un insieme ed è definita per ogni funzione. Se $f$
è iniettiva, la fibra di un elemento $y\in\operatorname{Im}(f)$ contiene
esattamente un elemento: non è vuota perché $y$ appartiene all'immagine,
e non ne può contenere due per l'iniettività. Questo consente di
definire una nuova funzione.

\begin{definizione}
Se $f\colon X\to Y$ è iniettiva, la \emph{funzione inversa} è
\[
f^{-1}\colon\operatorname{Im}(f)\longrightarrow X,
\qquad f^{-1}(y)=\text{l'unico }x\in X\text{ tale che }f(x)=y.
\]
\end{definizione}

La scrittura $f^{-1}(y)$ indica qui un elemento, mentre
$f^{-1}(\{y\})=\{f^{-1}(y)\}$ indica il corrispondente insieme con un
solo elemento. Inoltre $f^{-1}$ non indica il reciproco $1/f$.
L'inversa è definita sull'intero codominio $Y$ se $f$ è biettiva.

\begin{proposizione}
L'inversa di una funzione iniettiva soddisfa
\[
f^{-1}\circ f=\operatorname{Id}_X,
\qquad f\circ f^{-1}=\operatorname{Id}_{\operatorname{Im}(f)},
\]
dove $\operatorname{Id}_A\colon A\to A$, $a\mapsto a$, è la funzione
identità su $A$.
\end{proposizione}
\begin{proof}
Per $x\in X$, l'unico elemento che $f$ manda in $f(x)$ è $x$ stesso;
quindi $f^{-1}(f(x))=x$. Per $y\in\operatorname{Im}(f)$, la definizione
di $f^{-1}(y)$ dà direttamente $f(f^{-1}(y))=y$.
\end{proof}

Quando $X,Y\subseteq\R$, il punto $(x,y)$ appartiene al grafico di $f$
se e solo se $(y,x)$ appartiene al grafico dell'inversa: i due grafici
sono simmetrici rispetto alla bisettrice $y=x$.

\begin{schema}
\begin{tikzpicture}[x=1cm,y=1cm]
  \draw[asse] (-.3,0) -- (4.7,0) node[right] {$x$};
  \draw[asse] (0,-.3) -- (0,4.7) node[above] {$y$};
  \draw[dashed,black!45] (0,0) -- (4.4,4.4) node[above right] {$y=x$};
  \draw[intervallo,domain=0:2.05,samples=60] plot (\x,{\x*\x});
  \draw[thick,orange!80!black,domain=0:2.05,samples=60]
    plot ({\x*\x},\x);
  \node[coloreschema,left] at (1.8,3.7) {$y=x^2,\ x\ge0$};
  \node[orange!80!black,below] at (3.4,1.7) {$y=\sqrt x$};
  \fill[coloreschema] (1.5,2.25) circle (1.5pt);
  \fill[orange!80!black] (2.25,1.5) circle (1.5pt);
  \draw[dotted] (1.5,2.25) -- (2.25,1.5);
\end{tikzpicture}
\end{schema}

% !TeX root = ../main.tex
\section{Funzioni reali e dominio naturale}
% Fonte: Lezione 6, p.~3

Una \emph{funzione reale di variabile reale} ha dominio e codominio contenuti in $\R$: $f\colon X\to Y$, con $X,Y\subseteq\R$. Quando è assegnata mediante una formula, il suo \emph{dominio naturale} è il più grande sottoinsieme di $\R$ in cui tutte le operazioni della formula hanno significato reale.

\begin{esempio}
Per $f(x)=1/(x^2-4)$ occorre $x^2-4\ne0$. Dunque
\[
\operatorname{dom}(f)=\R\setminus\{-2,2\}.
\]
I punti esclusi dal dominio non tornano a farne parte per una semplice manipolazione algebrica della formula.
\end{esempio}

\section{Operazioni tra funzioni}
% Fonte: Lezione 6, p.~4

Date $f,g\colon D\subseteq\R\to\R$, le operazioni sono definite \emph{punto per punto}:
\[
\begin{aligned}
(f+g)(x)&=f(x)+g(x), & (f-g)(x)&=f(x)-g(x),\\
(fg)(x)&=f(x)g(x), & \left(\frac fg\right)(x)&=\frac{f(x)}{g(x)}.
\end{aligned}
\]
\textbf{Domini delle operazioni.}
\begin{itemize}
  \item Somma, differenza e prodotto hanno dominio $D$.
  \item Il quoziente è definito in $\{x\in D:g(x)\ne0\}$.
  \item Se i domini iniziali sono diversi, si usa la loro intersezione; per il quoziente si escludono inoltre gli zeri del denominatore.
\end{itemize}
Il prodotto $fg$ e la composizione $f\circ g$ sono operazioni diverse.

\section{Monotonia e composizione}
% Fonte: Lezione 6, pp.~5--7

\begin{definizione}
Sia $f\colon D\subseteq\R\to\R$.
\begin{itemize}
  \item È \emph{crescente}, o \emph{non decrescente}, se
\[
\forall x_1,x_2\in D,\quad x_1<x_2\ \Longrightarrow\ f(x_1)\le f(x_2).
\]
È \emph{strettamente crescente} se, sotto la stessa ipotesi, $f(x_1)<f(x_2)$.
  \item È \emph{decrescente}, o \emph{non crescente}, se
\[
x_1<x_2\ \Longrightarrow\ f(x_1)\ge f(x_2),
\]
ed è \emph{strettamente decrescente} se vale $f(x_1)>f(x_2)$.
\end{itemize}
Si dice \emph{monotona} se è crescente oppure decrescente.
\end{definizione}

\textbf{Confronto sull’intero dominio.} La monotonia confronta tutte le coppie ordinate di punti del dominio. Una funzione costante è sia crescente sia decrescente nel senso non stretto; su un dominio con almeno due punti non è strettamente monotona.

\begin{schema}
\begin{tikzpicture}[x=1.5cm,y=.85cm]
  \draw[asse] (-.2,0) -- (4.6,0) node[right] {$x$};
  \draw[asse] (0,-.2) -- (0,3.7) node[above] {$y$};
  \draw[intervallo] (.2,.4) -- (4.1,2.9);
  \draw[thick,orange!80!black] (.2,3.2) -- (4.1,.6);
  \draw[thick,black!55] (.2,3.4) -- (4.1,3.4);
  \node[coloreschema,right] at (4.1,2.9) {crescente};
  \node[orange!80!black,right] at (4.1,.6) {decrescente};
  \node[black!65,right] at (4.1,3.4) {costante};
\end{tikzpicture}
\end{schema}

\begin{proposizione}
\begin{enumerate}[label=\arabic*.]
\item La composizione di due funzioni monotone è monotona: se hanno lo stesso verso è crescente, se hanno versi opposti è decrescente.
\item Se entrambe sono strettamente monotone, anche la composta è strettamente monotona.
\end{enumerate}
\end{proposizione}
\begin{proof}
Scriviamo la composta come $g\circ f$ e scegliamo $x_1<x_2$ nel suo dominio.

\begin{enumerate}[label=\arabic*.]
\item \textbf{Stesso verso.} Se $f$ e $g$ sono crescenti, da $f(x_1)\le f(x_2)$ segue $g(f(x_1))\le g(f(x_2))$. Se entrambe sono decrescenti, si ha $f(x_1)\ge f(x_2)$; applicando $g$ si inverte nuovamente l'ordine e si ottiene ancora $g(f(x_1))\le g(f(x_2))$.

\textbf{Versi opposti.} Se $f$ è crescente e $g$ decrescente, da $f(x_1)\le f(x_2)$ segue $g(f(x_1))\ge g(f(x_2))$. Se $f$ è decrescente e $g$ crescente, da $f(x_1)\ge f(x_2)$ si ottiene la stessa conclusione. La composta è quindi decrescente.
\item \textbf{Caso stretto.} Quando entrambe le funzioni sono strettamente monotone, tutte le disuguaglianze corrispondenti sono strette.
\end{enumerate}
\end{proof}

\section{Stretta monotonia e invertibilità}
% Fonte: Lezione 6, pp.~7--8

\begin{proposizione}
\begin{enumerate}[label=\arabic*.]
\item Ogni funzione strettamente monotona è iniettiva.
\item Se una funzione monotona è iniettiva, è strettamente monotona e la sua inversa sull'immagine ha lo stesso verso di monotonia.
\end{enumerate}
\end{proposizione}
\begin{proof}
\begin{enumerate}[label=\arabic*.]
\item \textbf{Iniettività.} Per $x_1\ne x_2$, uno dei due punti è minore dell'altro. La stretta monotonia impone valori distinti: questo prova l'iniettività.
\item \textbf{Stretta monotonia.} In una funzione monotona iniettiva, l'uguaglianza fra due valori in punti distinti è esclusa, quindi le disuguaglianze diventano strette.

\textbf{Inversa crescente.} Supponiamo $f$ strettamente crescente e siano $y_1<y_2$ nella sua immagine. Posti $x_i=f^{-1}(y_i)$, non può essere $x_1\ge x_2$: si avrebbe $y_1=f(x_1)\ge f(x_2)=y_2$. Dunque $x_1<x_2$ e l'inversa è strettamente crescente.

\textbf{Inversa decrescente.} Se $f$ è strettamente decrescente e $y_1<y_2$, non può essere $x_1\le x_2$: si avrebbe $y_1=f(x_1)\ge f(x_2)=y_2$. Dunque $x_1>x_2$ e l'inversa è strettamente decrescente.
\end{enumerate}
\end{proof}

\begin{osservazione}
L'iniettività da sola non impone la monotonia. Per $f(x)=1/x$ su $\R\setminus\{0\}$:
\begin{itemize}
  \item \textbf{Iniettività.} Da $1/x_1=1/x_2$ segue $x_1=x_2$.
  \item \textbf{Decrescenza sulle semirette.} La funzione è decrescente su ciascuna delle due semirette separatamente; non lo è sull'intero dominio, perché $-1<1$ ma $f(-1)=-1<1=f(1)$.
  \item \textbf{Mancanza di crescenza.} Non è neppure crescente: $1<2$ ma $f(1)>f(2)$.
\end{itemize}
\end{osservazione}

\begin{schema}
\begin{tikzpicture}[scale=.85]
  \draw[asse] (-3.5,0) -- (3.5,0) node[right] {$x$};
  \draw[asse] (0,-3.1) -- (0,3.2) node[above] {$y$};
  \draw[intervallo,domain=.35:3.2,samples=80] plot (\x,{1/\x});
  \draw[intervallo,domain=-3.2:-.35,samples=80] plot (\x,{1/\x});
  \node[coloreschema,right] at (1.2,1.4) {$y=1/x$};
\end{tikzpicture}
\end{schema}

\section{Simmetrie: funzioni pari e dispari}
% Fonte: Lezione 6, p.~9

\textbf{Ipotesi sul dominio.} Per parlare di parità o disparità richiediamo che il dominio sia \emph{simmetrico rispetto a zero}: $x\in D$ implica $-x\in D$.

\begin{definizione}
La funzione $f\colon D\to\R$ è:
\begin{itemize}
  \item \emph{pari} se $f(-x)=f(x)$ per ogni $x\in D$;
  \item \emph{dispari} se $f(-x)=-f(x)$ per ogni $x\in D$.
\end{itemize}
\end{definizione}

\textbf{Interpretazione geometrica.}
\begin{itemize}
  \item Per una funzione pari, da $(x,f(x))$ sul grafico si ottiene $(-x,f(x))$: la simmetria è rispetto all'asse $y$.
  \item Per una funzione dispari si ottiene $(-x,-f(x))$: la simmetria è rispetto all'origine. Se $0\in D$, necessariamente $f(0)=0$, perché $f(0)=-f(0)$.
\end{itemize}
\begin{esempio}
Le funzioni $x^2$ e $x^3$ sono rispettivamente pari e dispari; i grafici sono riportati nel capitolo sulle funzioni elementari.
\end{esempio}

\section{Periodicità}
% Fonte: Lezione 6, p.~10

\begin{definizione}
Una funzione non costante $f\colon D\subseteq\R\to\R$ si dice \emph{periodica}\footnote{Si escludono qui le funzioni costanti. Per una funzione costante su $\R$, ogni $T>0$ sarebbe un periodo, senza un minimo positivo.} se esiste $T>0$ tale che, per ogni $x\in D$,
\[
x+T\in D,\quad x-T\in D,\qquad f(x+T)=f(x).
\]
Ogni tale $T$ è un \emph{periodo}. Se esiste il più piccolo periodo positivo, esso si chiama \emph{periodo principale} o fondamentale.
\end{definizione}

\textbf{Conseguenze della periodicità.}
\begin{itemize}
  \item \textbf{Traslazioni del dominio e dei valori.} Applicando la proprietà più volte, $x+kT\in D$ e $f(x+kT)=f(x)$ per ogni $k\in\Z$. Ogni multiplo intero positivo di un periodo è quindi ancora un periodo.
  \item \textbf{Mancanza di iniettività.} Una funzione periodica non costante non è iniettiva: $x$ e $x+T$ sono punti distinti con la stessa immagine.
\end{itemize}
\begin{esempio}
Seno e coseno hanno periodo principale $2\pi$; la tangente ha periodo principale $\pi$.
\end{esempio}

\section{Limitatezza ed estremi di una funzione}
% Fonte: Lezione 6, pp.~11--13

Per parlare degli estremi consideriamo un dominio $D$ non vuoto.

\begin{definizione}
Una funzione $f\colon D\to\R$ è \emph{limitata} se il suo insieme immagine è limitato; equivalentemente, esistono $m,M\in\R$ tali che
\[
m\le f(x)\le M\qquad\forall x\in D.
\]
Si definiscono gli estremi attraverso l'immagine:
\[
\sup_D f=\sup f(D),\qquad \inf_D f=\inf f(D).
\]
Quando esistono, $\max_D f=\max f(D)$ e $\min_D f=\min f(D)$.
\end{definizione}

\textbf{Valori e punti estremanti.}
\begin{itemize}
  \item $M=\max_D f$ è un \emph{valore massimo}; ogni $x_M\in D$ con $f(x_M)=M$ è un \emph{punto di massimo}.
  \item $m=\min_D f$ è un \emph{valore minimo}; ogni $x_m\in D$ con $f(x_m)=m$ è un \emph{punto di minimo}.
\end{itemize}
Sono estremi \emph{assoluti}: il confronto riguarda tutti i valori sul dominio. Il valore massimo o minimo, se esiste, è unico; può però essere raggiunto in più punti, anche agli estremi del dominio.

\begin{schema}
\begin{tikzpicture}[x=1.65cm,y=.85cm]
  \draw[asse] (-.2,0) -- (6,0) node[right] {$x$};
  \draw[asse] (0,-1.6) -- (0,2.9) node[above] {$y$};
  \draw[dashed,black!45] (0,2.3) node[left] {$M$} -- (5.7,2.3);
  \draw[dashed,black!45] (0,-1.1) node[left] {$m$} -- (5.7,-1.1);
  \draw[intervallo] (.8,2.3) .. controls (1.1,1.2) and (2,1.2) .. (3,2.3)
    .. controls (3.6,2.3) and (3.9,-1.1) .. (4.5,-1.1)
    .. controls (4.9,-1.1) and (5.3,-.5) .. (5.6,0);
  \foreach \x/\y in {.8/2.3,3/2.3,4.5/-1.1,5.6/0}
    \fill[coloreschema] (\x,\y) circle (1.8pt);
  \draw[dotted] (.8,0) -- (.8,2.3);
  \draw[dotted] (3,0) -- (3,2.3);
  \draw[dotted] (4.5,0) -- (4.5,-1.1);
  \node[below] at (.8,0) {$a=x_M^{(1)}$};
  \node[below] at (3,0) {$x_M^{(2)}$};
  \node[above] at (4.5,0) {$x_m$};
  \node[below] at (5.6,0) {$b$};
\end{tikzpicture}

Grafico schematico su $[a,b]$: due punti di massimo e un punto di minimo; ogni valore resta nella fascia $m\le y\le M$.
\end{schema}

\begin{esempio}
La funzione $f(x)=x$ su $[0,1)$ è limitata, ha minimo $0$ nel punto $x=0$ e supremo $1$, ma non ha massimo. La limitatezza non garantisce quindi che entrambi gli estremi siano raggiunti.
\end{esempio}


% !TeX root = ../main.tex
\chapter{Funzioni elementari}
\label{cap:elementari}
\fonti{Lezioni 6 e 7}

% !TeX root = ../main.tex
\section{Potenze a esponente intero positivo}
% Fonte: Lezione 6, pp.~14--15

Sia $n\in\N^*$. La potenza $x^n$ è il prodotto di $n$ fattori uguali a $x$ ed è definita per ogni $x\in\R$. Il comportamento dipende dalla parità dell'esponente.
\begin{itemize}
  \item \textbf{Esponente pari.} Se $n$ è pari, $x^n$ è non negativa su $\R$, nulla soltanto in $0$ e strettamente positiva per $x\ne0$. È una funzione pari, strettamente decrescente su $(-\infty,0]$ e strettamente crescente su $[0,+\infty)$. La sua immagine è $[0,+\infty)$:
  \[
  x\longmapsto x^n\colon\R\longrightarrow[0,+\infty).
  \]
  \item \textbf{Esponente dispari.} Se $n$ è dispari, $x^n$ ha lo stesso segno di $x$: è non negativa su $[0,+\infty)$ e non positiva su $(-\infty,0]$, con segno stretto fuori da $0$. È dispari e strettamente crescente su tutto $\R$. La sua immagine è $\R$:
  \[
  x\longmapsto x^n\colon\R\longrightarrow\R.
  \]
\end{itemize}

\begin{proof}[Verifica del segno e della monotonia]
\textbf{Segno e simmetria.} Seguono dal numero di fattori negativi e dall'identità $(-x)^n=(-1)^n x^n$.

\textbf{Crescita sui non negativi.} Per $0\le u<v$, la fattorizzazione
\[
v^n-u^n=(v-u)\sum_{k=0}^{n-1}v^{n-1-k}u^k>0
\]
prova la stretta crescita sulla semiretta positiva; per $n=1$ la somma vale $1$.

\textbf{Semiretta negativa e confronto fra segni opposti.} Per $u<v\le0$, si ha $0\le-v<-u$. La crescita sui non negativi, insieme alla simmetria, dà $u^n>v^n$ nel caso pari e $u^n<v^n$ nel caso dispari. Nel caso dispari, punti di segno opposto hanno anche immagini di segno opposto: la crescita vale su tutta la retta.

Le immagini dichiarate seguono dall'esistenza delle radici, dimostrata nella prossima sezione.
\end{proof}

\section{Radici e potenze a esponente razionale}
% Fonte: Lezione 6, pp.~16--17

Per ogni $n\in\N^*$ la restrizione della potenza alla semiretta $[0,+\infty)$ è biettiva su $[0,+\infty)$. La sua inversa si chiama \emph{radice $n$-esima}:
\[
\sqrt[n]{\,\cdot\,}\colon[0,+\infty)\longrightarrow[0,+\infty),
\qquad \sqrt[n]{x}=x^{1/n}.
\]
Più precisamente, $\sqrt[n]{x}$ è l'unico numero \emph{non negativo} la cui potenza $n$-esima è $x$; è strettamente positivo quando $x>0$, mentre $\sqrt[n]{0}=0$.

\begin{proposizione}[Esistenza e unicità della radice]
Per $a\ge0$ e $n\in\N^*$ esiste uno e un solo $c\ge0$ tale che $c^n=a$.
\end{proposizione}
\begin{proof}
\textbf{Esistenza: scelta del candidato.} Per $a=0$ basta $c=0$. Per $a>0$ consideriamo
\[
S=\{t\ge0:t^n\le a\}.
\]
È non vuoto perché contiene $0$ ed è limitato superiormente: se $t>\max\{1,a\}$, allora $t^n\ge t>a$. Esiste dunque $c=\sup S\ge0$.

\textbf{Stima della variazione.} Poniamo $B=c+1$ e $K=nB^{n-1}>0$. La fattorizzazione della differenza di potenze dà, per $u,v\in[0,B]$,
\[
|u^n-v^n|\le K|u-v|,
\]
perché ciascuno dei $n$ prodotti della somma è al più $B^{n-1}$.

\textbf{Caso}\enspace\passomath{c^n<a} Scegliamo $0<h<\min\{1,(a-c^n)/K\}$. Allora $(c+h)^n\le c^n+Kh<a$, perciò $c+h\in S$, in contraddizione con il fatto che $c$ è un maggiorante.

\textbf{Caso}\enspace\passomath{c^n>a} In questo caso $c>0$; scegliamo $0<h<\min\{c,(c^n-a)/K\}$. Si ha $(c-h)^n\ge c^n-Kh>a$. Per la stretta crescita della potenza sui non negativi, nessun $t\ge c-h$ appartiene a $S$. Quindi $c-h$ sarebbe un maggiorante di $S$ più piccolo di $c$, ancora impossibile. Resta necessariamente $c^n=a$.

\textbf{Unicità.} La potenza è strettamente crescente sui non negativi: due numeri non negativi distinti non possono avere la stessa potenza $n$-esima. Questo vale anche per $a=0$.
\end{proof}

\textbf{Indice dispari.} Se $n$ è dispari, la potenza è invertibile su tutta la retta e la radice si estende a una funzione $\R\to\R$ ponendo, per $x<0$,
\[
\sqrt[n]{x}=-\sqrt[n]{-x}.
\]
È strettamente crescente e dispari.

\textbf{Indice pari.} Se $n$ è pari, nessun reale ha potenza $n$-esima negativa, quindi il dominio reale della radice è soltanto $[0,+\infty)$. Per $a>0$, l'equazione $t^n=a$ con $n$ pari ha due soluzioni, $\pm\sqrt[n]{a}$: il simbolo di radice sceglie quella non negativa. In particolare
\[
(\sqrt[n]{x})^n=x\quad(x\ge0),\qquad
\sqrt[n]{x^n}=\begin{cases}|x|&n\text{ pari},\\x&n\text{ dispari}.
\end{cases}
\]

\begin{schema}
\begin{tikzpicture}[x=1.35cm,y=.9cm]
  \draw[asse] (-1.7,0) -- (1.8,0) node[right] {$x$};
  \draw[asse] (0,-2.8) -- (0,3) node[above] {$y$};
  \draw[intervallo,domain=-1.55:1.55,samples=90] plot (\x,{\x*\x});
  \draw[thick,orange!80!black,domain=-1.38:1.38,samples=90]
    plot (\x,{\x*\x*\x});
  \node[coloreschema,left,fill=white,inner sep=1pt] at (-1.2,1.9) {$x^2$};
  \node[orange!80!black,right] at (1.1,1.2) {$x^3$};
  \node[below] at (0,-3) {Potenze: esponente pari e dispari};
\end{tikzpicture}
\hspace{.9cm}
\begin{tikzpicture}[x=1.15cm,y=1.25cm]
  \draw[asse] (-2.3,0) -- (2.5,0) node[right] {$x$};
  \draw[asse] (0,-1.8) -- (0,2.15) node[above] {$y$};
  \draw[intervallo,domain=0:1.5,samples=75] plot ({\x*\x},\x);
  \draw[thick,orange!80!black,domain=-1.3:1.3,samples=90]
    plot ({\x*\x*\x},\x);
  \node[coloreschema,above left] at (1.9,1.45) {$\sqrt x$};
  \node[orange!80!black,below right] at (1.4,1.1) {$\sqrt[3]x$};
  \node[below] at (0,-2.16) {Radici: indice pari e dispari};
\end{tikzpicture}
\end{schema}

\textbf{Potenze razionali positive.} Combinando potenze intere e radici si definiscono le potenze a esponente razionale positivo. Per $p,q\in\N^*$ e $x\ge0$ poniamo
\[
x^{p/q}=(\sqrt[q]{x})^p=\sqrt[q]{x^p}.
\]
\textbf{Compatibilità delle due scritture.} Le due espressioni coincidono perché sono non negative e, elevate alla $q$-esima potenza, danno entrambe $x^p$.

\textbf{Indipendenza dalla frazione.} Se due frazioni rappresentano lo stesso razionale, cioè se $p/q=p'/q'$, le due potenze ottenute, elevate alla potenza $qq'$, danno rispettivamente $x^{pq'}$ e $x^{p'q}$, uguali poiché $pq'=p'q$. La stretta crescita della potenza $qq'$ sui non negativi conclude.

Per basi negative occorre verificare separatamente il dominio reale delle radici coinvolte; la definizione appena data riguarda esplicitamente le basi non negative.

% !TeX root = ../main.tex
\section{Valore assoluto e proprietà fondamentali}
% Fonte: Lezione 6, pp.~18--19

\begin{definizione}
Il \emph{valore assoluto}, o \emph{modulo}, di $x\in\R$ è
\[
|x|=\max\{x,-x\}=\begin{cases}
x,&x\ge0,\\
-x,&x<0.
\end{cases}
\]
\end{definizione}

La funzione $x\mapsto|x|$ ha dominio $\R$ e immagine $[0,+\infty)$,
è pari, decresce su $(-\infty,0]$ e cresce su $[0,+\infty)$.

\begin{schema}
\begin{tikzpicture}[x=1cm,y=.8cm]
  \draw[asse] (-3.2,0) -- (3.3,0) node[right] {$x$};
  \draw[asse] (0,-.3) -- (0,3.4) node[above] {$y$};
  \draw[intervallo] (-2.8,2.8) -- (0,0) -- (2.8,2.8);
  \node[coloreschema,above left] at (2.7,2.7) {$y=|x|$};
\end{tikzpicture}
\end{schema}

Per il modulo di una funzione la scelta del ramo dipende dal segno
di \emph{$f(x)$}, non dal segno di $x$:
\[
|f(x)|=\begin{cases}
f(x),&f(x)\ge0,\\
-f(x),&f(x)<0.
\end{cases}
\]
Graficamente si conserva la parte del grafico sopra l'asse $x$ e si
riflette rispetto a tale asse la parte che sta sotto.

\begin{proposizione}
Per ogni $a\in\R$ e ogni $b\ge0$ valgono
\begin{gather*}
a\le|a|,\qquad |a|=0\ \Longleftrightarrow\ a=0,\\
|a|\le b\ \Longleftrightarrow\ -b\le a\le b,\\
|a|\ge b\ \Longleftrightarrow\ a\ge b\ \text{ oppure }\ a\le-b.
\end{gather*}
Per ogni $a,b\in\R$ vale inoltre la \emph{disuguaglianza triangolare}
\[
|a+b|\le|a|+|b|.
\]
\end{proposizione}
\begin{proof}
Poiché $|a|=\max\{a,-a\}$, si ha sia $a\le|a|$ sia $-a\le|a|$;
il massimo è zero esattamente quando $a=0$.
La condizione $|a|\le b$ significa contemporaneamente $a\le b$ e
$-a\le b$, cioè $-b\le a\le b$. La condizione $|a|\ge b$ significa
che almeno uno tra $a$ e $-a$ è maggiore o uguale a $b$, dando la
seconda equivalenza.

Infine $-|a|\le a\le|a|$ e $-|b|\le b\le|b|$; sommando si ottiene
\[
-(|a|+|b|)\le a+b\le|a|+|b|.
\]
La prima equivalenza, con soglia non negativa $|a|+|b|$, conclude
la prova della disuguaglianza triangolare.
\end{proof}

Le versioni strette sono $|a|<b\iff-b<a<b$ e
$|a|>b\iff a>b\text{ oppure }a<-b$, per $b\ge0$.
Nelle disequazioni è essenziale distinguere l'intersezione delle due
condizioni nel primo caso dall'unione nel secondo.

\section{Distanza sulla retta e nozione di metrica}
% Fonte: Lezione 6, pp.~19--20

La distanza usuale fra due numeri reali è
\[
d(x,y)=|x-y|.
\]
È la lunghezza del segmento che li separa e non dipende dall'ordine
in cui si considerano i punti.

\begin{definizione}
Una \emph{distanza}, o \emph{metrica}, su un insieme $E$ (in
particolare su un sottoinsieme di $\R$) è una funzione
$d\colon E\times E\to[0,+\infty)$ tale che, per ogni $x,y,z\in E$,
\begin{enumerate}
  \item $d(x,y)\ge0$;
  \item $d(x,y)=0$ se e solo se $x=y$;
  \item $d(x,y)=d(y,x)$ \emph{(simmetria)};
  \item $d(x,y)\le d(x,z)+d(z,y)$ \emph{(disuguaglianza triangolare)}.
\end{enumerate}
\end{definizione}

\begin{proof}[Verifica per la distanza usuale]
Il modulo è non negativo e si annulla soltanto in zero: quindi
$|x-y|\ge0$ e $|x-y|=0\iff x=y$. Inoltre
$|x-y|=|-(y-x)|=|y-x|$, che dà la simmetria. Infine
\[
|x-y|=|(x-z)+(z-y)|\le|x-z|+|z-y|.
\]
Sono verificate tutte le proprietà della definizione.
\end{proof}

% !TeX root = ../main.tex
\section{Seno, coseno e tangente}
% Fonte: Lezione 6, pp.~10 e 21

Gli angoli sono misurati in \emph{radianti}: un giro completo misura $2\pi$. Sulla circonferenza di centro l'origine e raggio $1$, il punto individuato dall'angolo orientato $x$ ha coordinate
\[
P=(\cos x,\sin x).
\]
Il coseno è dunque l'ascissa e il seno l'ordinata del punto.

\begin{schema}
\begin{tikzpicture}[scale=2.1]
  \draw[asse] (-1.25,0) -- (1.55,0) node[right] {$u$};
  \draw[asse] (0,-1.15) -- (0,1.25) node[above] {$v$};
  \draw[black!55] (0,0) circle (1);
  \draw[orange!80!black,thick] (0,0) -- (1,{tan(35)});
  \draw[coloreschema,thick] (0,0) -- ({cos(35)},{sin(35)});
  \draw[black!45] (1,-.15) -- (1,1.05);
  \draw[dotted] ({cos(35)},0) -- ({cos(35)},{sin(35)})
    -- (0,{sin(35)});
  \draw[->] (.3,0) arc[start angle=0,end angle=35,radius=.3];
  \node at (.41,.12) {$x$};
  \fill[coloreschema] ({cos(35)},{sin(35)}) circle (.8pt)
    node[above left] {$P$};
  \fill[orange!80!black] (1,{tan(35)}) circle (.8pt);
  \node[below] at ({cos(35)},0) {$\cos x$};
  \node[left] at (0,{sin(35)}) {$\sin x$};
  \node[right] at (1.03,.35) {$\tan x$};
  \node[below left] at (0,0) {$O$};
  \node[below right] at (1,-.12) {$1$};
\end{tikzpicture}

Nel disegno $0<x<\pi/2$. La retta di pendenza $\tan x$ incontra la tangente verticale $u=1$ all'altezza $\tan x$.
\end{schema}

\textbf{Dominio, immagine e periodicità.} Seno e coseno sono definiti su tutto $\R$, hanno immagine $[-1,1]$ e periodo principale $2\pi$. La periodicità esprime il ritorno allo stesso punto dopo un giro completo.

\textbf{Parità.} Riflettendo il punto della circonferenza rispetto all'asse orizzontale si ottiene
\[
\sin(-x)=-\sin x,\qquad \cos(-x)=\cos x:
\]
il seno è dispari, il coseno è pari.

\textbf{Identità fondamentale.} Dall'equazione della circonferenza unitaria, o dal teorema di Pitagora, segue l'identità fondamentale
\[
\boxed{\sin^2x+\cos^2x=1}\qquad\forall x\in\R.
\]


\textbf{Estremi.} I valori $\pm1$ sono raggiunti, quindi sono anche minimo e massimo delle due funzioni.

\textbf{Tangente.} La \emph{tangente}, indicata anche con $\operatorname{tg}$, è
\[
\tan x=\frac{\sin x}{\cos x},\qquad
\operatorname{dom}(\tan)=\R\setminus
\left\{\frac\pi2+k\pi:k\in\Z\right\}.
\]
Sono esclusi esattamente gli zeri del denominatore.

\textbf{Parità e periodo.} La tangente è dispari e ha periodo principale $\pi$: seno e coseno cambiano entrambi segno dopo mezzo giro, lasciando invariato il loro rapporto.

\textbf{Immagine e monotonia.} La sua immagine è $\R$, quindi non è limitata. Su ogni intervallo $(-\pi/2+k\pi,\pi/2+k\pi)$ è strettamente crescente e assume tutti i valori reali; non è monotona sull'intero dominio, che contiene più rami distinti.

\section{Restrizioni invertibili e arcoseno}
% Fonte: Lezione 6, p.~22

Seno e coseno sono periodici e perciò non sono iniettivi su $\R$. Per costruirne le inverse si sceglie un intervallo in cui ciascuna funzione sia strettamente monotona e assuma l'intera immagine $[-1,1]$.

\textbf{Ramo invertibile del seno.} Il seno è strettamente crescente da $-1$ a $1$ su $[-\pi/2,\pi/2]$. Sulla semicirconferenza destra, infatti, l'ordinata del punto aumenta al crescere dell'angolo in tale intervallo.

\textbf{Arcoseno.} L'inversa della restrizione è l'\emph{arcoseno}:
\[
\arcsin\colon[-1,1]\longrightarrow
\left[-\frac\pi2,\frac\pi2\right],
\qquad \arcsin y=\text{l'unico }x\in[-\pi/2,\pi/2]\text{ con }\sin x=y.
\]
È strettamente crescente e dispari; il suo grafico è quello del ramo scelto del seno riflesso rispetto a $y=x$.

\begin{schema}
\begin{tikzpicture}[x=.85cm,y=1.15cm]
  \draw[asse] (-3.6,0) -- (3.6,0) node[right] {$x$};
  \draw[asse] (0,-1.4) -- (0,1.5) node[above] {$y$};
  \draw[black!40,domain=-3.4:3.4,samples=100] plot (\x,{sin(\x r)});
  \draw[intervallo,domain=-pi/2:pi/2,samples=70] plot (\x,{sin(\x r)});
  \draw[dotted] (-1.571,0) -- (-1.571,-1) -- (0,-1);
  \draw[dotted] (1.571,0) -- (1.571,1) -- (0,1);
  \fill[coloreschema] (-1.571,-1) circle (1.4pt);
  \fill[coloreschema] (1.571,1) circle (1.4pt);
  \node[above] at (-1.571,0) {$-\frac\pi2$};
  \node[below] at (1.571,0) {$\frac\pi2$};
  \node[above right] at (0,1) {$1$};
  \node[below right] at (0,-1) {$-1$};
  \node[below] at (0,-1.8) {Seno; in colore il ramo invertibile};
\end{tikzpicture}
\hspace{.6cm}
\begin{tikzpicture}[x=1.75cm,y=1cm]
  \draw[asse] (-1.35,0) -- (1.4,0) node[right] {$x$};
  \draw[asse] (0,-1.9) -- (0,1.95) node[above] {$y$};
  \draw[intervallo,domain=-pi/2:pi/2,samples=100] plot ({sin(\x r)},\x);
  \draw[dotted] (-1,0) -- (-1,-1.571) -- (0,-1.571);
  \draw[dotted] (1,0) -- (1,1.571) -- (0,1.571);
  \fill[coloreschema] (-1,-1.571) circle (1.4pt);
  \fill[coloreschema] (1,1.571) circle (1.4pt);
  \node[above] at (-1,0) {$-1$};
  \node[below] at (1,0) {$1$};
  \node[left] at (0,1.571) {$\frac\pi2$};
  \node[right] at (0,-1.571) {$-\frac\pi2$};
  \node[below] at (0,-2.07) {Arcoseno};
\end{tikzpicture}
\end{schema}

\textbf{Identità di inversione e domini.} Si ha
\[
\sin(\arcsin y)=y\quad(y\in[-1,1]),\qquad
\arcsin(\sin x)=x\quad(x\in[-\pi/2,\pi/2]).
\]
\begin{esempio}
La seconda identità non vale per ogni reale: $\arcsin(\sin\pi)=\arcsin0=0\ne\pi$.
\end{esempio}

\section{Arcocoseno}
% Fonte: Lezione 6, p.~23

\textbf{Ramo invertibile del coseno.} Il coseno è strettamente decrescente da $1$ a $-1$ su $[0,\pi]$: percorrendo la semicirconferenza superiore in senso antiorario, l'ascissa diminuisce.

\textbf{Arcocoseno.} Si definisce l'\emph{arcocoseno} come l'inversa di questa restrizione:
\[
\arccos\colon[-1,1]\longrightarrow[0,\pi],\qquad
\arccos y=\text{l'unico }x\in[0,\pi]\text{ con }\cos x=y.
\]
\textbf{Monotonia e identità di inversione.} È strettamente decrescente e soddisfa
\[
\cos(\arccos y)=y\quad(y\in[-1,1]),\qquad
\arccos(\cos x)=x\quad(x\in[0,\pi]).
\]
\textbf{Valori notevoli.} Si ha $\arccos1=0$, $\arccos0=\pi/2$ e $\arccos(-1)=\pi$.

\begin{schema}
\begin{tikzpicture}[x=.85cm,y=1.15cm]
  \draw[asse] (-1.9,0) -- (4.5,0) node[right] {$x$};
  \draw[asse] (0,-1.4) -- (0,1.5) node[above] {$y$};
  \draw[black!40,domain=-1.7:4.3,samples=100] plot (\x,{cos(\x r)});
  \draw[intervallo,domain=0:pi,samples=70] plot (\x,{cos(\x r)});
  \draw[dotted] (3.142,0) -- (3.142,-1) -- (0,-1);
  \fill[coloreschema] (0,1) circle (1.4pt);
  \fill[coloreschema] (3.142,-1) circle (1.4pt);
  \node[above] at (3.142,0) {$\pi$};
  \node[left] at (0,1) {$1$};
  \node[left] at (0,-1) {$-1$};
  \node[below] at (1.2,-1.8) {Coseno; in colore il ramo invertibile};
\end{tikzpicture}
\hspace{.6cm}
\begin{tikzpicture}[x=1.75cm,y=1cm]
  \draw[asse] (-1.35,0) -- (1.4,0) node[right] {$x$};
  \draw[asse] (0,-.3) -- (0,3.55) node[above] {$y$};
  \draw[intervallo,domain=0:pi,samples=100] plot ({cos(\x r)},\x);
  \draw[dotted] (-1,0) -- (-1,3.142) -- (0,3.142);
  \fill[coloreschema] (-1,3.142) circle (1.4pt);
  \fill[coloreschema] (1,0) circle (1.4pt);
  \node[below] at (-1,0) {$-1$};
  \node[below] at (1,0) {$1$};
  \node[right] at (0,3.142) {$\pi$};
  \node[left] at (0,1.571) {$\frac\pi2$};
  \node[below] at (0,-.92) {Arcocoseno};
\end{tikzpicture}
\end{schema}

\Needspace{31\baselineskip}
\section{Arcotangente}
% Fonte: Lezione 6, p.~24

\textbf{Ramo invertibile della tangente.} La restrizione $\tan\colon(-\pi/2,\pi/2)\to\R$ è biettiva. Geometricamente, mentre il punto percorre la semicirconferenza destra, la retta che passa per l'origine incontra la verticale $u=1$ a un'altezza che cresce da valori negativi arbitrariamente grandi in modulo a valori positivi arbitrariamente grandi. Ogni altezza reale è raggiunta una sola volta.

\textbf{Arcotangente.} Si definisce l'\emph{arcotangente}:
\[
\arctan\colon\R\longrightarrow\left(-\frac\pi2,\frac\pi2\right),
\qquad \arctan y=\text{l'unico }x\in(-\pi/2,\pi/2)\text{ con }\tan x=y.
\]
\textbf{Proprietà e identità di inversione.} È strettamente crescente, dispari e limitata. Vale
\[
\tan(\arctan y)=y\quad(y\in\R),\qquad
\arctan(\tan x)=x\quad(x\in(-\pi/2,\pi/2)).
\]
\textbf{Estremi non raggiunti.} Il suo supremo è $\pi/2$ e il suo infimo è $-\pi/2$, ma non ha né massimo né minimo: gli estremi dell'intervallo immagine sono esclusi. Le rette $y=\pm\pi/2$ indicano le altezze a cui il grafico si avvicina per argomenti rispettivamente molto grandi positivi e negativi.

\begin{schema}
\begin{tikzpicture}[x=.85cm,y=.65cm]
  \draw[asse] (-2.3,0) -- (4.9,0) node[right] {$x$};
  \draw[asse] (0,-3) -- (0,3.2) node[above] {$y$};
  \foreach \a in {-1.571,1.571,4.712}
    \draw[dashed,black!40] (\a,-2.8) -- (\a,2.8);
  \draw[intervallo,domain=-1.22:1.22,samples=90] plot (\x,{tan(\x r)});
  \draw[black!45,domain=1.93:4.36,samples=90] plot (\x,{tan(\x r)});
  \node[below left] at (-1.571,0) {$-\frac\pi2$};
  \node[below right] at (1.571,0) {$\frac\pi2$};
  \node[below] at (4.712,0) {$\frac{3\pi}2$};
  \node[below] at (1.2,-3.5) {Tangente; in colore il ramo invertibile};
\end{tikzpicture}
\hspace{.6cm}
\begin{tikzpicture}[x=.65cm,y=1.1cm]
  \draw[asse] (-3.4,0) -- (3.5,0) node[right] {$x$};
  \draw[asse] (0,-1.95) -- (0,1.95) node[above] {$y$};
  \draw[dashed,black!45] (-3.2,1.571) -- (3.2,1.571);
  \draw[dashed,black!45] (-3.2,-1.571) -- (3.2,-1.571);
  \draw[intervallo,domain=-3.2:3.2,samples=100] plot (\x,{atan(\x)*pi/180});
  \node[above right] at (0,1.571) {$\frac\pi2$};
  \node[below right] at (0,-1.571) {$-\frac\pi2$};
  \node[below] at (0,-2.17) {Arcotangente};
\end{tikzpicture}
\end{schema}

% !TeX root = ../main.tex
% Fonte: Lezione 7.pdf, pp. 2--9. La costruzione delle potenze reali
% è rinviata nella fonte; la nota ne dà una possibile definizione.
\section{Esponenziali e logaritmi}
\label{sec:exp-log}

\subsection{La funzione esponenziale}

\textbf{Esponenti nulli e negativi.} Per una base $a>0$, le potenze razionali si estendono all'esponente zero ponendo $a^0=1$ e agli esponenti negativi ponendo $a^{-r}=1/a^r$ per $r>0$ razionale.

\textbf{Estensione agli esponenti reali.} La funzione $x\mapsto a^x$ estende a tutti gli esponenti reali le potenze già definite per esponenti razionali.\footnote{Per $a>1$ si può porre $a^x=\sup\{a^r:r\in\Q,\ r\le x\}$; per $0<a<1$ si pone $a^x=(1/a)^{-x}$, mentre $1^x=1$. La verifica completa di questa costruzione e delle proprietà dell'esponenziale reale non è svolta qui: le regole delle potenze, la stretta monotonia e l'immagine $(0,+\infty)$, per $a\ne1$, sono assunte come proprietà note.} Ha dominio $\R$, assume soltanto valori strettamente positivi e passa per $(0,1)$, perché $a^0=1$.
\begin{itemize}
  \item Se $a>1$, è strettamente crescente e biettiva da $\R$ a $(0,+\infty)$.
  \item Se $0<a<1$, è strettamente decrescente e ancora biettiva da $\R$ a $(0,+\infty)$.
  \item Se $a=1$, è la funzione costante $1$ e non è invertibile su $\R$.
\end{itemize}
\textbf{Comportamento ed estremi.} Al crescere dell'argomento, nel caso $a>1$ i valori diventano arbitrariamente grandi; nel caso $0<a<1$ si avvicinano a zero senza mai raggiungerlo. Per argomenti molto negativi i due comportamenti si scambiano. Perciò, quando $a\ne1$, l'infimo è $0$ senza minimo, mentre non esiste un maggiorante reale.

\textbf{Regole delle potenze.} Si ha
\[
a^{x+y}=a^xa^y,\qquad a^{-x}=\frac1{a^x},\qquad a^0=1.
\]
% Raccordo con Lezione 15, pp. 13--15: definizione rigorosa e rappresentazioni nel capitolo sui limiti notevoli.
\textbf{Base $e$.} Una base particolarmente importante è il \emph{numero di Nepero} $e=2{,}718281828\ldots$, definito nella sezione~\ref{sec:numero-nepero}. Con la scrittura \emph{funzione esponenziale}, senza precisare la base, si indica di solito
\[
\exp\colon\R\longrightarrow(0,+\infty),\qquad \exp(x)=e^x,
\]
che è strettamente crescente e biettiva.

\begin{schema}
\begin{tikzpicture}[x=1cm,y=.65cm]
  \draw[asse] (-2.5,0) -- (2.6,0) node[right] {$x$};
  \draw[asse] (0,-.3) -- (0,5.2) node[above] {$y$};
  \draw[intervallo,domain=-2.2:2.2,samples=70] plot (\x,{pow(2,\x)});
  \draw[thick,orange!80!black,domain=-2.2:2.2,samples=70]
    plot (\x,{pow(2,-\x)});
  \fill[coloreschema] (0,1) circle (1.6pt) node[above right] {$1$};
  \node[coloreschema,right] at (1.6,3.1) {$2^x$};
  \node[orange!80!black,left] at (-1.6,3.1) {$(1/2)^x$};
\end{tikzpicture}
\hspace{.6cm}
\begin{tikzpicture}[x=1cm,y=.85cm]
  \draw[asse] (-.3,0) -- (4.8,0) node[right] {$x$};
  \draw[asse] (0,-2) -- (0,2.15) node[above] {$y$};
  \draw[intervallo,domain=.17:4.5,samples=80] plot (\x,{ln(\x)/ln(2)});
  \draw[thick,orange!80!black,domain=.17:4.5,samples=80]
    plot (\x,{-ln(\x)/ln(2)});
  \fill[coloreschema] (1,0) circle (1.6pt) node[below left] {$1$};
  \node[coloreschema,above left] at (4.4,2.1) {$\log_2x$};
  \node[orange!80!black,below left] at (4.4,-1.9) {$\log_{1/2}x$};
\end{tikzpicture}

Esponenziali con basi maggiore e minore di $1$, e rispettivi logaritmi.
\end{schema}

\subsection{Il logaritmo come funzione inversa}

Per $a>0$, $a\ne1$, si definisce il \emph{logaritmo in base $a$} come l'inversa dell'esponenziale:
\[
\log_a\colon(0,+\infty)\longrightarrow\R,\qquad
\log_a x=y\ \Longleftrightarrow\ a^y=x.
\]
\textbf{Dominio, immagine e monotonia.} Il dominio contiene soltanto numeri positivi; l'immagine è tutta $\R$. Il logaritmo è strettamente crescente se $a>1$, strettamente decrescente se $0<a<1$, perché l'inversa conserva il verso della monotonia.

\textbf{Identità del logaritmo.} Per ogni base ammissibile,
\[
\log_a1=0,\qquad \log_a a=1,\qquad
\log_a(a^y)=y\ (y\in\R),\qquad a^{\log_a x}=x\ (x>0).
\]
\textbf{Logaritmo naturale.} Se la base è $e$, si parla di \emph{logaritmo naturale} e si scrive $\ln x=\log_e x$. In particolare $\ln\colon(0,+\infty)\to\R$ è strettamente crescente e inversa di $\exp$.

\section{Funzioni iperboliche}
\label{sec:iperboliche}

\subsection{Definizioni, identità e grafici}
% Fonte: Lezione 7.pdf, pp. 5--7.

Si definiscono \emph{seno iperbolico}, \emph{coseno iperbolico} e \emph{tangente iperbolica} mediante
\[
\sinh x=\frac{e^x-e^{-x}}2,\qquad
\cosh x=\frac{e^x+e^{-x}}2,\qquad
\tanh x=\frac{\sinh x}{\cosh x}.
\]
Sono tutte definite su $\R$: infatti $\cosh x>0$ e il quoziente non ha denominatore nullo. La tangente iperbolica si indica anche con $\operatorname{tgh}$.

\Needspace{6\baselineskip}
\textbf{Proprietà delle tre funzioni.}
\begin{itemize}
  \item \textbf{Seno iperbolico:} dispari, strettamente crescente, con immagine $\R$.
  \item \textbf{Coseno iperbolico:} pari, strettamente decrescente su $(-\infty,0]$ e strettamente crescente su $[0,+\infty)$; ha minimo $1$ nel punto $0$ e immagine $[1,+\infty)$.
  \item \textbf{Tangente iperbolica:} dispari, strettamente crescente, con immagine $(-1,1)$; è limitata, senza massimo né minimo.
\end{itemize}

\begin{proposizione}[Identità iperbolica fondamentale]
Per ogni $t\in\R$ si ha
\[
\boxed{\cosh^2t-\sinh^2t=1}.
\]
\end{proposizione}
\begin{proof}
Sostituendo le definizioni e usando $e^te^{-t}=1$ si ottiene
\[
\cosh^2t-\sinh^2t
=\frac{(e^t+e^{-t})^2-(e^t-e^{-t})^2}{4}
=\frac{4e^te^{-t}}4=1.
\]
\end{proof}

\textbf{Interpretazione geometrica.} Il nome richiama una costruzione geometrica: mentre $P(t)=(\cos t,\sin t)$ percorre la circonferenza $u^2+v^2=1$, il punto $Q(t)=(\cosh t,\sinh t)$ percorre il \emph{ramo destro} dell'iperbole $u^2-v^2=1$, quello con $u\ge1$.

\begin{schema}
\begin{tikzpicture}[scale=1.45]
  \draw[asse] (-1.3,0) -- (1.45,0) node[right] {$u$};
  \draw[asse] (0,-1.25) -- (0,1.4) node[above] {$v$};
  \draw[intervallo] (0,0) circle (1);
  \draw[dotted] (0,0) -- ({cos(40)},{sin(40)});
  \fill[coloreschema] ({cos(40)},{sin(40)}) circle (1pt)
    node[above right] {$P(t)$};
  \node[below] at (0,-1.5) {$u^2+v^2=1$};
\end{tikzpicture}
\hspace{1.3cm}
\begin{tikzpicture}[scale=1.45]
  \draw[asse] (-.3,0) -- (2.25,0) node[right] {$u$};
  \draw[asse] (0,-1.8) -- (0,1.8) node[above] {$v$};
  \draw[intervallo,domain=-1.2:1.2,samples=90]
    plot ({(exp(\x)+exp(-\x))/2},{(exp(\x)-exp(-\x))/2});
  \fill[coloreschema] (1,0) circle (1pt) node[below left] {$1$};
  \fill[coloreschema] ({(exp(.8)+exp(-.8))/2},{(exp(.8)-exp(-.8))/2})
    circle (1pt) node[above left] {$Q(t)$};
  \node[below] at (1,-1.99) {$u^2-v^2=1,\quad u\ge1$};
\end{tikzpicture}
\end{schema}

\begin{proof}[Verifica delle proprietà]
\textbf{Simmetrie.} La parità di $\cosh$ e la disparità di $\sinh$ seguono scambiando $t$ con $-t$.

\textbf{Crescita del seno iperbolico.} Se $u<v$, si ha $e^u<e^v$ ed $e^{-u}>e^{-v}$: dalla differenza delle due quantità segue $\sinh u<\sinh v$.

\textbf{Minimo del coseno iperbolico.} Si ha
\[
\cosh x-1=\frac{(e^{x/2}-e^{-x/2})^2}{2}\ge0,
\]
con uguaglianza soltanto in $x=0$.

\textbf{Monotonia del coseno iperbolico.} Per $0\le u<v$, sviluppando le definizioni si verifica
\[
\cosh v-\cosh u
=2\sinh\frac{u+v}2\,\sinh\frac{v-u}2>0.
\]
La parità dà poi la decrescenza sulla semiretta negativa.

\textbf{Tangente iperbolica.} È utile scrivere $\tanh x=(e^{2x}-1)/(e^{2x}+1)$. Per $0<s<t$,
\[
\frac{t-1}{t+1}-\frac{s-1}{s+1}
=\frac{2(t-s)}{(t+1)(s+1)}>0;
\]
quindi $\tanh$ è strettamente crescente. Il rapporto è sempre compreso fra $-1$ e $1$ e, sostituendo $-x$, cambia segno. Le immagini esatte delle tre funzioni seguono anche dalle formule inverse ricavate qui sotto.
\end{proof}

\begin{schema}
\begin{tikzpicture}[x=.8cm,y=.65cm]
  \draw[asse] (-2.3,0) -- (2.4,0) node[right] {$x$};
  \draw[asse] (0,-3.4) -- (0,3.6) node[above] {$y$};
  \draw[intervallo,domain=-1.9:1.9,samples=80]
    plot (\x,{(exp(\x)-exp(-\x))/2});
  \node[below] at (0,-3.7) {$\sinh x$};
\end{tikzpicture}
\hspace{.3cm}
\begin{tikzpicture}[x=.8cm,y=.65cm]
  \draw[asse] (-2.3,0) -- (2.4,0) node[right] {$x$};
  \draw[asse] (0,-.4) -- (0,3.8) node[above] {$y$};
  \draw[dashed,black!40] (-2.1,1) -- (2.1,1);
  \draw[intervallo,domain=-1.9:1.9,samples=80]
    plot (\x,{(exp(\x)+exp(-\x))/2});
  \node[below right] at (0,1) {$1$};
  \node[below] at (0,-.7) {$\cosh x$};
\end{tikzpicture}
\hspace{.3cm}
\begin{tikzpicture}[x=.8cm,y=1.15cm]
  \draw[asse] (-2.3,0) -- (2.4,0) node[right] {$x$};
  \draw[asse] (0,-1.5) -- (0,1.5) node[above] {$y$};
  \foreach \y in {-1,1} \draw[dashed,black!40] (-2.2,\y) -- (2.2,\y);
  \draw[intervallo,domain=-2.2:2.2,samples=80]
    plot (\x,{(exp(2*\x)-1)/(exp(2*\x)+1)});
  \node[above right] at (0,1) {$1$};
  \node[below right] at (0,-1) {$-1$};
  \node[below] at (0,-1.75) {$\tanh x$};
\end{tikzpicture}
\end{schema}

\subsection{Funzioni iperboliche inverse}
% Fonte: Lezione 7.pdf, pp. 8--9.

\textbf{Settore seno iperbolico.} Il seno iperbolico è biettivo da $\R$ a $\R$. La sua inversa è il \emph{settore seno iperbolico}, indicato con $\operatorname{arsinh}$ o $\operatorname{settsinh}$, e si esprime mediante il logaritmo:
\[
\operatorname{arsinh}x=\ln\!\left(x+\sqrt{x^2+1}\right),
\qquad \operatorname{arsinh}\colon\R\longrightarrow\R.
\]
\begin{proof}[Derivazione]
Risolviamo $x=\sinh y$ e poniamo $t=e^y>0$:
\begin{align*}
x=\frac{e^y-e^{-y}}2
&\ \Longleftrightarrow\ 2x=t-\frac1t\\
&\ \Longleftrightarrow\ t^2-2xt-1=0\\
&\ \Longleftrightarrow\ t=x\pm\sqrt{x^2+1}.
\end{align*}
Poiché $\sqrt{x^2+1}>|x|$, il segno meno darebbe $t<0$ e va scartato; il segno più dà un valore strettamente positivo per ogni reale $x$. Dunque $e^y=x+\sqrt{x^2+1}$ e si applica $\ln$.
\end{proof}

\textbf{Settore coseno iperbolico.} Il coseno iperbolico è invertibile dopo averne ristretto il dominio a $[0,+\infty)$. L'inversa è il \emph{settore coseno iperbolico}:
\[
\operatorname{arcosh}x=\ln\!\left(x+\sqrt{x^2-1}\right),
\qquad \operatorname{arcosh}\colon[1,+\infty)\longrightarrow[0,+\infty).
\]
La notazione alternativa è $\operatorname{settcosh}$.

\begin{proof}[Scelta del ramo]
Da $x=\cosh y$ e $t=e^y$ segue $t^2-2xt+1=0$, con soluzioni $t=x\pm\sqrt{x^2-1}$. Per $x>1$ il segno più dà $t>1$, dunque $y>0$; il segno meno dà $0<t<1$, dunque $y<0$, escluso dalla restrizione scelta. Per $x=1$ le due soluzioni coincidono in $t=1$, cioè $y=0$.
\end{proof}

\textbf{Settore tangente iperbolica.} La tangente iperbolica è biettiva da $\R$ a $(-1,1)$. La sua inversa, il \emph{settore tangente iperbolica}, è
\[
\operatorname{artanh}x=\frac12\ln\!\left(\frac{1+x}{1-x}\right),
\qquad \operatorname{artanh}\colon(-1,1)\longrightarrow\R,
\]
indicata anche con $\operatorname{setttgh}$.

\begin{proof}[Derivazione]
L'equazione $x=(e^{2y}-1)/(e^{2y}+1)$ equivale a $e^{2y}=(1+x)/(1-x)$. Per $-1<x<1$ il rapporto è positivo, quindi si può applicare il logaritmo e dividere per $2$.
\end{proof}


% !TeX root = ../main.tex
% Fonte: Lezione 7.pdf, pp. 10--18.
\section{Funzioni definite a tratti}

\subsection{Parte positiva e parte negativa}

Per ogni reale $x$ si definiscono la \emph{parte positiva} e la \emph{parte negativa} mediante
\[
x^+=\max\{x,0\}=\begin{cases}x,&x\ge0,\\0,&x<0,\end{cases}
\qquad
x^-=-\min\{x,0\}=\begin{cases}-x,&x\le0,\\0,&x>0.\end{cases}
\]
\textbf{Segno e identità.} Entrambe sono non negative: la parte negativa misura il modulo della componente negativa e non è un numero negativo. Distinguendo i casi $x\ge0$ e $x<0$ si verificano le identità
\[
\boxed{x=x^+-x^-,\qquad |x|=x^++x^-}.
\]


\textbf{Parti di una funzione.} Le stesse definizioni si applicano punto per punto a una funzione: $f^+(x)=\max\{f(x),0\}$ e $f^-(x)=\max\{-f(x),0\}$, con $f=f^+-f^-$ e $|f|=f^++f^-$.

\begin{schema}
\begin{tikzpicture}[x=1cm,y=.85cm]
  \draw[asse] (-2.4,0) -- (2.5,0) node[right] {$x$};
  \draw[asse] (0,-.3) -- (0,2.6) node[above] {$y$};
  \draw[intervallo] (-2.1,0) -- (0,0) -- (2.1,2.1);
  \node[above left] at (2.1,2.1) {$x^+$};
\end{tikzpicture}
\hspace{1cm}
\begin{tikzpicture}[x=1cm,y=.85cm]
  \draw[asse] (-2.4,0) -- (2.5,0) node[right] {$x$};
  \draw[asse] (0,-.3) -- (0,2.6) node[above] {$y$};
  \draw[intervallo] (-2.1,2.1) -- (0,0) -- (2.1,0);
  \node[above right] at (-2.1,2.1) {$x^-$};
\end{tikzpicture}
\end{schema}

\subsection{Parte intera e parte frazionaria}
\enlargethispage{2pt}

\begin{definizione}
La \emph{parte intera} di $x\in\R$, o \emph{floor}, è l'unico intero $m$ tale che $m\le x<m+1$. Si indica con $\lfloor x\rfloor$.\footnote{Si usano anche $[x]$ per la parte intera e $(x)$ per la parte frazionaria; qui preferiamo $\lfloor x\rfloor$ e $\{x\}$ per distinguere chiaramente le due funzioni dalle altre parentesi.} La \emph{parte frazionaria} è $\{x\}=x-\lfloor x\rfloor$.
\end{definizione}

\begin{proof}[Esistenza e unicità]
\textbf{Esistenza.} Archimede e il principio del minimo intero forniscono, come nella prova della densità dei razionali, un intero $m$ con $x-1<m\le x$: basta applicare a $x-1$ il fatto che per ogni reale $t$ esiste un intero in $(t,t+1]$. Si ottiene $m\le x<m+1$.

\textbf{Unicità.} Se due interi $m<k$ soddisfacessero entrambe le condizioni, si avrebbe $m+1\le k\le x<m+1$, assurdo: due interi distinti distano almeno $1$.
\end{proof}

Per ogni reale $x$ si ha $0\le\{x\}<1$; dunque
\[
\lfloor\,\cdot\,\rfloor\colon\R\to\Z,\qquad
\{\,\cdot\,\}\colon\R\to[0,1).
\]
\textbf{Grafici.} Per $m\le x<m+1$, il grafico della parte intera è il segmento orizzontale $y=m$ e quello della parte frazionaria è il segmento $y=x-m$, di pendenza $1$. In entrambi, l'estremo sinistro è incluso e il destro escluso.

\begin{proof}[Periodo principale]
\textbf{Periodo $1$ della parte frazionaria.} Se $m=\lfloor x\rfloor$, allora $m+1\le x+1<m+2$: dunque $\lfloor x+1\rfloor=\lfloor x\rfloor+1$ e $\{x+1\}=\{x\}$.

\textbf{Minimalità del periodo.} Ogni periodo $T>0$ soddisfa $\{T\}=\{0\}=0$, quindi $T\in\N^*$ e $T\ge1$. Il periodo $1$ è pertanto il periodo principale.
\end{proof}

\begin{esempio}
Si ha $\lfloor2{,}7\rfloor=2$ e $\{2{,}7\}=0{,}7$, ma $\lfloor-1{,}3\rfloor=-2$ e $\{-1{,}3\}=0{,}7$. La parte intera arrotonda verso il basso, anche per i negativi; non coincide con l'eliminazione delle cifre dopo la virgola.
\end{esempio}

\begin{schema}
\begin{tikzpicture}[x=.8cm,y=.65cm]
  \draw[asse] (-2.3,0) -- (3.4,0) node[right] {$x$};
  \draw[asse] (0,-2.4) -- (0,3.4) node[above] {$y$};
  \foreach \m in {-2,-1,0,1,2} {
    \draw[intervallo] (\m,\m) -- ({\m+1},\m);
    \fill[coloreschema] (\m,\m) circle (1.7pt);
    \draw[estremo] ({\m+1},\m) circle (1.7pt);
  }
  \draw[intervallo] (3,3) -- (3.25,3);
  \fill[coloreschema] (3,3) circle (1.7pt);
  \foreach \m in {-2,-1,1,2,3} \node[below=4pt] at (\m,0) {$\m$};
  \node[below] at (.5,-2.9) {$y=\lfloor x\rfloor$};
\end{tikzpicture}
\hspace{.6cm}
\begin{tikzpicture}[x=.8cm,y=1.3cm]
  \draw[asse] (-2.3,0) -- (3.4,0) node[right] {$x$};
  \draw[asse] (0,-.4) -- (0,1.6) node[above] {$y$};
  \draw[dotted,black!40] (-2.1,1) -- (3.1,1);
  \foreach \m in {-2,-1,0,1,2} {
    \draw[intervallo] (\m,0) -- ({\m+1},1);
    \fill[coloreschema] (\m,0) circle (1.7pt);
    \draw[estremo] ({\m+1},1) circle (1.7pt);
  }
  \draw[intervallo] (3,0) -- (3.25,.25);
  \fill[coloreschema] (3,0) circle (1.7pt);
  \foreach \m in {-2,-1,1,2,3} \node[below=4pt] at (\m,0) {$\m$};
  \node[above left] at (0,1) {$1$};
  \node[below] at (.5,-.8) {$y=\{x\}$};
\end{tikzpicture}
\end{schema}

\section{Trasformazioni dei grafici}
% Fonte: Lezione 7.pdf, p. 13.

Sia $D$ il dominio di $f$. Dal grafico di $y=f(x)$ si possono ottenere altri grafici seguendo come cambia un suo punto $(x,f(x))$:
\begin{itemize}
  \item \textbf{Traslazione verticale.} $y=f(x)+a$: ogni punto diventa $(x,f(x)+a)$; il grafico si trasla verticalmente di $a$, verso l'alto se $a>0$ e verso il basso se $a<0$;
  \item \textbf{Traslazione orizzontale.} $y=f(x-b)$: ogni punto $(u,f(u))$ diventa $(u+b,f(u))$; la traslazione è orizzontale, verso destra se $b>0$ e verso sinistra se $b<0$. Il dominio diventa $D+b=\{u+b:u\in D\}$;
  \item \textbf{Riflessione rispetto all’asse $x$.} $y=-f(x)$: ogni punto diventa $(x,-f(x))$; il grafico viene riflesso rispetto all'asse $x$.
\end{itemize}
Per $b>0$, il segno meno \emph{dentro} $f(x-b)$ corrisponde a uno spostamento verso destra: per conservare l'argomento $u$ si prende $x=u+b$. Il verso della traslazione dipende quindi dal segno di $b$.

\begin{schema}
\begin{tikzpicture}[x=.8cm,y=.55cm]
  \draw[asse] (-.7,0) -- (3.5,0) node[right] {$x$};
  \draw[asse] (0,-2.3) -- (0,4.2) node[above] {$y$};
  \draw[black!45,domain=-.4:3.05,samples=90] plot (\x,{\x*(\x-2)^2});
  \draw[intervallo,domain=-.4:3.05,samples=90] plot (\x,{\x*(\x-2)^2+1});
  \node[below] at (1.3,-2.9) {$f(x)+1$};
\end{tikzpicture}
\hspace{.35cm}
\begin{tikzpicture}[x=.8cm,y=.55cm]
  \draw[asse] (-.7,0) -- (4.5,0) node[right] {$x$};
  \draw[asse] (0,-2.3) -- (0,4.2) node[above] {$y$};
  \draw[black!45,domain=-.4:3.05,samples=90] plot (\x,{\x*(\x-2)^2});
  \draw[intervallo,domain=.6:4.05,samples=90] plot (\x,{(\x-1)*(\x-3)^2});
  \node[below] at (1.8,-2.9) {$f(x-1)$};
\end{tikzpicture}
\hspace{.35cm}
\begin{tikzpicture}[x=.8cm,y=.55cm]
  \draw[asse] (-.7,0) -- (3.5,0) node[right] {$x$};
  \draw[asse] (0,-3.6) -- (0,3.5) node[above] {$y$};
  \draw[black!45,domain=-.4:3.05,samples=90] plot (\x,{\x*(\x-2)^2});
  \draw[intervallo,domain=-.4:3.05,samples=90] plot (\x,{-\x*(\x-2)^2});
  \node[below] at (1.3,-4.1) {$-f(x)$};
\end{tikzpicture}

In grigio il grafico iniziale; in colore la trasformazione indicata.
\end{schema}

\Needspace{16\baselineskip}
\section{Un esercizio su dominio, monotonia e inversa}
% Fonte: Lezione 7.pdf, pp. 14--18. Corretto il titolo «non decr.»
% della p. 16, incompatibile con il calcolo decrescente svolto sotto.

\begin{esercizio}
Studiamo la funzione
\[
f(x)=4\arcsin\!\bigl(1-\ln(x-1)\bigr).
\]
Cerchiamo il dominio naturale, il verso di monotonia e l'insieme dei valori $y\in\R$ per i quali $f(x)=y$ ammette una soluzione.
\end{esercizio}

\subsection{Determinazione del dominio}

\begin{svolgimento}
\textbf{Condizioni di definizione.} Il logaritmo richiede $x-1>0$; l'arcoseno richiede che il proprio argomento sia in $[-1,1]$. Quindi
\begin{align*}
x\in D
&\ \Longleftrightarrow\
\begin{cases}x>1,\\-1\le1-\ln(x-1)\le1\end{cases}\\
&\ \Longleftrightarrow\
\begin{cases}x>1,\\0\le\ln(x-1)\le2\end{cases}\\
&\ \Longleftrightarrow\
\begin{cases}x>1,\\1\le x-1\le e^2.\end{cases}
\end{align*}
L'ultima doppia disuguaglianza equivale a $2\le x\le1+e^2$ e implica già $x>1$. Pertanto
\[
\boxed{D=[2,1+e^2]}.
\]
\end{svolgimento}

\subsection{Monotonia}

\begin{svolgimento}
Siano $x_1,x_2\in D$ con $x_1<x_2$. Il logaritmo è strettamente crescente, perciò
\[
x_1-1<x_2-1
\ \Longrightarrow\ \ln(x_1-1)<\ln(x_2-1)
\ \Longrightarrow\ 1-\ln(x_1-1)>1-\ln(x_2-1).
\]
L'arcoseno è strettamente crescente e moltiplicare per $4>0$ conserva il verso: ne segue $f(x_1)>f(x_2)$. Dunque $f$ è \emph{strettamente decrescente}, quindi iniettiva sul proprio dominio.

\end{svolgimento}

\Needspace{30\baselineskip}

\subsection{Immagine, equazione e inversa}

\begin{svolgimento}
\textbf{Necessità.} Poiché l'arcoseno assume valori in $[-\pi/2,\pi/2]$, una condizione necessaria per $f(x)=y$ è $-2\pi\le y\le2\pi$.

\textbf{Sufficienza: costruzione della soluzione.} Sotto questa condizione $y/4$ appartiene proprio all'intervallo in cui seno e arcoseno sono inversi, e possiamo procedere mediante equivalenze:
\begin{align*}
4\arcsin\!\bigl(1-\ln(x-1)\bigr)=y
&\ \Longleftrightarrow\ 1-\ln(x-1)=\sin(y/4)\\
&\ \Longleftrightarrow\ \ln(x-1)=1-\sin(y/4)\\
&\ \Longleftrightarrow\ x=1+e^{\,1-\sin(y/4)}.
\end{align*}
Per ogni $y\in[-2\pi,2\pi]$ il valore $1-\sin(y/4)$ è in $[0,2]$; quindi la soluzione trovata appartiene a $[2,1+e^2]$ ed è valida. Per i valori di $y$ esterni all'intervallo non ci sono soluzioni. In conclusione,
\[
\operatorname{Im}(f)=[-2\pi,2\pi],\qquad
f^{-1}(y)=1+e^{\,1-\sin(y/4)}\quad(-2\pi\le y\le2\pi).
\]
\textbf{Unicità.} La soluzione è unica per l'iniettività.
\end{svolgimento}

\textbf{Massimo e minimo.} Gli estremi sono raggiunti: $f(2)=2\pi$ è il massimo e $f(1+e^2)=-2\pi$ è il minimo.

\begin{schema}
\begin{tikzpicture}[x=.85cm,y=.4cm]
  \draw[asse] (-.4,0) -- (9.2,0) node[right] {$x$};
  \draw[asse] (0,-7.5) -- (0,7.6) node[above] {$y$};
  \draw[intervallo,domain=-pi/2:pi/2,samples=100]
    plot ({1+exp(1-sin(\x r))},{4*\x});
  \draw[dotted] (2,0) -- (2,{2*pi}) -- (0,{2*pi});
  \draw[dotted] ({1+exp(2)},0) -- ({1+exp(2)},{-2*pi}) -- (0,{-2*pi});
  \fill[coloreschema] (2,{2*pi}) circle (1.8pt);
  \fill[coloreschema] ({1+exp(2)},{-2*pi}) circle (1.8pt);
  \node[below] at (2,0) {$2$};
  \node[above] at ({1+exp(2)},0) {$1+e^2$};
  \node[left] at (0,{2*pi}) {$2\pi$};
  \node[left] at (0,{-2*pi}) {$-2\pi$};
\end{tikzpicture}

Il grafico collega gli estremi inclusi $(2,2\pi)$ e $(1+e^2,-2\pi)$.
\end{schema}



% !TeX root = ../main.tex
% Contenuto della Lezione 8 importato dalla trascrizione già revisionata.
% Convenzioni e preambolo sono ora centralizzati; i tre TikZ sono conservati.
% L'introduzione dalla Lezione 7 precede le proprietà della Lezione 8.
\chapter{Topologia della retta}
\label{cap:topologia}
\fonti{Lezioni 7 e 8}

% !TeX root = ../main.tex
% Fonte: Lezione 7.pdf, pp. 19--21.
\section{Insiemi aperti e insiemi chiusi}

\begin{definizione}[Insieme aperto]
Un insieme $A\subseteq\R$ si dice \emph{aperto} se è un'unione di intervalli aperti: esiste una famiglia $(I_\alpha)_{\alpha\in J}$ di intervalli aperti tale che
\[
A=\bigcup_{\alpha\in J}I_\alpha.
\]
\end{definizione}

\textbf{Famiglie ammesse.} La famiglia può avere un numero finito o infinito di elementi; si ammette anche la famiglia vuota, la cui unione è $\varnothing$.

\textbf{Esempi di aperti.}
\begin{itemize}
  \item $\R=(-\infty,+\infty)$ è aperto;
  \item $\varnothing$ è aperto, come unione della famiglia vuota;
  \item $\R\setminus\Z=\bigcup_{n\in\Z}(n,n+1)$ è aperto;
  \item $\R^*=\R\setminus\{0\}=(-\infty,0)\cup(0,+\infty)$ è aperto.
\end{itemize}
L'insieme di tutti gli aperti si chiama \emph{topologia usuale} di $\R$.

\begin{definizione}[Insieme chiuso]
Un insieme $B\subseteq\R$ si dice \emph{chiuso} se il complemento $B^c=\R\setminus B$ è aperto.
\end{definizione}

\textbf{Esempi di chiusi.} Per verificare che un insieme è chiuso basta quindi descrivere il complemento come unione di intervalli aperti:
\begin{itemize}
  \item $\R$ è chiuso, perché $\R^c=\varnothing$ è aperto;
  \item $\varnothing$ è chiuso, perché $\varnothing^c=\R$ è aperto;
  \item per $a\le b$, $[a,b]$ è chiuso, perché $[a,b]^c=(-\infty,a)\cup(b,+\infty)$;
  \item $\{0\}$ è chiuso, perché il suo complemento è $(-\infty,0)\cup(0,+\infty)$;
  \item $[a,+\infty)$ è chiuso, perché ha complemento $(-\infty,a)$.
\end{itemize}
\textbf{Rapporto fra le due proprietà.} Aperto e chiuso non sono proprietà mutuamente esclusive: $\varnothing$ e $\R$ le possiedono entrambe. Esistono anche insiemi che non possiedono nessuna delle due, come vedremo per $[0,1)$.

\section{Il criterio locale di apertura}

\begin{proposizione}[Caratterizzazione degli aperti]
\label{prop:aperti-criterio}
Un insieme $A\subseteq\R$ è aperto se e solo se
\[
\forall x\in A\ \exists\varepsilon>0:
\quad (x-\varepsilon,x+\varepsilon)\subseteq A.
\]
\end{proposizione}
\begin{proof}
\passomath{\Rightarrow} Se $A$ è un'unione di intervalli aperti e $x\in A$, allora $x$ appartiene a uno di tali intervalli, diciamo $I$. Poiché $x$ non è un estremo di $I$, si può scegliere $\varepsilon>0$ abbastanza piccolo affinché $(x-\varepsilon,x+\varepsilon)\subseteq I\subseteq A$. Per $I=(a,b)$ con estremi finiti basta, per esempio, $\varepsilon=\min\{x-a,b-x\}/2$; per una semiretta basta considerare la distanza dall'unico estremo finito, e per $I=\R$ ogni raggio va bene.

\passomath{\Leftarrow} Per ogni $x\in A$ scegliamo un intervallo $I_x=(x-\varepsilon_x,x+\varepsilon_x)\subseteq A$ come nell'ipotesi. Allora
\[
A=\bigcup_{x\in A}I_x:
\]
\passomath{A\subseteq\bigcup_{x\in A}I_x} Ogni $x\in A$ appartiene al proprio intervallo $I_x$.

\passomath{\bigcup_{x\in A}I_x\subseteq A} Tutti gli intervalli $I_x$ sono contenuti in $A$. È dunque una rappresentazione di $A$ come unione di intervalli aperti. Se $A=\varnothing$, entrambi i criteri valgono con la convenzione della famiglia vuota.
\end{proof}

La proprietà è locale: intorno a ogni punto di $A$ si deve poter ritagliare un piccolo intervallo che resti interamente in $A$. Il raggio può dipendere dal punto; non serve un unico raggio valido per tutti.

\begin{esempio}
\textbf{Non apertura.} L'insieme $A=[0,1)$ non è aperto: $0\in A$, ma ogni intervallo $(-\varepsilon,\varepsilon)$ contiene punti negativi, esterni ad $A$.

\textbf{Non chiusura.} Non è neppure chiuso, perché
\[
A^c=(-\infty,0)\cup[1,+\infty)
\]
non è aperto: contiene $1$, ma nessun intervallo aperto centrato in $1$ è contenuto nel complemento. Basta infatti scegliere, per ogni $\varepsilon>0$, il punto $1-\min\{\varepsilon/2,1/2\}\in A\cap(1-\varepsilon,1+\varepsilon)$.
\end{esempio}


\section{Proprietà degli insiemi aperti}
% Fonte: Lezione 8, pp.~2--3

Ricordiamo il criterio fondamentale: un insieme $A\subseteq\R$ è \emph{aperto}
se, per ogni $x\in A$, esiste $\varepsilon>0$ tale che
$(x-\varepsilon,x+\varepsilon)\subseteq A$. In altre parole, ogni punto di $A$
può essere circondato da un piccolo intervallo interamente contenuto in $A$.

\begin{proposizione}[Operazioni con gli aperti]
Valgono le seguenti proprietà.
\begin{enumerate}
  \item $\varnothing$ e $\R$ sono aperti.
  \item L'intersezione di un numero \emph{finito} di aperti è aperta.
  \item L'unione di una famiglia \emph{arbitraria} di aperti è aperta.
\end{enumerate}
\end{proposizione}

\begin{proof}
Per l'insieme vuoto non ci sono punti sui quali verificare il criterio;
per $\R$, qualunque intervallo centrato in un suo punto resta in $\R$.

Siano $A_1,\ldots,A_m$ aperti e sia $x\in\bigcap_{j=1}^{m}A_j$.
Per ogni $j$ esiste $\varepsilon_j>0$ per cui
$(x-\varepsilon_j,x+\varepsilon_j)\subseteq A_j$.
Poiché gli insiemi sono in numero finito, il minimo
$\varepsilon=\min\{\varepsilon_1,\ldots,\varepsilon_m\}$ è ancora positivo.
L'intervallo $(x-\varepsilon,x+\varepsilon)$ è quindi contenuto in tutti
gli $A_j$, e dunque nella loro intersezione.

Se invece $x$ appartiene a un'unione $\bigcup_{\alpha\in I}A_\alpha$ di aperti,
allora appartiene ad almeno uno di essi, diciamo $A_{\alpha_0}$. Un intervallo
aperto intorno a $x$ contenuto in $A_{\alpha_0}$ è automaticamente contenuto
nell'intera unione. Non serve che la famiglia sia finita.
\end{proof}

L'ipotesi di finitezza nell'intersezione è essenziale. Per esempio,
\[
  \bigcap_{n\in\N^*}\left(-\frac1n,\frac1n\right)=\{0\}.
\]
Il numero $0$ appartiene a tutti gli intervalli. Se $x\ne0$, scegliendo
$n>1/|x|$ si ottiene $1/n<|x|$, perciò $x$ non appartiene all'$n$-esimo
intervallo. Tuttavia $\{0\}$ non è aperto: nessun intervallo aperto centrato
in $0$ può essere contenuto nel solo punto $0$. Un'intersezione infinita
di aperti può quindi non essere aperta.

\Needspace{12\baselineskip}
\section{Proprietà degli insiemi chiusi}
% Fonte: Lezione 8, p.~4

Un insieme $C\subseteq\R$ è \emph{chiuso} quando il suo complemento $C^c$
è aperto. Applicando le leggi di De Morgan alle proprietà appena dimostrate
si ottiene il risultato duale.

\begin{proposizione}[Operazioni con i chiusi]
\begin{enumerate}
  \item $\varnothing$ e $\R$ sono chiusi.
  \item L'intersezione di una famiglia arbitraria di chiusi è chiusa.
  \item L'unione di un numero finito di chiusi è chiusa.
\end{enumerate}
\end{proposizione}

\begin{proof}
I complementi di $\varnothing$ e $\R$ sono rispettivamente $\R$ e
$\varnothing$, che sono aperti. Se $C_\alpha$ sono chiusi, allora
\[
  \left(\bigcap_{\alpha\in I}C_\alpha\right)^c
  =\bigcup_{\alpha\in I}C_\alpha^c
\]
è aperto, perché è un'unione arbitraria di aperti. Analogamente, per
$C_1,\ldots,C_m$ chiusi,
\[
  \left(\bigcup_{j=1}^{m}C_j\right)^c
  =\bigcap_{j=1}^{m}C_j^c
\]
è aperto, perché l'intersezione è finita. In entrambi i casi l'insieme
di partenza è dunque chiuso.
\end{proof}

\section{Esempi nella retta reale}
% Fonte: Lezione 8, pp.~5--9

\subsection{I numeri naturali e gli interi}

L'insieme $\N$ è chiuso. Un modo diretto per vederlo è scrivere il suo
complemento come unione di intervalli aperti:
\[
  \R\setminus\N=(-\infty,0)\ \cup\
  \bigcup_{n\in\N}(n,n+1).
\]
Ogni reale non naturale è negativo oppure si trova strettamente tra due
naturali consecutivi. Il membro destro è aperto, essendo un'unione di aperti;
per definizione, quindi, $\N$ è chiuso.

Lo stesso ragionamento vale per gli interi:
\[
  \R\setminus\Z=\bigcup_{n\in\Z}(n,n+1).
\]
Fra due interi consecutivi non si trovano altri interi, e ogni reale non
intero cade in uno di questi intervalli. Ne segue che $\Z$ è chiuso.
Né $\N$ né $\Z$ sono aperti: ciascuno dei loro punti è isolato rispetto
all'insieme, cioè ogni piccolo intervallo centrato in esso contiene anche
reali che non gli appartengono.

\subsection{I numeri razionali}

$\Q$ non è aperto. La ragione è che in ogni intervallo aperto non vuoto
si trova almeno un numero irrazionale: dunque non esiste alcun intervallo
aperto centrato in un razionale e interamente contenuto in $\Q$.
Questa proprietà si chiama \emph{densità degli irrazionali in $\R$}.

Per chiarire il passaggio, siano $a<b$. Dalla densità di $\Q$ in $\R$
possiamo scegliere un razionale $q\ne0$ tale che
$a/\sqrt2<q<b/\sqrt2$. Allora $a<q\sqrt2<b$ e $q\sqrt2$ è irrazionale:
se fosse razionale, dividendo per il razionale non nullo $q$ si otterrebbe
che $\sqrt2$ è razionale. Abbiamo così trovato un irrazionale in $(a,b)$.

Neppure $\Q$ è chiuso. Infatti, se lo fosse, il complemento
$\R\setminus\Q$ sarebbe aperto. Ma, per la densità dei razionali,
ogni intervallo aperto centrato in un irrazionale contiene razionali;
quindi $\R\setminus\Q$ non è aperto. In sintesi, $\Q$ non è né aperto
né chiuso. La densità di entrambi gli insiemi spiega le due conclusioni.

\subsection{La successione \texorpdfstring{$1/n$}{1/n} vista come insieme}

Consideriamo
\[
  X=\left\{\frac1n:n\in\N^*\right\}
   =\left\{1,\frac12,\frac13,\ldots\right\}.
\]
$X$ non è aperto: per esempio, ogni intervallo aperto centrato in $1$
contiene numeri diversi da tutti gli elementi di $X$.

$X$ non è neppure chiuso. Il punto $0$ non appartiene a $X$, ma è
arbitrariamente vicino ai suoi elementi: per ogni $\varepsilon>0$ esiste
$n\in\N^*$ con $1/n<\varepsilon$. Il complemento di $X$ contiene dunque
$0$, ma non contiene nessun intervallo aperto centrato in $0$;
di conseguenza non è aperto. Esplicitamente,
\[
  X^c=(-\infty,0]\ \cup\
  \bigcup_{n\in\N^*}\left(\frac1{n+1},\frac1n\right)
  \ \cup\ (1,+\infty).
\]
L'estremo $0$, incluso nel primo pezzo, è proprio ciò che impedisce a
$X^c$ di essere aperto.

Se aggiungiamo il punto mancante, otteniamo invece l'insieme chiuso
\[
  Y=X\cup\{0\}.
\]
Infatti
\[
  Y^c=(-\infty,0)\ \cup\
  \bigcup_{n\in\N^*}\left(\frac1{n+1},\frac1n\right)
  \ \cup\ (1,+\infty)
\]
è un'unione di intervalli aperti. L'esempio mostra che l'unione infinita
dei singoletti chiusi $\{1/n\}$ non è necessariamente chiusa, mentre
l'aggiunta del suo punto di accumulazione $0$ la rende tale.

\begin{schema}
\begin{tikzpicture}[x=10cm,y=1cm]
  \foreach \quota/\nome in {.6/X,-.6/Y} {
    \draw[asse] (-.045,\quota) -- (1.065,\quota);
    \node[anchor=east] at (-.075,\quota) {$\nome$};
    \foreach \n in {1,...,9} {
      \fill[coloreschema] ({1/\n},\quota) circle[radius=1.7pt];
    }
    \node[coloreschema,inner sep=0pt] at (.047,\quota) {$\cdots$};
  }
  \draw[estremo] (0,.6) circle[radius=2.5pt];
  \fill[coloreschema] (0,-.6) circle[radius=2.5pt];
  \foreach \pos/\testo in {0/0,.25/{\frac14},.333333/{\frac13},.5/{\frac12},1/1} {
    \node[above=6pt] at (\pos,.6) {$\testo$};
  }
  \node[below=5pt] at (0,-.6) {$0$};
  \node[below=5pt] at (1,-.6) {$1$};
\end{tikzpicture}
\par\smallskip
I punti $1/n$ si avvicinano a $0$: il cerchio vuoto indica che
$0\notin X$, il punto pieno che $0\in Y$. Aggiungendo $0$ si ottiene
l'insieme chiuso $Y=X\cup\{0\}$.
\end{schema}

\section{Intorni di un punto}
% Fonte: Lezione 8, pp.~10--11

\paragraph{Definizione generale.}
Dato $x_0\in\R$, un insieme $U\subseteq\R$ si dice \emph{intorno di $x_0$}
se esiste un aperto $A$ tale che
\[
  x_0\in A\subseteq U.
\]
L'insieme $U$ non deve necessariamente essere aperto: deve soltanto
contenere un insieme aperto che comprenda $x_0$. Per esempio,
$[x_0-1,x_0+1]$ è un intorno di $x_0$, perché contiene
$(x_0-1,x_0+1)$.

Se $U$ è esso stesso aperto e $x_0\in U$, si parla di \emph{intorno aperto}
di $x_0$. In particolare, ogni intervallo $(a,b)$ con $a<x_0<b$
è un intorno aperto di $x_0$. L'insieme di tutti gli intorni di $x_0$
si può indicare con $\mathcal I_{x_0}$.

La definizione precedente equivale a una forma molto concreta:
\[
  U\text{ è un intorno di }x_0
  \quad\Longleftrightarrow\quad
  \exists\,\varepsilon>0:\ (x_0-\varepsilon,x_0+\varepsilon)\subseteq U.
\]
Infatti, se $x_0\in A\subseteq U$ con $A$ aperto, il criterio di apertura
fornisce un tale $\varepsilon$; nel verso opposto, l'intervallo stesso
può fare da insieme aperto $A$.

\paragraph{Intorni sferici.}
Per $r>0$ si definisce l'\emph{intorno sferico}
di centro $x_0$ e raggio $r$ come
\[
  I_r(x_0)=(x_0-r,x_0+r)
  =\{x\in\R:|x-x_0|<r\}.
\]
Sulla retta si tratta semplicemente di un intervallo aperto simmetrico.
Ogni intorno $U$ di $x_0$ contiene un intorno sferico $I_r(x_0)$.
Per esempio, se $(a,b)\subseteq U$ e $a<x_0<b$, basta scegliere
$0<r\le\min\{x_0-a,b-x_0\}$.

\begin{schema}
\begin{tikzpicture}[x=1.2cm,y=1cm]
  \draw[asse] (-.5,0) -- (8.5,0);
  \draw[draw=black!25,line width=4pt] (0,0) -- (8,0);
  \draw[draw=black!55,fill=white,line width=.9pt] (0,0) circle[radius=2.6pt];
  \draw[draw=black!55,fill=white,line width=.9pt] (8,0) circle[radius=2.6pt];
  \draw[intervallo] (1.4,0) -- (4.6,0);
  \foreach \pos in {1.4,4.6} {
    \draw[estremo] (\pos,0) circle[radius=2.6pt];
    \draw[densely dotted,black!45] (\pos,.1) -- (\pos,.65);
  }
  \fill[coloreschema] (3,0) circle[radius=2pt];
  \draw[densely dotted,black!45] (3,.1) -- (3,.65);
  \draw[{Stealth[length=1.5mm]}-{Stealth[length=1.5mm]},coloreschema]
    (1.4,.55) -- node[above] {$r$} (3,.55);
  \draw[{Stealth[length=1.5mm]}-{Stealth[length=1.5mm]},coloreschema]
    (3,.55) -- node[above] {$r$} (4.6,.55);
  \foreach \pos/\testo in {0/a,1.4/{x_0-r},3/{x_0},4.6/{x_0+r},8/b} {
    \node[below=6pt] at (\pos,0) {$\testo$};
  }
  \node[black!70] at (6.5,.5) {$(a,b)\subseteq U$};
\end{tikzpicture}
\par\smallskip
L'intorno sferico $I_r(x_0)$, in blu, è simmetrico rispetto a $x_0$
e resta dentro $(a,b)$; quindi è contenuto anche in $U$.
\end{schema}

\begin{proposizione}[Aperti e intorni]
Un insieme $A\subseteq\R$ è aperto se e solo se è un intorno di ciascuno
dei suoi punti.
\end{proposizione}
\begin{proof}
Se $A$ è aperto e $x\in A$, nella definizione di intorno possiamo
scegliere come aperto proprio $A$. Viceversa, se $A$ è un intorno di
ogni $x\in A$, per ogni tale $x$ esiste $\varepsilon_x>0$ con
$(x-\varepsilon_x,x+\varepsilon_x)\subseteq A$; questo è esattamente
il criterio che definisce un insieme aperto.
\end{proof}

\paragraph{Intorni di $+\infty$ e $-\infty$.}
Quando si lavora nella retta estesa
$\overline\R=\R\cup\{-\infty,+\infty\}$, gli intorni aperti di
$+\infty$ contengono una coda $(a,+\infty]$, e quelli di $-\infty$
contengono una coda $[-\infty,b)$, con $a,b\in\R$.
Se si resta nella sola $\R$, le rispettive code si scrivono
$(a,+\infty)$ e $(-\infty,b)$: i simboli $\pm\infty$ non sono allora
punti dell'insieme ambiente.

\section{Intorni laterali}
% Fonte: Lezione 8, p.~12

Fissato $x_0\in\R$, un \emph{intorno sinistro}
di $x_0$ è un intervallo $(a,x_0)$ con $a<x_0$; un \emph{intorno destro}
è un intervallo $(x_0,b)$ con $b>x_0$. Si possono ammettere anche
$a=-\infty$ e $b=+\infty$.\footnote{Gli intorni laterali qui definiti
escludono il punto $x_0$: non sono quindi intorni di $x_0$ nel senso
della definizione generale. La denominazione serve a descrivere il
comportamento vicino al punto da sinistra o da destra.}

Per esempio, $(-1,1)$ è un intorno sferico di $0$ di raggio $1$.
È anche un intorno ordinario di ogni suo punto, fra cui $0$, $1/2$,
$1/3$, $-1/10$ e $1/5$. Visto invece rispetto ai suoi estremi, $(-1,1)$ è un
intorno destro di $-1$ e un intorno sinistro di $1$. In particolare,
gli intorni laterali di un punto non contengono quel punto.

\begin{schema}
\begin{tikzpicture}[x=1.15cm,y=1cm]
  \foreach \quota in {.7,-.7} {
    \draw[asse] (0,\quota) -- (8,\quota);
    \draw[densely dotted,black!40] (4,{\quota-.25}) -- (4,{\quota+.3});
    \node[below=6pt] at (4,\quota) {$x_0$};
  }
  \draw[intervallo] (1,.7) -- (4,.7);
  \draw[intervallo] (4,-.7) -- (7,-.7);
  \foreach \pos/\quota in {1/.7,4/.7,4/-.7,7/-.7} {
    \draw[estremo] (\pos,\quota) circle[radius=2.6pt];
  }
  \node[below=6pt] at (1,.7) {$a$};
  \node[below=6pt] at (7,-.7) {$b$};
  \node[above=6pt,coloreschema] at (2.5,.7) {intorno sinistro $(a,x_0)$};
  \node[above=6pt,coloreschema] at (5.5,-.7) {intorno destro $(x_0,b)$};
\end{tikzpicture}
\par\smallskip
Negli intorni laterali il punto $x_0$ è escluso,
come mostrano i cerchi vuoti.
\end{schema}

% !TeX root = ../main.tex
\chapter{Successioni e limiti}
\label{cap:successioni}
\fonti{Lezioni 10, 11 e 12}

% !TeX root = ../main.tex
% Fonte: Lezione 10, pp. 7--17. Prove implicite completate redazionalmente.
\section{Proprietà definitive e frequenti}

Nello studio di una successione interessa ciò che accade per indici arbitrariamente grandi. Una proprietà può valere da un certo indice in poi, oppure continuare a ripresentarsi senza diventare stabile: le due situazioni hanno quantificatori diversi.

\begin{definizione}[Definitivamente e frequentemente]
Sia $P(n)$ un predicato definito per i naturali di una semiretta $\{n\in\N:n\ge n_0\}$. Si dice che $P(n)$ è:
\begin{itemize}
  \item \emph{definitivamente vero} se esiste $N\ge n_0$ tale che $P(n)$ è vero per ogni $n\ge N$;
  \item \emph{frequentemente vero} se per ogni $N\ge n_0$ esiste $n\ge N$ tale che $P(n)$ è vero.
\end{itemize}
\end{definizione}

Definitivamente significa che ci sono al più un numero finito di eccezioni; frequentemente significa che gli indici per cui la proprietà vale formano un insieme infinito, quindi illimitato in $\N$. Ogni proprietà definitiva è frequente, ma non vale il contrario. Inoltre, negando i quantificatori,
\[
\neg(P\text{ definitivamente vero})
\ \Longleftrightarrow\ (\neg P)\text{ frequentemente vero}.
\]

\begin{esempio}
Per $n\ge1$, la proprietà $1/n\le1/10$ è definitivamente vera: vale per tutti gli indici $n\ge10$. La proprietà $(-1)^n>0$ è invece vera per tutti gli indici pari, perciò è frequente ma non definitiva, perché è falsa per tutti gli indici dispari.
\end{esempio}

\begin{proposizione}[Congiunzione di proprietà definitive]
\label{prop:congiunzione-definitiva}
Se $P_1(n),\ldots,P_k(n)$, con $k\ge1$, sono definitivamente vere, anche la loro congiunzione è definitivamente vera.
\end{proposizione}
\begin{proof}
Per ogni $j$ esiste una soglia $N_j$ oltre la quale $P_j(n)$ è sempre vera. Ponendo $N=\max\{N_1,\ldots,N_k\}$, per $n\ge N$ risultano contemporaneamente vere tutte le proprietà.
\end{proof}

\begin{esempio}[Una soglia comune]
Le proprietà $n^2\ge25$ e $1/n\le1/10$, per $n\ge1$, valgono rispettivamente da $N_1=5$ e $N_2=10$. La loro congiunzione vale da $N=10$. Per proprietà soltanto frequenti il ragionamento non funziona: «$n$ è pari» e «$n$ è dispari» sono entrambe frequenti, ma non sono mai vere insieme.
\end{esempio}

\section{Successioni e insieme dei valori}
% Fonte: Lezione 10, pp. 10--11 e 16--17.

\begin{definizione}[Successione]
Una \emph{successione a valori in $A$} è una funzione
\[
a\colon S\longrightarrow A,\qquad n\longmapsto a_n,
\quad S=\{n\in\N:n\ge n_0\}.
\]
Il numero $n$ è l'\emph{indice}, $a_n$ il termine di indice $n$. Si usano le notazioni $(a_n)_{n\ge n_0}$, $(a_n)$ e, quando non crea ambiguità, $\{a_n\}$. Per $A=\R$ si parla di \emph{successione reale}.
\end{definizione}

Una successione conserva l'ordine dei termini e le loro ripetizioni; l'insieme dei suoi valori $\{a_n:n\in S\}$ dimentica entrambe le informazioni. Per esempio la successione costante $a_n=1$ ha infiniti termini, ma il suo insieme immagine è il singoletto $\{1\}$.

\begin{esempio}[Formule e dominio degli indici]
Le formule $a_n=1$ e $b_n=2n^2$ definiscono successioni su tutta $\N$. La formula $c_n=1/(n-4)$ può definire una successione sulla semiretta $n\ge5$, sulla quale il denominatore è sempre diverso da zero. Specificare il dominio evita di considerare un termine non definito.
\end{esempio}

\begin{schema}
\begin{tikzpicture}[x=.72cm,y=.9cm]
  \draw[asse] (-.4,0) -- (10.7,0) node[right] {$n$};
  \draw[asse] (0,-.2) -- (0,1.8) node[above] {$a_n$};
  \foreach \n in {0,...,10} {
    \fill[coloreschema] (\n,1) circle[radius=1.8pt];
  }
  \foreach \n in {0,2,4,6,8,10} {\node[below=4pt] at (\n,0) {$\n$};}
  \node[left] at (0,1) {$1$};
\end{tikzpicture}
\par\smallskip Il grafico di una successione è formato da punti con ascissa intera; qui ogni indice ha valore $1$.
\end{schema}

Per una successione reale si riprendono le nozioni già date per le funzioni: è \emph{superiormente limitata} se $a_n\le M$ per ogni indice e per uno stesso $M\in\R$; è \emph{inferiormente limitata} se $a_n\ge m$ per ogni indice. È \emph{limitata} se esiste $K>0$ con $|a_n|\le K$ per ogni indice.

\begin{proposizione}[Da una coda limitata all'intera successione]
\label{prop:coda-limitata}
Una successione reale definitivamente superiormente limitata è superiormente limitata su tutto il dominio. Valgono anche le versioni con limitatezza inferiore e con limitatezza in valore assoluto.
\end{proposizione}
\begin{proof}
Supponiamo $a_n\le M$ per $n\ge N$. Il numero
\[
M'=\max\bigl(\{M\}\cup\{a_n:n_0\le n<N\}\bigr)
\]
esiste perché l'insieme è finito e non vuoto. Esso maggiora i termini iniziali e anche tutti quelli della coda, quindi maggiora l'intera successione. Per la limitatezza inferiore si usa il minimo insieme a un minorante della coda. Per quella in valore assoluto si applica lo stesso ragionamento alla successione $|a_n|$.
\end{proof}

La proposizione usa il fatto che prima di una soglia naturale ci sono solo finiti termini. Un vincolo su un solo lato non dà invece limitatezza in valore assoluto: per esempio $a_n=-n$ è superiormente limitata, ma non inferiormente limitata.

\begin{schema}
\begin{tikzpicture}[x=.72cm,y=.65cm]
  \draw[asse] (-.4,0) -- (10.7,0) node[right] {$n$};
  \draw[asse] (0,-2.5) -- (0,3.3);
  \draw[dashed,black!50] (0,1.2) -- (10.3,1.2) node[right] {$K$};
  \draw[dashed,black!50] (0,-1.2) -- (10.3,-1.2) node[right] {$-K$};
  \draw[dotted] (4,-2.2) -- (4,3);
  \node[above] at (4,3) {$N$};
  \foreach \n/\v in {0/2.6,1/-2,2/1.9,3/2.3,4/.8,5/-.7,6/.5,7/-.6,8/.4,9/-.3,10/.2} {
    \fill[coloreschema] (\n,\v) circle[radius=1.8pt];
  }
\end{tikzpicture}
\par\smallskip Dopo $N$ i valori restano fra $-K$ e $K$. I finiti termini precedenti si includono ampliando il vincolo.
\end{schema}

\section{Sottosuccessioni}
% Fonte: Lezione 10, pp. 12--15.

\begin{definizione}[Sottosuccessione o successione estratta]
Data $a\colon S\to A$, una \emph{sottosuccessione} è la composta $a\circ k$, dove $k\colon T\to S$ è strettamente crescente e $T$ è una semiretta di $\N$. Scrivendo $k_j=k(j)$, i termini estratti sono $a_{k_j}$, con
\[
k_{j+1}>k_j.
\]
Si scelgono quindi infiniti indici preservando il loro ordine.
\end{definizione}

\begin{esempio}[Scelte di indici]
Se $(a_n)$ è definita su tutta $\N$, sono sottosuccessioni
\[
(a_j),\qquad (a_{2j}),\qquad (a_{13j+4}),\qquad (a_{j+1}),
\qquad j\in\N.
\]
Per esempio $(a_{13j+4})=(a_4,a_{17},a_{30},\ldots)$. Non sono scelte ammesse la lista $a_2,a_6,a_{11},a_3,\ldots$, che torna a un indice precedente, né la lista $a_0,a_5,a_{11},a_6,a_{21},a_{19},\ldots$, che non conserva l'ordine. La condizione riguarda gli indici: due termini estratti possono avere lo stesso valore.
\end{esempio}

\begin{schema}
\begin{tikzpicture}[x=.8cm,y=.65cm]
  \draw[asse] (-.4,0) -- (9.7,0) node[right] {$n$};
  \draw[asse] (0,-1.8) -- (0,2.8);
  \foreach \n/\v in {0/.6,1/1.1,2/1.8,3/.5,4/-1.2,5/-1,6/-.5,7/.8,8/1.2,9/2.4} {
    \pgfmathtruncatemacro{\pari}{mod(\n,2)}
    \ifnum\pari=0
      \fill[coloreschema] (\n,\v) circle[radius=2.4pt];
    \else
      \draw[draw=black!60,fill=white] (\n,\v) circle[radius=2.4pt];
    \fi
    \node[below=4pt] at (\n,-1.7) {$\n$};
  }
\end{tikzpicture}
\par\smallskip I punti pieni selezionano $(a_{2j})$, i punti vuoti $(a_{2j+1})$; ciascuna estrazione mantiene l'ordine degli indici.
\end{schema}

Gli indici estratti diventano arbitrariamente grandi. Infatti, se $T=\{j\ge j_0\}$, la stretta crescita fra interi implica per induzione
\[
k_j\ge k_{j_0}+j-j_0.
\]
In particolare, indicizzando da $0$, si ha $k_j\ge j$. La crescita richiesta è quella di $k_j$, non quella dei valori $a_{k_j}$.

\begin{proposizione}[Estrazione da una proprietà frequente]
Se $P(a_n)$ è frequentemente vera, esiste una sottosuccessione $(a_{k_j})$ per la quale $P(a_{k_j})$ è vera a ogni indice.
\end{proposizione}
\begin{proof}
Scegliamo un indice $k_0$ per cui $P(a_{k_0})$ è vera. Dopo aver scelto $k_j$, la frequenza applicata alla soglia $k_j+1$ fornisce un indice $k_{j+1}\ge k_j+1$ per il quale la proprietà è ancora vera. Si costruisce così una funzione di estrazione strettamente crescente.
\end{proof}

\begin{esempio}
Per $a_n=(-1)^n$ la proprietà $a_n>0$ è frequentemente vera. L'estrazione $k_j=2j$ dà $a_{k_j}=1>0$ per ogni $j$.
\end{esempio}

% !TeX root = ../main.tex
% Fonti: Lezione 10, p. 18; Lezione 11, pp. 2--17; Lezione 12, p. 10.
% La definizione comune alle lezioni 10 e 11 è esposta una sola volta.
\Needspace{22\baselineskip}
\section{Limite di una successione}

\begin{definizione}[Limite]
Sia $(a_n)$ una successione reale. Si dice che ha \emph{limite} $\ell\in\overline\R$ e si scrive
\[
a_n\longrightarrow\ell,\qquad
\lim_{n\to+\infty}a_n=\ell,
\]
se, per ogni intorno $U$ di $\ell$ nella retta estesa, i termini $a_n$ appartengono definitivamente a $U$:
\[
\forall U\in\mathcal I_\ell\ \exists N\in\N:
\quad \forall n\ge N,\quad a_n\in U.
\]
Si sottintende che gli indici considerati appartengano al dominio della successione.
\end{definizione}

Il limite descrive una proprietà definitiva per \emph{ogni} intorno del candidato $\ell$. La soglia può dipendere dall'intorno: non occorre trovare un unico indice che funzioni contemporaneamente per tutti gli intorni. Cambiare o eliminare finiti termini iniziali non modifica il limite, perché le soglie possono sempre essere aumentate.

\begin{definizione}[Successioni convergenti, divergenti e irregolari]
Una successione si dice \emph{convergente} se ha un limite finito $\ell\in\R$; se $\ell=0$ è detta anche \emph{infinitesima}. Si dice \emph{divergente a $+\infty$} o \emph{a $-\infty$} se ha il corrispondente limite infinito. Le successioni convergenti o divergenti sono dette \emph{regolari}; quelle che non hanno limite neppure in $\overline\R$ si dicono \emph{irregolari}.
\end{definizione}

Con queste convenzioni, «non convergente» non significa necessariamente «divergente»: può mancare qualsiasi limite oppure esistere un limite infinito.

\subsection{La forma con epsilon e le code a infinito}

Per un limite finito, ogni intorno di $\ell$ contiene un intervallo $(\ell-\varepsilon,\ell+\varepsilon)$ con $\varepsilon>0$. Usare questi soli intorni è quindi sufficiente, e la definizione diventa
\[
a_n\to\ell\in\R
\quad\Longleftrightarrow\quad
\forall\varepsilon>0\ \exists N_\varepsilon\in\N:
\quad \forall n\ge N_\varepsilon,\quad |a_n-\ell|<\varepsilon.
\]
Infatti la definizione per tutti gli intorni dà in particolare questa condizione; viceversa, dato un intorno qualunque, basta prendere un intervallo sferico contenuto in esso.

Analogamente, gli intorni a infinito si riducono alle code. Si ottiene
\[
\begin{aligned}
a_n\to+\infty
&\Longleftrightarrow\ \forall M>0\ \exists N_M\in\N:
\quad \forall n\ge N_M,\quad a_n>M,\\
a_n\to-\infty
&\Longleftrightarrow\ \forall M>0\ \exists N_M\in\N:
\quad \forall n\ge N_M,\quad a_n<-M.
\end{aligned}
\]
Richiedere solo $M>0$ è sufficiente: aumentando una soglia positiva si coprono anche le code con estremo reale arbitrario.

\begin{schema}
\begin{tikzpicture}[x=.65cm,y=1cm]
  \draw[asse] (-.4,0) -- (12,0) node[right] {$n$};
  \draw[asse] (0,-.2) -- (0,2.7);
  \draw[black!45] (0,2) -- (11.6,2);
  \draw[dashed,coloreschema] (0,1.6) -- (11.6,1.6);
  \draw[dashed,coloreschema] (0,2.4) -- (11.6,2.4);
  \node[left] at (0,2) {$\ell$};
  \node[left] at (0,1.6) {$\ell-\varepsilon$};
  \node[left] at (0,2.4) {$\ell+\varepsilon$};
  \draw[dotted] (4,0) -- (4,2.6);
  \node[below] at (4,0) {$N_\varepsilon$};
  \foreach \n/\v in {1/.5,2/.9,3/1.4,4/1.7,5/1.8,6/2.1,7/1.9,8/2.08,9/1.94,10/2.03,11/1.98} {
    \fill[coloreschema] (\n,\v) circle[radius=1.8pt];
  }
\end{tikzpicture}
\par\smallskip Per il raggio scelto, dopo $N_\varepsilon$ tutti i termini restano nella fascia; una fascia più stretta può richiedere una soglia maggiore.
\end{schema}

\begin{schema}
\begin{tikzpicture}[x=.65cm,y=.45cm]
  \draw[asse] (-.4,0) -- (12,0) node[right] {$n$};
  \draw[asse] (0,-.4) -- (0,6.2);
  \draw[dashed,coloreschema] (0,2.6) -- (11.5,2.6) node[right] {$M$};
  \draw[dotted] (6,0) -- (6,5.7);
  \node[below] at (6,0) {$N_M$};
  \foreach \n/\v in {1/.3,2/.8,3/1.2,4/1.6,5/2.2,6/3,7/3.4,8/4,9/4.5,10/5,11/5.5} {
    \fill[coloreschema] (\n,\v) circle[radius=1.8pt];
  }
\end{tikzpicture}
\par\smallskip Per tendere a $+\infty$ i termini devono superare definitivamente ogni livello $M$, non soltanto assumere valori arbitrariamente grandi.
\end{schema}

\subsection{Primi limiti e verifiche dalla definizione}
% Fonte: Lezione 11, pp. 5--9; soglie naturali rese esplicite.

\begin{esempio}[Costante, reciproco e indice]
La successione costante $a_n=1$ converge a $1$: ogni intorno di $1$ contiene tutti i suoi termini. Per $n\ge1$, la successione $b_n=1/n$ converge a $0$: dato $\varepsilon>0$, scegliamo un naturale $N_\varepsilon>1/\varepsilon$, così per $n\ge N_\varepsilon$ si ha $|b_n|=1/n<\varepsilon$. Infine $c_n=n$ diverge a $+\infty$: dato $M>0$, basta un naturale $N_M>M$.
\end{esempio}

\begin{esercizio}[Un limite finito]
Verificare dalla definizione che $1-1/n^2\to1$, per $n\ge1$, e mostrare direttamente perché il candidato limite $2$ è impossibile.
\end{esercizio}
\begin{svolgimento}
Per $\varepsilon>0$ si ha
\[
\left|\left(1-\frac1{n^2}\right)-1\right|<\varepsilon
\quad\Longleftrightarrow\quad n>\frac1{\sqrt\varepsilon}.
\]
Scegliamo quindi un \emph{naturale} $N_\varepsilon>1/\sqrt\varepsilon$; per ogni $n\ge N_\varepsilon$ vale la disuguaglianza richiesta. La scelta, per esempio, $N_\varepsilon=\lfloor1/\sqrt\varepsilon\rfloor+1$ garantisce anche la stretta disuguaglianza quando $1/\sqrt\varepsilon$ è intero.

Per il candidato $2$ si avrebbe invece
\[
\left|\left(1-\frac1{n^2}\right)-2\right|
=1+\frac1{n^2}>1.
\]
Con $\varepsilon=1/2$ nessun termine soddisfa la condizione, perciò non esiste la soglia richiesta e $2$ non è un limite.
\end{svolgimento}

\begin{esercizio}[Un limite infinito]
Verificare dalla definizione che $n^3-1\to+\infty$.
\end{esercizio}
\begin{svolgimento}
Per $M>0$ si ha $n^3-1>M$ se e solo se $n>\sqrt[3]{M+1}$. Scegliendo un naturale $N_M>\sqrt[3]{M+1}$, la condizione vale per ogni $n\ge N_M$.
\end{svolgimento}

\Needspace{16\baselineskip}
\section{Unicità ed estrazione di sottosuccessioni}
% Fonte: Lezione 11, pp. 10--14; separazione ripresa dal capitolo 7.

\begin{teorema}[Unicità del limite]
Una successione reale ha al più un limite in $\overline\R$.
\end{teorema}
\begin{proof}
Supponiamo per assurdo che $a_n\to\ell_1$ e $a_n\to\ell_2$, con $\ell_1\ne\ell_2$. La proprietà di separazione nella retta estesa fornisce intorni disgiunti $U_1$ e $U_2$ dei due punti. Dalle definizioni di limite, $a_n\in U_1$ e $a_n\in U_2$ sono entrambe proprietà definitive. Per la proposizione~\ref{prop:congiunzione-definitiva}, oltre la massima delle due soglie si avrebbe $a_n\in U_1\cap U_2=\varnothing$, assurdo.
\end{proof}

\begin{schema}
\begin{tikzpicture}[x=.65cm,y=.85cm]
  \draw[asse] (-.4,0) -- (12,0) node[right] {$n$};
  \draw[asse] (0,-.2) -- (0,4.7);
  \foreach \v in {1,2,3,4} {\draw[dashed,black!45] (0,\v) -- (11.5,\v);}
  \node[left] at (0,1.5) {$U_1$};
  \node[left] at (0,3.5) {$U_2$};
  \draw[dotted,coloreschema] (5,0) -- (5,4.5);
  \draw[dotted,black!60] (8,0) -- (8,4.5);
  \node[below] at (5,0) {$N_1$};
  \node[below] at (8,0) {$N_2$};
  \draw[{Stealth}-{Stealth},coloreschema] (9,2.1) -- (9,2.9);
\end{tikzpicture}
\par\smallskip Oltre $\max\{N_1,N_2\}$ la definizione imporrebbe a ogni termine di stare in entrambe le fasce disgiunte, condizione impossibile.
\end{schema}

\begin{proposizione}[Limite delle sottosuccessioni]
Se $a_n\to\ell\in\overline\R$, ogni sottosuccessione $(a_{k_j})$ tende allo stesso limite $\ell$.
\end{proposizione}
\begin{proof}
Sia $U$ un intorno di $\ell$. Esiste $N$ tale che $a_n\in U$ per ogni $n\ge N$. Gli indici $k_j$ sono strettamente crescenti, dunque diventano arbitrariamente grandi: per la stima $k_j\ge k_{j_0}+j-j_0$ esiste $J$ tale che $k_j\ge N$ per ogni $j\ge J$. Allora $a_{k_j}\in U$ per ogni $j\ge J$, cioè la sottosuccessione tende a $\ell$.
\end{proof}

\begin{corollario}[Criterio di non esistenza del limite]
Se due sottosuccessioni di $(a_n)$ hanno limiti diversi in $\overline\R$, la successione $(a_n)$ non ha limite.
\end{corollario}
\begin{proof}
Un eventuale limite dell'intera successione sarebbe il limite di entrambe le sottosuccessioni. L'unicità del limite rende impossibile che esse abbiano limiti diversi.
\end{proof}

\begin{esempio}[Successioni oscillanti]
La successione $a_n=(-1)^n$ non ha limite: $a_{2j}=1\to1$ e $a_{2j+1}=-1\to-1$. È dunque irregolare, pur essendo limitata.

Anche $b_n=\sin(n\pi/2)$ è irregolare: le estrazioni $b_{4j}=0$, $b_{4j+1}=1$ e $b_{4j+3}=-1$ sono costanti con tre limiti diversi. Le ripetizioni dei valori non impediscono l'estrazione: sono gli indici a essere strettamente crescenti.
\end{esempio}

\section{Limitatezza e permanenza del segno}
% Fonte: Lezione 11, pp. 15--17.

\begin{teorema}[Una successione convergente è limitata]
Se $a_n\to\ell\in\R$, esiste $K>0$ con $|a_n|\le K$ per ogni indice del dominio.
\end{teorema}
\begin{proof}
Usando $\varepsilon=1$ nella definizione di limite, definitivamente si ha $|a_n-\ell|<1$, dunque $|a_n|\le|a_n-\ell|+|\ell|<1+|\ell|$. La coda è limitata e la proposizione~\ref{prop:coda-limitata} estende il vincolo all'intera successione.
\end{proof}

\begin{osservazione}
La limitatezza non implica convergenza, come mostra $(-1)^n$. Neppure l'illimitatezza implica divergenza a uno dei due infiniti: per $a_n=(-1)^n n$, la sottosuccessione pari tende a $+\infty$ e quella dispari a $-\infty$, perciò il limite non esiste neanche nella retta estesa.
\end{osservazione}

\begin{teorema}[Permanenza del segno]
\label{teo:permanenza-segno}
\begin{enumerate}[label=\arabic*.]
\item Se $a_n\to\ell\in\R$ e $\ell>0$, allora $a_n>0$ definitivamente.
\item Se $\ell<0$, allora $a_n<0$ definitivamente.
\end{enumerate}
\end{teorema}
\begin{proof}
\begin{enumerate}[label=\arabic*.]
\item Per $\ell>0$, scegliamo $\varepsilon=\ell/2$. Definitivamente
\[
|a_n-\ell|<\ell/2
\quad\Longrightarrow\quad \ell/2<a_n<3\ell/2,
\]
quindi $a_n>0$.
\item Per $\ell<0$, scegliendo $\varepsilon=|\ell|/2$ si ottiene definitivamente $a_n<\ell+|\ell|/2=\ell/2<0$.
\end{enumerate}
\end{proof}

La conclusione vale anche per i limiti infiniti: se $a_n\to+\infty$, definitivamente $a_n>1>0$; se $a_n\to-\infty$, definitivamente $a_n<-1<0$.

\begin{corollario}[Segno del limite]
\begin{enumerate}[label=\arabic*.]
\item Se $a_n\ge0$ definitivamente e $a_n\to\ell\in\R$, allora $\ell\ge0$.
\item Analogamente, da $a_n\le0$ definitivamente segue $\ell\le0$.
\end{enumerate}
\end{corollario}
\begin{proof}
\begin{enumerate}[label=\arabic*.]
\item Se il limite fosse negativo, per la permanenza del segno i termini sarebbero definitivamente negativi, in contraddizione con l'ipotesi.
\item Il secondo caso è simmetrico.
\end{enumerate}
\end{proof}

Non si può concludere la stretta positività del limite dalla stretta positività dei termini: $1/n>0$ per ogni $n\ge1$, ma $1/n\to0$.

% !TeX root = ../main.tex
% Fonte: Lezione 12, pp. 2--10; caso sin(n) da Lezione 11, p. 14.
% Completate le prove lasciate implicite: confronti a infinito, prodotto,
% casi infiniti dell'algebra. Il teorema sui quozienti non è nelle fonti.
\section{Teoremi di confronto}

\begin{proposizione}[Confronto con una successione divergente]
Siano $a_n\le b_n$ definitivamente. Se $a_n\to+\infty$, allora $b_n\to+\infty$. Se $b_n\to-\infty$, allora $a_n\to-\infty$.
\end{proposizione}
\begin{proof}
Nel primo caso fissiamo $M>0$. Definitivamente si hanno sia $a_n>M$ sia $a_n\le b_n$; oltre la massima delle due soglie risulta $b_n\ge a_n>M$. Questo è il criterio per $b_n\to+\infty$.

Nel secondo caso si hanno definitivamente $b_n<-M$ e $a_n\le b_n$, quindi $a_n<-M$ definitivamente, per ogni $M>0$.
\end{proof}

\begin{teorema}[Teorema dei carabinieri]
\label{teo:carabinieri}
Siano $(a_n)$, $(b_n)$ e $(c_n)$ successioni reali con $a_n\le b_n\le c_n$ definitivamente. Se $a_n\to\ell$ e $c_n\to\ell$, con $\ell\in\overline\R$, allora anche $b_n\to\ell$.
\end{teorema}
\begin{proof}
\textbf{Limite finito.} Per $\varepsilon>0$, definitivamente $\ell-\varepsilon<a_n<\ell+\varepsilon$ e $\ell-\varepsilon<c_n<\ell+\varepsilon$. Oltre una soglia comune, che assicuri anche i confronti, si ha
\[
\ell-\varepsilon<a_n\le b_n\le c_n<\ell+\varepsilon.
\]
Quindi $|b_n-\ell|<\varepsilon$ definitivamente.

\textbf{Limiti infiniti.} Se $\ell=+\infty$, basta il confronto $a_n\le b_n$; se $\ell=-\infty$, basta $b_n\le c_n$. Si applica la proposizione precedente.
\end{proof}

\begin{schema}
\begin{tikzpicture}[x=.72cm,y=.85cm]
  \draw[asse] (-.3,0) -- (10.7,0) node[right] {$n$};
  \draw[asse] (0,-.2) -- (0,3.2);
  \draw[dashed,black!45] (0,1.6) -- (10.3,1.6) node[right] {$\ell$};
  \foreach \n in {1,...,10} {
    \pgfmathsetmacro{\d}{1.25/\n}
    \draw[draw=black!50,fill=white] (\n,{1.6-\d}) circle[radius=1.8pt];
    \draw[draw=black!50,fill=white] (\n,{1.6+\d}) circle[radius=1.8pt];
    \fill[coloreschema] (\n,{1.6+cos(\n*65)*\d*.7}) circle[radius=1.8pt];
  }
\end{tikzpicture}
\par\smallskip I valori della successione centrale restano fra due successioni che tendono allo stesso limite; anche la fascia che li contiene si restringe verso $\ell$.
\end{schema}

\section{Somme con successioni divergenti}
% Fonte: Lezione 12, p. 7; prova assegnata per casa.

\begin{esercizio}[Divergente più limitata da un lato]
Dimostrare che, se $a_n\to+\infty$ e $(b_n)$ è inferiormente limitata, allora $a_n+b_n\to+\infty$.
\end{esercizio}
\begin{svolgimento}
Sia $m\in\R$ un minorante per tutti i $b_n$. Dato $M>0$, dalla divergenza di $a_n$ si ottiene definitivamente $a_n>M-m$; se $M-m$ non è positivo, basta prima imporre una soglia positiva più grande. Pertanto
\[
a_n+b_n\ge a_n+m>M
\]
definitivamente. Questo vale per ogni $M>0$, quindi la somma tende a $+\infty$.
\end{svolgimento}

L'ipotesi su $b_n$ può valere soltanto definitivamente, perché si possono assorbire i termini iniziali mediante la proposizione~\ref{prop:coda-limitata}. Cambiando i segni si ottiene anche: se $a_n\to-\infty$ e $b_n$ è superiormente limitata, allora $a_n+b_n\to-\infty$.

\begin{esempio}
Per $a_n=n$ e $b_n=(-1)^n$ si ha $b_n\ge-1$, quindi $n+(-1)^n\to+\infty$. Il secondo addendo non ha limite, ma il suo vincolo inferiore è sufficiente.
\end{esempio}

\begin{schema}
\begin{tikzpicture}[x=.72cm,y=.35cm]
  \draw[asse] (-.4,0) -- (10.7,0) node[right] {$n$};
  \draw[asse] (0,-.5) -- (0,12);
  \foreach \n in {1,...,10} {
    \draw[draw=black!55,fill=white] (\n,\n) circle[radius=1.8pt];
    \fill[coloreschema] (\n,{\n+cos(\n*180)}) circle[radius=2pt];
  }
\end{tikzpicture}
\par\smallskip I punti pieni rappresentano $n+(-1)^n$, quelli vuoti $n$: le oscillazioni del secondo addendo non impediscono la divergenza della somma.
\end{schema}

\section{Prodotto di una limitata per un'infinitesima}
% Fonte: Lezione 12, pp. 8--9. La stima per K|a_n| è giustificata direttamente.

\begin{proposizione}
\label{prop:limitata-infinitesima}
Se $a_n\to0$ e $(b_n)$ è limitata, allora $a_nb_n\to0$.
\end{proposizione}
\begin{proof}
Scegliamo $K>0$ tale che $|b_n|\le K$ per ogni indice. Si ha
\[
0\le|a_nb_n|=|a_n|\,|b_n|\le K|a_n|.
\]
La successione $K|a_n|$ tende a $0$: dato $\varepsilon>0$, basta usare la definizione di $a_n\to0$ con il raggio $\varepsilon/K$. Per il teorema dei carabinieri segue $|a_nb_n|\to0$. Infine
\[
-|a_nb_n|\le a_nb_n\le|a_nb_n|,
\]
e le due successioni esterne tendono a $0$; applicando nuovamente i carabinieri si conclude $a_nb_n\to0$.
\end{proof}

La limitatezza è richiesta solo da un certo indice in poi, perché essa si estende a tutta la successione; analogamente, per la dimostrazione basta che le stime valgano definitivamente. Il risultato non richiede che $(b_n)$ abbia limite.

\begin{schema}
\begin{tikzpicture}[x=.72cm,y=.8cm]
  \draw[asse] (-.4,0) -- (10.7,0) node[right] {$n$};
  \draw[asse] (0,-1.8) -- (0,1.8);
  \foreach \n in {1,...,10} {
    \draw[draw=black!50,fill=white] (\n,{1.5/\n}) circle[radius=1.8pt];
    \draw[draw=black!50,fill=white] (\n,{-1.5/\n}) circle[radius=1.8pt];
    \fill[coloreschema] (\n,{cos(\n*180)/\n}) circle[radius=1.8pt];
  }
\end{tikzpicture}
\par\smallskip Un prodotto con fattore limitato può oscillare di segno, ma il suo modulo resta sotto un'infinitesima: i due vincoli tendono entrambi a $0$.
\end{schema}

\Needspace{18\baselineskip}
\section{Algebra dei limiti}
% Fonte: Lezione 12, pp. 4--6; completate le prove dei prodotti e dei casi infiniti.

\subsection{Operazioni nella retta estesa}

Nella retta estesa alcune operazioni si prolungano secondo le seguenti regole:
\[
\begin{aligned}
x+(+\infty)&=+\infty &&\text{se }x\ne-\infty,\\
x+(-\infty)&=-\infty &&\text{se }x\ne+\infty.
\end{aligned}
\]
Per i prodotti vale la consueta regola dei segni: un fattore positivo, finito o infinito, moltiplicato per $+\infty$ dà $+\infty$, e moltiplicato per $-\infty$ dà $-\infty$; un fattore negativo inverte i segni. In particolare
\[
\begin{array}{c|cc}
\cdot & +\infty & -\infty\\ \hline
x>0 & +\infty & -\infty\\
x<0 & -\infty & +\infty
\end{array}
\]
dove il fattore $x$ può essere anche infinito. Rimangono escluse le espressioni
\[
(+\infty)+(-\infty),\qquad (-\infty)+(+\infty),
\qquad 0\cdot(\pm\infty).
\]
La retta estesa non è quindi un campo: queste operazioni non sono definite per ogni coppia di elementi.

\subsection{Limite della somma}

\begin{teorema}
Siano $a_n\to\alpha$ e $b_n\to\beta$, con $\alpha,\beta\in\overline\R$. Se la somma $\alpha+\beta$ è definita, allora
\[
a_n+b_n\longrightarrow\alpha+\beta.
\]
\end{teorema}
\begin{proof}
\textbf{Due limiti finiti.} Dato $\varepsilon>0$, definitivamente valgono insieme $|a_n-\alpha|<\varepsilon/2$ e $|b_n-\beta|<\varepsilon/2$. Per la disuguaglianza triangolare,
\[
|(a_n+b_n)-(\alpha+\beta)|
\le|a_n-\alpha|+|b_n-\beta|<\varepsilon.
\]

\textbf{Almeno un limite infinito.} Se, per esempio, $\alpha=+\infty$ e $\beta\in\R\cup\{+\infty\}$, la successione $(b_n)$ è inferiormente limitata: se converge è limitata; se tende a $+\infty$ è definitivamente positiva e poi si applica la proposizione~\ref{prop:coda-limitata}. Il risultato sulla somma con una successione inferiormente limitata dà $a_n+b_n\to+\infty$. Si scambiano gli addendi se è $\beta$ a essere infinito. Per il limite $-\infty$ si usa la versione con limitatezza superiore. Sono così coperti tutti i casi in cui la somma dei limiti è definita.
\end{proof}

\subsection{Limite del prodotto}

\begin{teorema}
Siano $a_n\to\alpha$ e $b_n\to\beta$, con $\alpha,\beta\in\overline\R$. Se il prodotto $\alpha\beta$ è definito, allora
\[
a_nb_n\longrightarrow\alpha\beta.
\]
\end{teorema}
\begin{proof}
\textbf{Due limiti finiti.} La successione $(b_n)$ è limitata. Dalla definizione, $a_n-\alpha\to0$ e $b_n-\beta\to0$; per la proposizione~\ref{prop:limitata-infinitesima},
\[
(a_n-\alpha)b_n\to0,\qquad \alpha(b_n-\beta)\to0.
\]
Poiché
\[
a_nb_n-\alpha\beta=(a_n-\alpha)b_n+\alpha(b_n-\beta),
\]
la regola della somma già dimostrata dà $a_nb_n-\alpha\beta\to0$, ossia $a_nb_n\to\alpha\beta$.

\Needspace{14\baselineskip}
\textbf{Un limite finito non nullo e uno infinito.} Supponiamo prima $\alpha>0$ e $\beta=+\infty$. Definitivamente $a_n>\alpha/2$. Per ogni $M>0$, definitivamente $b_n>2M/\alpha>0$, quindi $a_nb_n>M$. Se $\alpha<0$ oppure $\beta=-\infty$, si applica lo stesso ragionamento alle successioni con i segni cambiati, ottenendo il segno previsto dalla tabella. Il cambio di segno nel limite segue direttamente da $|-a_n-(-\alpha)|=|a_n-\alpha|$ nei casi finiti e dalle definizioni con le code nei casi infiniti. Si scambiano i fattori quando necessario.

\textbf{Due limiti infiniti.} Se $a_n,b_n\to+\infty$, allora per ogni $M>0$ definitivamente entrambi i fattori sono maggiori di $\sqrt M$, dunque $a_nb_n>M$. Gli altri casi si riducono a questo cambiando i segni dei fattori; si ottiene $+\infty$ per segni uguali e $-\infty$ per segni opposti.
\end{proof}

\begin{osservazione}[Forme indeterminate]
Se le operazioni fra limiti non sono definite, i limiti dei singoli fattori o addendi non determinano il risultato. Per esempio, con $c\in\R$,
\[
n+(-n+c)=c,
\]
benché i due addendi tendano a $+\infty$ e $-\infty$. Per i prodotti con fattori tendenti a $0$ e a $+\infty$ si hanno, per $n\ge1$,
\[
\frac1{n^2}\,n=\frac1n\to0,\qquad
\frac1n\,n=1,\qquad
\frac1n\,n^2=n\to+\infty.
\]
\end{osservazione}

\begin{esercizio}[L'oscillazione di $\sin n$]
% Fonte: Lezione 11, p. 14: esempio senza dimostrazione.
% Svolgimento aggiunto; identità trigonometriche di addizione richiamate in nota.
Mostrare che la successione $a_n=\sin n$, con gli angoli misurati in radianti, non ha limite.
\end{esercizio}
\begin{svolgimento}
La successione è limitata fra $-1$ e $1$, quindi non può tendere a un infinito. Supponiamo per assurdo $\sin n\to L\in\R$. Le successioni traslate $\sin(n+1)$ e $\sin(n+2)$ sono sottosuccessioni e hanno lo stesso limite $L$. Dall'identità\footnote{Si usano qui le formule trigonometriche di addizione $\sin(u+v)=\sin u\cos v+\cos u\sin v$ e $\sin(u-v)=\sin u\cos v-\cos u\sin v$, con gli angoli in radianti.}
\[
\sin(n+2)+\sin n=2\cos1\,\sin(n+1)
\]
e dalle regole per somma e prodotto finiti segue $2L=2\cos1\,L$. Poiché $\cos1\ne1$, necessariamente $L=0$.

\Needspace{10\baselineskip}
L'altra identità
\[
\cos n=\frac{\sin(n+1)-\cos1\,\sin n}{\sin1}
\]
è lecita perché $\sin1\ne0$. Il suo membro destro è una combinazione lineare di successioni infinitesime, dunque $\cos n\to0$. Le regole del prodotto e della somma darebbero allora $\sin^2n+\cos^2n\to0$, in contraddizione con l'identità $\sin^2n+\cos^2n=1$ e con l'unicità del limite.
\end{svolgimento}

\Needspace{26\baselineskip}
\section{Avvicinamento da sopra e da sotto}
% Fonte: Lezione 12, p. 10. Le fonti terminano qui: nessun teorema del quoziente.

Per descrivere il segno della differenza $a_n-\ell$ occorre aggiungere un'informazione alla sola convergenza.

\begin{definizione}
Per $\ell\in\R$ si scrive $a_n\to\ell^+$ se $a_n\to\ell$ e $a_n>\ell$ definitivamente; si scrive $a_n\to\ell^-$ se $a_n\to\ell$ e $a_n<\ell$ definitivamente. Si dice che la successione tende a $\ell$ rispettivamente \emph{da sopra} o \emph{da sotto}.
\end{definizione}

\begin{esempio}
Per $n\ge1$, $\ell+1/n\to\ell^+$ e $\ell-1/n\to\ell^-$. La successione $\ell+(-1)^n/n$ converge a $\ell$, ma assume valori maggiori e minori di $\ell$ a indici arbitrariamente grandi: non tende a $\ell$ né da sopra né da sotto. La convergenza da un lato non richiede monotonia; richiede soltanto la posizione definitiva rispetto al limite.
\end{esempio}

\begin{schema}
\begin{tikzpicture}[x=.72cm,y=1cm]
  \draw[asse] (-.4,0) -- (10.7,0) node[right] {$n$};
  \draw[asse] (0,-.2) -- (0,3);
  \draw[black!60] (0,1.5) -- (10.3,1.5) node[right] {$\ell$};
  \foreach \n in {1,...,10} {
    \fill[coloreschema] (\n,{1.5+1.1/\n}) circle[radius=1.8pt];
    \draw[draw=black!60,fill=white] (\n,{1.5-1.1/\n}) circle[radius=1.8pt];
  }
\end{tikzpicture}
\par\smallskip I punti pieni tendono a $\ell$ da sopra, quelli vuoti da sotto. Una convergenza generica può invece attraversare ripetutamente il livello $\ell$.
\end{schema}


% !TeX root = ../main.tex
\chapter{Continuità}
\label{cap:continuita}
\fonti{Lezioni 13 e 14}
% Fonte: Lezione 13, pp. 13--17. Prova dell'algebra e chiarimenti su estremi/isolati aggiunti.

\section{Continuità e successioni}

Una funzione è continua in un punto quando avvicinare gli argomenti a quel punto avvicina anche i valori della funzione al valore assunto nel punto. Le successioni permettono di esprimere questa idea usando la nozione di limite già introdotta.

\begin{definizione}[Continuità in un punto e su un insieme]
\label{def:continuita-sequenziale}
Siano $f\colon A\subseteq\R\to\R$ e $x_0\in A$. Si dice che $f$ è \emph{continua in $x_0$} se, per ogni successione $(x_n)$ a valori in $A$,
\[
x_n\longrightarrow x_0
\quad\Longrightarrow\quad
f(x_n)\longrightarrow f(x_0).
\]
Se è continua in ogni punto di $A$, si dice \emph{continua su $A$} e si scrive $f\in C(A)$ oppure $f\in C^0(A)$.
\end{definizione}

La condizione riguarda tutte le successioni del dominio che tendono a $x_0$, comprese quelle che assumono il valore $x_0$. Il punto deve appartenere al dominio, perché compare $f(x_0)$; la notazione $C^0$ non richiede l'esistenza di derivate.

\begin{esempio}[Il valore assoluto]
Per ogni $x_0\in\R$, la proposizione~\ref{prop:limite-modulo} dà $x_n\to x_0\Rightarrow|x_n|\to|x_0|$. Pertanto $x\mapsto|x|$ è continua su tutta $\R$, compreso $0$.
\end{esempio}

\begin{schema}
\begin{tikzpicture}[x=1.6cm,y=1cm]
  \draw[asse] (-.3,0) -- (5,0) node[right] {$x$};
  \draw[asse] (0,-.2) -- (0,3.4) node[above] {$y$};
  \draw[intervallo,domain=.45:4.5,samples=80] plot (\x,{1+.35*\x+.06*(\x-2)^3});
  \draw[dashed,black!55] (3.5,0) -- (3.5,2.4275) -- (0,2.4275);
  \fill[coloreschema] (3.5,2.4275) circle[radius=2pt];
  \foreach \a in {2.8,3.1,3.3,3.4} {
    \pgfmathsetmacro{\b}{1+.35*\a+.06*(\a-2)^3}
    \draw[black!40] (\a,0) -- (\a,\b) -- (0,\b);
    \fill[coloreschema] (\a,\b) circle[radius=1.6pt];
  }
  \node[below] at (3.5,0) {$x_0$};
  \node[below] at (2.8,-.15) {$x_n$};
  \draw[asse,coloreschema] (2.85,-.55) -- (3.4,-.55);
  \node[left] at (0,2.4275) {$f(x_0)$};
  \node[left] at (0,1.7) {$f(x_n)$};
\end{tikzpicture}
\par\smallskip Schema di continuità: ai punti $x_n$ che si avvicinano a $x_0$ corrispondono ordinate che si avvicinano a $f(x_0)$. Il disegno mostra una sola scelta di successione; la definizione le richiede tutte.
\end{schema}

\section{Il ruolo del dominio}
% Fonte: Lezione 13, p. 16; preservati dominio [0,3] unito a {4}, valore 5 e x_n=3-1/n.

La continuità si valuta usando soltanto punti del dominio. A un estremo di un intervallo le successioni possono avvicinarsi soltanto dal lato disponibile; un punto isolato presenta una situazione ancora diversa.

\begin{esempio}[Un estremo e un punto isolato]
Consideriamo una funzione $f\colon[0,3]\cup\{4\}\to\R$, continua sul tratto $[0,3]$, e poniamo $f(4)=5$. La successione $x_n=3-1/n$, per $n\ge1$, appartiene al dominio e tende a $3$. La continuità in $3$ richiede $f(3-1/n)\to f(3)$; non impone che questi valori tendano a $f(4)=5$, perché $x_n$ non tende a $4$.

Il punto $4$ è isolato nel dominio. Se $x_n\in[0,3]\cup\{4\}$ e $x_n\to4$, usando il raggio $1/2$ si ha definitivamente $x_n\in(7/2,9/2)$. L'unico punto del dominio in tale intervallo è $4$, quindi $x_n=4$ e $f(x_n)=5$ definitivamente. La funzione è continua anche in $4$, qualunque sia il valore che le si assegna in quel punto.
\end{esempio}

\begin{schema}
\begin{tikzpicture}[x=1.65cm,y=.7cm]
  \draw[asse] (-.3,0) -- (5,0) node[right] {$x$};
  \draw[asse] (0,-.2) -- (0,5.8) node[above] {$y$};
  \draw[intervallo,domain=0:3,samples=70] plot (\x,{1+\x/3+\x*(\x-1.5)*(\x-3)/5});
  \fill[coloreschema] (0,1) circle[radius=2pt];
  \fill[coloreschema] (3,2) circle[radius=2pt];
  \fill[coloreschema] (4,5) circle[radius=2pt];
  \draw[dashed,black!50] (4,0) -- (4,5) -- (0,5);
  \draw[dotted,black!50] (3,0) -- (3,2);
  \foreach \a in {0,3,4} {\node[below] at (\a,0) {$\a$};}
  \node[left] at (0,5) {$5$};
  \node[right] at (3,2) {$f(3)$};
  \draw[asse,coloreschema] (2.05,.12) -- (2.95,.12);
  \node[below] at (2.3,-.25) {$3-1/n$};
\end{tikzpicture}
\par\smallskip Fra $3$ e $4$ il dominio non contiene punti: il grafico non deve essere congiunto attraverso questo intervallo. La continuità nel punto isolato $4$ non dipende dal valore all'estremo $3$.
\end{schema}

Il ragionamento sul punto $4$ vale in generale: se $x_0$ è isolato in $A$, ogni successione di $A$ che tende a $x_0$ è definitivamente uguale a $x_0$. Ogni funzione $A\to\R$ è dunque continua nei punti isolati del proprio dominio.

\section{Funzioni elementari e operazioni}
% Fonte: Lezione 13, p. 17. Primo teorema enunciato senza prova nella fonte; prova del secondo aggiunta.

\begin{teorema}[Continuità delle funzioni elementari]
\label{teo:continuita-elementari}
Le funzioni elementari già introdotte, potenze, funzioni trigonometriche, esponenziali, logaritmi e funzioni iperboliche, e le loro inverse sui rami in cui sono definite, sono continue nei rispettivi domini.\footnote{Questo risultato è assunto qui come proprietà delle funzioni elementari. Non è una conseguenza della sola definizione: le dimostrazioni generali, in particolare per trigonometria ed esponenziali, non sono sviluppate in questa trattazione. Restano le ipotesi dichiarate nella costruzione dell'esponenziale reale.}
\end{teorema}

L'affermazione riguarda soltanto punti del dominio: per esempio, il logaritmo è continuo per argomenti positivi; la tangente nei punti in cui il coseno non si annulla. Per le radici di indice pari la continuità in $0$ si valuta nel dominio $[0,+\infty)$.

\begin{teorema}[Algebra delle funzioni continue]
Siano $f,g\colon A\subseteq\R\to\R$ continue in $x_0\in A$. \begin{enumerate}[label=\arabic*.]
\item Allora $f+g$, $f-g$ e $fg$ sono continue in $x_0$.
\item Se inoltre $g(x_0)\ne0$, anche $f/g$, definita su $B=\{x\in A:g(x)\ne0\}$, è continua in $x_0$.
\end{enumerate}
\end{teorema}
\begin{proof}
\begin{enumerate}[label=\arabic*.]
\item \textbf{Somma, differenza e prodotto.} Per ogni successione $x_n\in A$ tendente a $x_0$, la continuità dà $f(x_n)\to f(x_0)$ e $g(x_n)\to g(x_0)$. L'algebra dei limiti finiti implica
\[
\begin{aligned}
f(x_n)\pm g(x_n)&\longrightarrow f(x_0)\pm g(x_0),\\
f(x_n)g(x_n)&\longrightarrow f(x_0)g(x_0).
\end{aligned}
\]
Sono esattamente le condizioni di continuità delle tre funzioni in $x_0$.
\item \textbf{Quoziente.} Poiché $g(x_0)\ne0$, si ha $x_0\in B$. Se $x_n\in B$ e $x_n\to x_0$, valgono gli stessi limiti dei valori di $f$ e $g$; il teorema~\ref{teo:limite-quoziente} dà
\[
\frac{f(x_n)}{g(x_n)}\longrightarrow\frac{f(x_0)}{g(x_0)}.
\]
Questa è la continuità di $f/g$ in $x_0$ sul suo dominio $B$.
\end{enumerate}
\end{proof}

\begin{esempio}
Ogni polinomio è continuo su $\R$, perché è una somma finita di costanti e potenze. Ogni funzione razionale $P/Q$ è continua dove $Q$ non si annulla. Per esempio $1/(x^2-4)$ è continua su $\R\setminus\{-2,2\}$: il risultato non assegna valori nei due punti esclusi.
\end{esempio}

% !TeX root = ../main.tex
% Fonte: Lezione 14, pp. 2--6. Accorpati i richiami alla definizione; prova della composizione aggiunta.
\section{Composizione e calcolo dei limiti}

\begin{proposizione}[Continuità della composta]
Siano $f\colon A\to B$ e $g\colon B\to\R$, con $A,B\subseteq\R$. Se $f$ è continua in $x_0\in A$ e $g$ è continua in $f(x_0)$, allora $g\circ f$ è continua in $x_0$.
\end{proposizione}
\begin{proof}
Per ogni successione $x_n\in A$ con $x_n\to x_0$, la continuità di $f$ dà $f(x_n)\to f(x_0)$. I valori $f(x_n)$ appartengono a $B$; la continuità di $g$ nel punto $f(x_0)$ implica quindi $g(f(x_n))\to g(f(x_0))$, come richiesto.
\end{proof}

La stessa definizione si applica direttamente al calcolo di un limite: se $x_n\to x_0$, i termini appartengono al dominio di $f$ e $f$ è continua in $x_0$, allora $f(x_n)\to f(x_0)$. Si può dunque calcolare prima il limite dell'argomento, controllare dominio e continuità, e poi applicare la funzione al limite trovato. Per una successione definita solo definitivamente si lavora sulla sua coda.

\begin{esercizio}[Continuità e fattore oscillante]
Calcolare, per $n\to+\infty$,
\[
\sin\!\left(\left|2\left[\cos\!\left(\frac1{n^2}\right)-1\right]\sin n-\pi\right|\right).
\]
\end{esercizio}
\begin{svolgimento}
Poiché $1/n^2\to0$ e il coseno è continuo, $\cos(1/n^2)\to1$. Quindi $2[\cos(1/n^2)-1]\to0$. La successione $\sin n$ è limitata, anche se non ha limite, e il prodotto di una limitata per un'infinitesima tende a zero:
\[
2[\cos(1/n^2)-1]\sin n-\pi\longrightarrow-\pi.
\]
Per la continuità del modulo, i valori assoluti tendono a $\pi$. Infine, per la continuità del seno, la successione richiesta tende a $\sin\pi=0$.
\end{svolgimento}

\begin{esercizio}[Una composizione di funzioni elementari]
Calcolare
\[
\lim_{n\to+\infty}\sqrt[3]{\cos(1/n+\sin(1/n^3))}.
\]
\end{esercizio}
\begin{svolgimento}
Si procede dall'argomento più interno verso l'esterno. Poiché $1/n^3\to0$ e il seno è continuo, $\sin(1/n^3)\to0$. Sommando $1/n\to0$ si ottiene $1/n+\sin(1/n^3)\to0$. La continuità del coseno dà
\[
\cos(1/n+\sin(1/n^3))\longrightarrow\cos0=1.
\]
La radice cubica è definita e continua su tutta $\R$, perciò il limite cercato è $\sqrt[3]1=1$.
\end{svolgimento}

\section{Quando la sostituzione non è lecita}

\begin{esempio}[Una funzione discontinua]
Consideriamo
\[
f(x)=\begin{cases}1,&x=0,\\0,&x\ne0.\end{cases}
\]
Per $n\ge1$, $f(1/n)=0$, quindi $f(1/n)\to0$. Sebbene $1/n\to0$, il limite non è $f(0)=1$: manca la continuità in $0$. La successione costante $x_n=0$ dà invece $f(x_n)=1$, confermando che le immagini di successioni tendenti allo stesso punto possono avere limiti diversi.
\end{esempio}

\begin{schema}
\begin{tikzpicture}[x=1.6cm,y=1.5cm]
  \draw[asse] (-2.2,0) -- (2.4,0) node[right] {$x$};
  \draw[asse] (0,-.3) -- (0,1.55) node[above] {$y$};
  \draw[intervallo] (-2,0) -- (2.1,0);
  \draw[estremo] (0,0) circle[radius=2pt];
  \fill[coloreschema] (0,1) circle[radius=2pt] node[left] {$1$};
  \node[below left] at (0,0) {$0$};
\end{tikzpicture}
\par\smallskip Il punto pieno indica $f(0)=1$; il punto vuoto esclude $(0,0)$ dal grafico. Avvicinarsi a zero attraverso punti non nulli dà sempre valore $0$.
\end{schema}

\begin{esempio}[Controllare prima il dominio]
La formula $\sqrt{\cos(1/n^2)-1}$, per $n\ge1$, non definisce una successione reale: infatti $0<1/n^2\le1<2\pi$, dunque $\cos(1/n^2)<1$ e il radicando è negativo. Il fatto che il radicando tenda a zero non rende definita la radice quadrata dei suoi termini. Non è quindi lecito concludere che la successione abbia limite $\sqrt0$ in $\R$.
\end{esempio}


\end{document}
